% !TEX root = print.tex

\part{Triebel--Lizorkin spaces}
\label{part:triebel-lizorkin}

This part treats the Triebel--Lizorkin spaces $\TL{s}{p}{q}$ on $\R^d$ and on the torus: the
Besov sandwich from the comparison of mixed norms, the spaces of tempered distributions for
every real $s$, and, with the Fefferman--Stein inequality, the characterisation by local means
on both, the independence of the Littlewood--Paley family, the characterisations by ball means
of differences and by atoms on $\R^d$, and the real interpolation spaces
(\Cref{ch:triebel-lizorkin}). \Cref{ch:singular-integrals} proves the
strong type of vector-valued singular integrals with Calder\'on--Zygmund kernel families and,
from it, the square function theorem $\TL{0}{p}{2} = L_p$ for $n \geq 1$ and $1 < p < \infty$
(\Cref{cor:square-function-theorem}), on $\R^d$ and on the torus
(\Cref{cor:torus-square-function-theorem}).

\chapter{Triebel--Lizorkin spaces}
\label{ch:triebel-lizorkin}

Besov quasi-norms take the $L_p$ norm of every Littlewood--Paley block and then the $\ell_q$ norm
over the levels. Triebel--Lizorkin quasi-norms take the two norms in the other order: the $\ell_q$
norm over the levels at every point, and then the $L_p$ norm. This chapter works on a
finite-dimensional real inner product space $V$ of dimension $n$ with Lebesgue measure, as in
\Cref{ch:rd}, and on the torus; the two mixed norms are compared in \Cref{thm:mixed-norms}, and
the maximal function and the Fefferman--Stein inequality are those of
\Cref{ch:maximal-functions}. For $s > 0$ the spaces
$\TL{s}{p}{q}$ and their descriptions by local means, ball means of differences and atoms are
those of Triebel~\cite{TriebelII}; the descriptions are proved here for $1 < p < \infty$ and
$1 < q \leq \infty$. The atoms here are normalised in $L_\infty$ and the sequence space
$f^s_{p,q}$ carries the weights $2^{js}$; Triebel states the atomic characterisation with
$(s,p)$-atoms normalised in $L_p$, and the two forms correspond under
$\lambda_{j,m} \mapsto 2^{j(s - n/p)} \lambda_{j,m}$. \Cref{sec:tl-distributions} defines
$\TL{s}{p}{q}$ for tempered distributions and every real $s$, and
\Cref{sec:lp-family-independence} shows that every admissible Littlewood--Paley family gives
the same spaces. The last two sections identify
the real interpolation spaces of Triebel--Lizorkin and of Bessel potential spaces with Besov spaces
(\Cref{sec:tl-interpolation}), and those of Triebel--Lizorkin spaces on the torus
(\Cref{sec:torus-tl-interpolation}).

\section{The spaces \texorpdfstring{$\TL{s}{p}{q}$}{F\^{}s\_\{p,q\}} and the Besov sandwich}
\label{sec:tl-besov-sandwich}

As in \Cref{ch:rd}, $V$ is a finite-dimensional real inner product space with Lebesgue measure
and $L_p = L_p(V)$ for $1 \leq p < \infty$. The Littlewood--Paley kernels $\psi_j$ are those of
\Cref{def:lp-kernels}, and the Besov space $\Besov{s}{p}{q}$ is the approximation space $A^s_q$
of the band-limited scheme in $L_p$ (\Cref{thm:lp-band}).

\begin{lemma}[Pointwise bound of a convolution]
\label{lem:conv-pointwise-bound}
\lean{Rd.iSup_enorm_lt_top_of_bddAbove, Rd.enorm_conv_le_of_enorm_le}
\leanok
\uses{def:conv, lem:weighted-jensen}
For every $1 \leq p < \infty$, every measurable $K : V \to \C$ and every $B \in [0, \infty]$ with
$|K(y)| \leq B$ for all $y \in V$, every measurable $f : V \to \C$ and every $x \in V$,
\[
  |(K * f)(x)| \leq \|K\|_1^{1 - 1/p} \, B^{1/p} \, \|f\|_p .
\]
For bounded $K$, the supremum $B = \sup_{y \in V} |K(y)|$ is finite.
\end{lemma}

\begin{proof}
\leanok
By the weighted Jensen inequality (\Cref{lem:weighted-jensen}) with the weight $|K|$,
\[
  \Bigl( \int |K(y)| \, |f(x - y)| \, dy \Bigr)^p
  \leq \|K\|_1^{p-1} \int |K(y)| \, |f(x - y)|^p \, dy ,
\]
and the last integral is at most $B \|f\|_p^p$ by the translation invariance of Lebesgue
measure. Take the $p$-th root.
\end{proof}

The bound covers $p = 1$, where it reads $|(K * f)(x)| \leq B \|f\|_1$.

\begin{definition}[Triebel--Lizorkin spaces]
\label{def:triebel-lizorkin}
\lean{Rd.blockFun, Rd.lpBlock_ae_eq_blockFun, Rd.aestronglyMeasurable_blockFun,
  Rd.enorm_lpBlock_eq_eLpNorm_blockFun, Rd.triebelLizorkinNorm, Rd.MemTriebelLizorkin,
  Rd.aemeasurable_lqNorm_blockFun, Rd.littlewoodPaleyNorm_eq_lqNorm_eLpNorm}
\leanok
\uses{def:lp-blocks, def:conv, def:lq-norm}
Let $1 \leq p < \infty$, $s \in \R$ and $1 \leq q \leq \infty$. For $f \in L_p$ and $j \in \N$ the
pointwise block is $\psi_j * f$, the convolution integral at every point of $V$; it represents
the Littlewood--Paley block $\Delta_j f \in L_p$. The Triebel--Lizorkin quasi-norm of $f$ is
\[
  \|f\|_{F^s_{p,q}} = \Bigl\| x \mapsto \bigl\| (2^{js} |(\psi_j * f)(x)|)_{j \in \N}
  \bigr\|_{\ell_q} \Bigr\|_{L_p},
\]
and $\TL{s}{p}{q}$ consists of the $f \in L_p$ with $\|f\|_{F^s_{p,q}} < \infty$.
\end{definition}

For $s < 0$ this is the part $\TL{s}{p}{q} \cap L_p$ of the classical space of tempered
distributions; \Cref{sec:tl-distributions} defines $\TL{s}{p}{q}$ for tempered
distributions.

The Littlewood--Paley quasi-norm of the Besov spaces (\Cref{def:lp-blocks}) takes the two norms
in the other order, $\|f\|^{\mathrm{LP}}_{s,q} = \|(2^{js} \|\psi_j * f\|_p)_j\|_{\ell_q}$; in the
notation of the mixed norms, $\|f\|_{F^s_{p,q}}$ is an $L_p(\ell_q)$ norm and
$\|f\|^{\mathrm{LP}}_{s,q}$ an $\ell_q(L_p)$ norm of the same family.

\begin{lemma}[The pointwise blocks]
\label{lem:tl-pointwise-blocks}
\lean{Rd.iSup_enorm_lpBlockKernel_lt_top, Rd.eLpNorm_lpBlockKernel_mul_iSup_lpBlockKernel_ne_top,
  Rd.enorm_blockFun_le, Rd.blockFun_add_ae, Rd.blockFun_smul, Rd.blockFun_eq_of_ae_eq,
  Rd.memLp_blockFun, Rd.eLpNorm_blockFun_zero_le}
\leanok
\uses{def:triebel-lizorkin, lem:conv-pointwise-bound, lem:lp-kernels-basic, lem:conv-assoc,
  thm:young, def:conv}
For every $1 \leq p < \infty$, all $f, g \in L_p$, every $c \in \C$, every $j \in \N$ and every
$x \in V$, $\psi_j * f \in L_p$, $\|\psi_0 * f\|_p \leq \|\Phi\|_1 \, \|f\|_p$,
\[
  |(\psi_j * f)(x)| \leq \|\psi_j\|_1^{1 - 1/p} \bigl( \sup |\psi_j| \bigr)^{1/p} \|f\|_p
\]
with $\|\psi_j\|_1^{1 - 1/p} ( \sup |\psi_j| )^{1/p} < \infty$, and
$(\psi_j * (c f))(x) = c \, (\psi_j * f)(x)$. For almost every $x \in V$,
$(\psi_j * (f + g))(x) = (\psi_j * f)(x) + (\psi_j * g)(x)$ for all $j \in \N$. For every
$1 \leq r \leq \infty$ and every $h \in L_r$ with $h = f$ almost everywhere,
$\psi_j * h = \psi_j * f$.
\end{lemma}

\begin{proof}
\leanok
Young's inequality (\Cref{thm:young}) gives $\psi_j * f \in L_p$, and with
$\|\psi_0\|_1 = \|\Phi_0\|_1 = \|\Phi\|_1$ (\Cref{lem:lp-kernels-basic}) the bound of
$\psi_0 * f$. The pointwise bound is \Cref{lem:conv-pointwise-bound}; its constant is finite
since $\psi_j$ is integrable and bounded (\Cref{lem:lp-kernels-basic}). Homogeneity holds for
the integrals, and additivity holds almost everywhere for each $j$ by \Cref{lem:conv-assoc},
hence for all $j$ at once outside a countable union of null sets. Convolutions depend only on
the class of a function (\Cref{def:conv}).
\end{proof}

\begin{proposition}[Elementary properties]
\label{prop:tl-elementary}
\lean{Rd.triebelLizorkinNorm_add_le, Rd.triebelLizorkinNorm_smul, Rd.triebelLizorkinNorm_zero,
  Rd.triebelLizorkinNorm_zero_left, Rd.triebelLizorkinNorm_mono_exponent,
  Rd.triebelLizorkinNorm_mono_smoothness, Rd.MemTriebelLizorkin.mono_exponent,
  Rd.MemTriebelLizorkin.mono_smoothness, Rd.MemTriebelLizorkin.add, Rd.MemTriebelLizorkin.smul,
  Rd.triebelLizorkinSpace, Rd.mem_triebelLizorkinSpace}
\leanok
\uses{def:triebel-lizorkin, lem:tl-pointwise-blocks, lem:lq-basic, lem:lp-norm-homogeneity,
  lem:mixed-norm-dyadic-family}
For every $1 \leq p < \infty$, all $s \leq t$ in $\R$, all $1 \leq q \leq r \leq \infty$, all
$f, g \in L_p$ and every $c \in \C$,
\[
  \|f + g\|_{F^s_{p,q}} \leq \|f\|_{F^s_{p,q}} + \|g\|_{F^s_{p,q}}, \qquad
  \|c f\|_{F^s_{p,q}} = |c| \, \|f\|_{F^s_{p,q}},
\]
\[
  \|f\|_{F^s_{p,r}} \leq \|f\|_{F^s_{p,q}}, \qquad \|f\|_{F^s_{p,q}} \leq \|f\|_{F^t_{p,q}} ,
\]
and $\|f\|_{F^0_{p,q}} = \bigl\| x \mapsto \|((\psi_j * f)(x))_j\|_{\ell_q} \bigr\|_{L_p}$.
In particular $\TL{s}{p}{q}$ is a linear subspace of $L_p$,
$\TL{s}{p}{q} \subseteq \TL{s}{p}{r}$ and $\TL{t}{p}{q} \subseteq \TL{s}{p}{q}$.
\end{proposition}

\begin{proof}
\leanok
At almost every point the blocks of $f + g$ are the sums of the blocks of $f$ and $g$
(\Cref{lem:tl-pointwise-blocks}), so the triangle inequality in $\ell_q$
(\Cref{lem:lq-basic}) and then in $L_p$ gives the first claim. The blocks of $c f$ are $c$ times
those of $f$, and $|c|$ leaves the $\ell_q$ and the $L_p$ norm
(\Cref{lem:lp-norm-homogeneity}). The monotonicity follows from
$\|\cdot\|_{\ell_r} \leq \|\cdot\|_{\ell_q}$ and $2^{js} \leq 2^{jt}$ at every point. For
$s = 0$ the weights $2^{j \cdot 0}$ are $1$, for every $q$ and in particular for $q = 2$. These
are the properties of the mixed norms of \Cref{lem:mixed-norm-dyadic-family} for the families
$(|\psi_j * f|)_j$.
\end{proof}

\begin{lemma}[The weighted blocks and their tails]
\label{lem:weighted-block-tails}
\lean{Rd.eLpNorm_lqNorm_ite_blockFun_le, Rd.tendsto_eLpNorm_lqNorm_blockFun_tail}
\leanok
\uses{def:triebel-lizorkin, lem:mixed-norm-tails}
Let $1 \leq p < \infty$, $s \in \R$, $1 \leq q \leq \infty$ and $f \in L_p(V)$. For every set
$P \subseteq \N$,
\[
  \bigl\| \|(\mathbf 1_P(j) \, 2^{js} |(\psi_j * f)(x)|)_j\|_{\ell_q} \bigr\|_p
  \leq \|f\|_{\TL{s}{p}{q}} .
\]
If $q < \infty$ and $f \in \TL{s}{p}{q}$, then
\[
  \bigl\| \|(2^{js} |(\psi_j * f)(x)|)_{j \geq N}\|_{\ell_q} \bigr\|_p \to 0 \qquad
  (N \to \infty) ,
\]
where the $\ell_q$ norm is taken of the sequence with the terms $j < N$ replaced by $0$.
\end{lemma}

\begin{proof}
\leanok
Both claims are \Cref{lem:mixed-norm-tails} for the functions
$g_j = 2^{js} |\psi_j * f|$, whose mixed norm is $\|f\|_{\TL{s}{p}{q}}$
(\Cref{def:triebel-lizorkin}).
\end{proof}

\begin{theorem}[$\TL{s}{p}{p} = \Besov{s}{p}{p}$]
\label{thm:tl-diagonal}
\lean{Rd.triebelLizorkinNorm_eq_littlewoodPaleyNorm, Rd.memTriebelLizorkin_iff_memApprox,
  Rd.memTriebelLizorkin_iff_modulus}
\leanok
\uses{def:triebel-lizorkin, cor:mixed-norm-sandwich, thm:lp-band, thm:rd-besov}
For every $1 \leq p < \infty$, every $s \in \R$ and every $f \in L_p$,
$\|f\|_{F^s_{p,p}} = \|f\|^{\mathrm{LP}}_{s,p}$. If moreover $s > 0$, then $f \in \TL{s}{p}{p}$ if
and only if $f \in \Besov{s}{p}{p}$, and for every $r \in \N$ with $s < r$, $f \in \TL{s}{p}{p}$ if
and only if $\|f\|_p + \|(2^{ks} \omega_r(f, 2^{-k})_p)_k\|_{\ell_p} < \infty$.
\end{theorem}

\begin{proof}
\leanok
The $L_p$ norm of the pointwise block $\psi_j * f$ is $\|\Delta_j f\|_p$, so the identity is
$\|g\|_{L_p(\ell_p)} = \|g\|_{\ell_p(L_p)}$ for $g_j = 2^{js} |\psi_j * f|$
(\Cref{cor:mixed-norm-sandwich}). The two characterisations are \Cref{thm:lp-band} and
\Cref{thm:rd-besov} with $q = p$.
\end{proof}

\begin{theorem}[The Besov sandwich]
\label{thm:tl-besov-sandwich}
\lean{Rd.littlewoodPaleyNorm_max_le_triebelLizorkinNorm,
  Rd.triebelLizorkinNorm_le_littlewoodPaleyNorm_min, Rd.approxNormDyadic_le_triebelLizorkinNorm,
  Rd.triebelLizorkinNorm_le_approxNormDyadic, Rd.memApprox_of_memTriebelLizorkin,
  Rd.memTriebelLizorkin_of_memApprox}
\leanok
\uses{def:triebel-lizorkin, cor:mixed-norm-sandwich, thm:lp-band, def:dyadic-norm}
For every $1 \leq p < \infty$, every $s \in \R$, every $1 \leq q \leq \infty$ and every
$f \in L_p$,
\[
  \|f\|^{\mathrm{LP}}_{s,\max(p,q)} \leq \|f\|_{F^s_{p,q}} \leq \|f\|^{\mathrm{LP}}_{s,\min(p,q)} .
\]
If $s \geq 0$, then
\[
  \|f\|_{F^s_{p,q}} \leq 8 \cdot 8^s \, \|\Phi\|_1 \, \|f\|^{\mathrm{dy}}_{A^s_{\min(p,q)}}
\]
with the band-limited scheme. If $s > 0$, then
\[
  \|f\|^{\mathrm{dy}}_{A^s_{\max(p,q)}} \leq 2 (1 - 2^{-s})^{-1} \, \|f\|_{F^s_{p,q}}
\]
and
\[
  \Besov{s}{p}{\min(p,q)} \subseteq \TL{s}{p}{q} \subseteq \Besov{s}{p}{\max(p,q)} .
\]
\end{theorem}

\begin{proof}
\leanok
The first two inequalities are \Cref{cor:mixed-norm-sandwich} for the family
$g_j = 2^{js} |\psi_j * f|$, whose $L_p$ norms are $2^{js} \|\Delta_j f\|_p$. The bounds by the
approximation norms follow with \Cref{thm:lp-band} at the exponents $\max(p,q)$ and
$\min(p,q)$, and so do the inclusions.
\end{proof}

For $q = 2$ the theorem places $\TL{s}{p}{2}$ between $\Besov{s}{p}{\min(p,2)}$ and
$\Besov{s}{p}{\max(p,2)}$: for $p \geq 2$ between $\Besov{s}{p}{2}$ and $\Besov{s}{p}{p}$, for
$p \leq 2$ between $\Besov{s}{p}{p}$ and $\Besov{s}{p}{2}$.

\section{Triebel--Lizorkin spaces of tempered distributions}
\label{sec:tl-distributions}

The quasi-norm of \Cref{def:triebel-lizorkin} is defined on $L_p$. Here the Littlewood--Paley
blocks act on tempered distributions as Fourier multipliers, $\Delta_j u = (\widehat{\psi_j})(D) u$,
and a tempered distribution lies in $\TL{s}{p}{q}$ when every block is an $L_p$ function and the
mixed norm of the weighted blocks is finite; for $p < \infty$ this describes $\TL{s}{p}{q}$ for
every real $s$ (Triebel defines $F^s_{p,q}$ for $p < \infty$; the standard space
$F^s_{\infty,q}$, $q < \infty$, of Frazier and Jawerth has a different quasi-norm). The
blocks need the kernels $\psi_j = (G^\psi_j)_{2^j}$ as Schwartz functions (\Cref{lem:psi-base},
\Cref{def:schwartz-dilate}). The Fourier multiplier $m(D)$ of a tempered distribution is defined
in the paragraph before \Cref{lem:fourier-multiplier-fourier}; Fourier multipliers on these spaces
are the subject of \Cref{sec:tl-multipliers}.

\paragraph{Convolutions as Fourier multipliers.}

\begin{lemma}[Convolutions as Fourier multipliers]
\label{lem:conv-fourier-multiplier}
\lean{Rd.fourierMultiplierCLM_fourier_eq_convolution, Rd.fourierMultiplierCLM_fourier_eq_conv,
  Rd.fourierMultiplierCLM_fourier_toLp_eq_conv, Rd.coeFn_sum_convLp_ae}
\leanok
\uses{def:conv, lem:conv-op, lem:ae-equal-classes-distribution, lem:lp-kernels-basic}
For all Schwartz functions $\eta, \varphi \colon V \to \C$,
$(\widehat\eta)(D) \varphi = \eta * \varphi$, Mathlib's convolution of Schwartz functions. For
every $1 \leq r < \infty$, every Schwartz function $\eta$ and every $G \in L_r(V)$,
$(\widehat\eta)(D) G = \eta * G$ as tempered distributions. For every $1 \leq r < \infty$, all
Schwartz functions $K_j$, all $g_j \in L_r(V)$ and every finite set $t \subseteq \N$, the sum
$\sum_{j \in t} K_j * g_j$ of the classes in $L_r(V)$ is almost everywhere the pointwise sum of
the convolutions.
\end{lemma}

\begin{proof}
\leanok
Both sides of the first identity are Schwartz functions with the Fourier transform $\widehat\eta \,
\widehat\varphi$ (Mathlib's \texttt{SchwartzMap.\allowbreak fourier\_\allowbreak convolution}), and
the Fourier transform is injective on Schwartz functions; Mathlib's convolution of Schwartz
functions is the convolution integral (\texttt{SchwartzMap.\allowbreak convolution\_\allowbreak
apply}). For the second, both sides are continuous in $G \in L_r(V)$, as maps into the tempered
distributions: $(\widehat\eta)(D)$ is continuous on $\mathcal S'(V)$, $L_r(V)$ embeds continuously,
and $G \mapsto \eta * G$ is bounded on $L_r(V)$ (\Cref{lem:conv-op}). They agree on Schwartz
functions by the first identity, and Schwartz functions are dense in $L_r(V)$ for $r < \infty$
(Mathlib's \texttt{SchwartzMap.\allowbreak denseRange\_\allowbreak toLpCLM}); two continuous maps
into a Hausdorff space that agree on a dense set are equal.
\end{proof}

\paragraph{The blocks of tempered distributions.}

\begin{definition}[Blocks of tempered distributions]
\label{def:distrib-block}
\lean{Rd.fourier_psiS_apply, Rd.distribBlock}
\leanok
\uses{def:psi-base, lem:psi-base, def:schwartz-dilate, def:lp-kernels}
For every $j \in \N$ and every tempered distribution $u$ on $V$ let
$\Delta_j u = (\widehat{\psi_j})(D) u$, with the Schwartz function $\psi_j = (G^\psi_j)_{2^j}$
of \Cref{lem:psi-base}.
\end{definition}

\begin{lemma}[Blocks of $L_r$ functions]
\label{lem:distrib-block-lp}
\lean{Rd.distribBlock_toTemperedDistribution}
\leanok
\uses{def:distrib-block, lem:conv-fourier-multiplier, def:lp-blocks, def:triebel-lizorkin}
For every $1 \leq r < \infty$, every $f \in L_r(V)$ and every $j \in \N$, the block $\Delta_j f$
of the tempered distribution of $f$ is the tempered distribution of the block
$\psi_j * f \in L_r(V)$ of \Cref{def:lp-blocks}.
\end{lemma}

\begin{proof}
\leanok
\Cref{lem:conv-fourier-multiplier} with $\eta = \psi_j$.
\end{proof}

\begin{theorem}[Littlewood--Paley decomposition of tempered distributions]
\label{thm:distrib-block-decomposition}
\lean{Rd.sum_fourier_psiS_apply, Rd.sum_distribBlock, Rd.smulLeftCLM_sum_fourier_psiS,
  Rd.tendsto_sum_distribBlock}
\leanok
\uses{def:distrib-block, lem:fourier-psi-window, lem:lp-kernels-frequency, def:schwartz-cutoff,
  lem:schwartz-cutoff-tendsto}
For every $J \in \N$ and every $\xi \in V$, $\sum_{j \leq J} \widehat{\psi_j}(\xi) =
\chi(2^{-J} \xi)$ with the cut-off $\chi = \chi_{1,2}$ of \Cref{def:lp-kernels}, so
$\sum_{j \leq J} \Delta_j u$ is the Fourier multiplier with $\chi(2^{-J} \cdot)$ for every
tempered distribution $u$. For every tempered distribution $u$ on $V$,
\[
  \sum_{j < J} \Delta_j u \to u \qquad (J \to \infty)
\]
in $\mathcal S'(V)$, that is, $\sum_{j < J} \Delta_j u(\varphi) \to u(\varphi)$ for every
Schwartz function $\varphi$.
\end{theorem}

\begin{proof}
\leanok
The Fourier transforms of the blocks telescope to $\widehat{\Phi_J}(\xi) = \chi(2^{-J} \xi)$
(\Cref{lem:fourier-psi-window}, \Cref{lem:lp-kernels-frequency}), and Fourier multipliers add
(Mathlib's \texttt{TemperedDistribution.\allowbreak fourierMultiplierCLM\_\allowbreak sum}). For a
Schwartz function $\varphi$, with $w = \mathcal F^{-1} \varphi$,
$\sum_{j \leq J} \Delta_j u(\varphi) = u\bigl(\mathcal F(\chi(2^{-J} \cdot) \, w)\bigr)$, and
$\chi(2^{-J} \cdot) \, w$ is the cut-off of $w$ at radius $2^J$ (\Cref{def:schwartz-cutoff}). The
cut-offs converge to $w$ in the Schwartz space (\Cref{lem:schwartz-cutoff-tendsto}), and
$\mathcal F$ and $u$ are continuous, so the values converge to $u(\mathcal F w) = u(\varphi)$.
\end{proof}

\begin{lemma}[Limits of $L_p$ blocks]
\label{lem:distrib-block-limit}
\lean{Rd.eq_of_tendsto_sum_distribBlock}
\leanok
\uses{thm:distrib-block-decomposition}
Let $1 \leq p \leq \infty$, $w$ a tempered distribution on $V$, $G_j \in L_p(V)$ with
$\Delta_j w = G_j$ for every $j \in \N$, and $h \in L_p(V)$. If $\sum_{j < J} G_j \to h$ in
$L_p(V)$, then $w = h$ as tempered distributions.
\end{lemma}

\begin{proof}
\leanok
The partial sums tend to $w$ in $\mathcal S'(V)$ (\Cref{thm:distrib-block-decomposition}) and to
$h$, since $L_p(V)$ embeds continuously into $\mathcal S'(V)$; limits in $\mathcal S'(V)$ are
unique.
\end{proof}

\paragraph{Admissible kernels.}
The quasi-norm below is defined once, for the blocks of an admissible kernel; the kernel $\Phi$
of \Cref{def:lp-kernels} is one of them and gives the blocks $\Delta_j$. Every admissible kernel
gives the same spaces for $1 < p < \infty$ and $1 < q \leq \infty$
(\Cref{sec:lp-family-independence}).

\begin{definition}[Admissible kernels and their blocks]
\label{def:lp-kernel-admissible}
\lean{Rd.LPKernel, Rd.LPKernel.base, Rd.LPKernel.psiS, Rd.LPKernel.coe_psiS,
  Rd.LPKernel.distribBlock, Rd.LPKernel.standard}
\leanok
\uses{def:schwartz-dilate, def:lp-kernels, lem:lp-kernels-basic, def:psi-base, def:distrib-block}
An \emph{admissible kernel} $K$ is a Schwartz function $\Phi' \colon V \to \C$ with reals
$0 < a < b$ such that $\widehat{\Phi'}(\xi) = 1$ for $\|\xi\| \leq a$ and
$\widehat{\Phi'}(\xi) = 0$ for $\|\xi\| \geq b$. Its bases are $G'_0 = \Phi'$ and
$G'_j = \Phi' - \delta_{1/2} \Phi'$ for $j \geq 1$, its blocks are $\psi'_j = (G'_j)_{2^j}$, and
for every tempered distribution $u$ its blocks are $\Delta'_j u = (\widehat{\psi'_j})(D) u$. The
kernel $\Phi$ of \Cref{def:lp-kernels}, with $\widehat\Phi = \chi_{1,2}$, $a = 1$ and $b = 2$, is
the \emph{standard} admissible kernel.
\end{definition}

\begin{lemma}[The blocks of the standard kernel]
\label{lem:lp-kernel-standard-blocks}
\lean{Rd.LPKernel.standard_psiS, Rd.LPKernel.standard_distribBlock,
  Rd.LPKernel.standard_distribBlock_eq_iff}
\leanok
\uses{def:lp-kernel-admissible, def:psi-base, def:distrib-block}
For the standard kernel and every $j \in \N$, $\psi'_j = \psi_j$ and $\Delta'_j = \Delta_j$.
\end{lemma}

\begin{proof}
\leanok
The bases of the standard kernel are those of \Cref{def:psi-base}.
\end{proof}

\paragraph{The quasi-norm.}

\begin{definition}[Triebel--Lizorkin spaces of tempered distributions]
\label{def:tl-distr}
\lean{Rd.LPKernel.triebelLizorkinNormDistr, Rd.triebelLizorkinNormDistr,
  Rd.MemTriebelLizorkinDistr, Rd.LPKernel.standard_triebelLizorkinNormDistr,
  Rd.aemeasurable_lqNorm_dyadicWeight_Lp}
\leanok
\uses{def:lp-kernel-admissible, def:lq-norm}
Let $K$ be an admissible kernel, $1 \leq p \leq \infty$, $s \in \R$ and $1 \leq q \leq \infty$.
For every tempered distribution $u$ on $V$ let
\[
  \|u\|_{F^s_{p,q}, K} = \inf \Bigl\| x \mapsto \bigl\| (2^{js} |g_j(x)|)_{j \in \N}
  \bigr\|_{\ell_q} \Bigr\|_{L_p} ,
\]
the infimum over the families $(g_j)$ of functions in $L_p(V)$ with $\Delta'_j u = g_j$ as
tempered distributions for every $j$; it is $\infty$ if there is no such family. For the
standard kernel write $\|u\|_{F^s_{p,q}} = \|u\|_{F^s_{p,q}, K}$, with the blocks
$\Delta'_j = \Delta_j$. The space $\TL{s}{p}{q}$ consists of the tempered distributions $u$ with
$\|u\|_{F^s_{p,q}} < \infty$. For every family $(g_j)$ in $L_p(V)$ the function
$x \mapsto \| (2^{js} |g_j(x)|)_j \|_{\ell_q}$ is almost everywhere measurable.
\end{definition}

\begin{lemma}[The quasi-norm of the representing family]
\label{lem:tl-distr-eq}
\lean{Rd.LPKernel.triebelLizorkinNormDistr_le, Rd.LPKernel.triebelLizorkinNormDistr_eq,
  Rd.LPKernel.triebelLizorkinNormDistr_eq_top, Rd.LPKernel.le_mul_triebelLizorkinNormDistr,
  Rd.triebelLizorkinNormDistr_le, Rd.triebelLizorkinNormDistr_eq,
  Rd.triebelLizorkinNormDistr_eq_top, Rd.le_mul_triebelLizorkinNormDistr,
  Rd.MemTriebelLizorkinDistr.exists_blocks}
\leanok
\uses{def:tl-distr, lem:lp-distribution-determined, lem:lp-kernel-standard-blocks}
Let $K$ be an admissible kernel, $1 \leq p \leq \infty$, $s \in \R$, $1 \leq q \leq \infty$ and
$u$ a tempered distribution. For every family $(g_j)$ in $L_p(V)$ with $\Delta'_j u = g_j$ for
every $j$, $\|u\|_{F^s_{p,q}, K} = \| x \mapsto \| (2^{js} |g_j(x)|)_j \|_{\ell_q} \|_{L_p}$. If
there is no such family, $\|u\|_{F^s_{p,q}, K} = \infty$; so for every $u \in \TL{s}{p}{q}$ there
is such a family for the standard kernel. Hence, for every $X \in [0, \infty]$ and every
$C \in (0, \infty]$ such that $X \leq C \| x \mapsto \| (2^{js} |g_j(x)|)_j \|_{\ell_q} \|_{L_p}$
for every such family, $X \leq C \|u\|_{F^s_{p,q}, K}$.
\end{lemma}

\begin{proof}
\leanok
$L_p$ functions are determined by their tempered distributions
(\Cref{lem:lp-distribution-determined}), so there is at most one such family. Without a
family, $C \cdot \infty = \infty$ since $C \neq 0$. For the standard kernel the blocks are
$\Delta_j$ (\Cref{lem:lp-kernel-standard-blocks}).
\end{proof}

On $L_p$ each admissible kernel also gives a quasi-norm with the convolutions $\psi'_j * f$.

\begin{definition}[The quasi-norm of an admissible kernel on $L_p$]
\label{def:lp-kernel-tl-norm}
\lean{Rd.LPKernel.triebelLizorkinNorm}
\leanok
\uses{def:lp-kernel-admissible}
Let $K$ be an admissible kernel, $1 \leq p \leq \infty$, $s \in \R$ and $1 \leq q \leq \infty$.
For $f \in L_p(V)$,
\[
  \|f\|_{F^s_{p,q}, K} = \Bigl\| x \mapsto \bigl\| (2^{js} |(\psi'_j * f)(x)|)_j
  \bigr\|_{\ell_q} \Bigr\|_{L_p} .
\]
\end{definition}

\begin{lemma}[The quasi-norm of an admissible kernel on $L_p$]
\label{lem:lp-kernel-tl-norm-basic}
\lean{Rd.LPKernel.triebelLizorkinNormDistr_toTemperedDistribution,
  Rd.LPKernel.standard_triebelLizorkinNorm}
\leanok
\uses{def:lp-kernel-tl-norm, def:tl-distr, lem:tl-distr-eq, lem:lp-kernel-standard-blocks,
  lem:conv-fourier-multiplier, def:triebel-lizorkin}
Let $K$ be an admissible kernel, $1 \leq p \leq \infty$, $s \in \R$ and $1 \leq q \leq \infty$.
\begin{enumerate}
  \item If $p < \infty$, then for every $f \in L_p(V)$, with $T_f$ the tempered distribution of
    $f$, $\|T_f\|_{F^s_{p,q}, K} = \|f\|_{F^s_{p,q}, K}$, the left side from \Cref{def:tl-distr}
    and the right side from \Cref{def:lp-kernel-tl-norm}.
  \item For the standard kernel $K$ and every $f \in L_p(V)$,
    $\|f\|_{F^s_{p,q}, K} = \|f\|_{F^s_{p,q}}$, the quasi-norm of \Cref{def:triebel-lizorkin}.
\end{enumerate}
\end{lemma}

\begin{proof}
\leanok
(1) The blocks of $f$ are the $L_p$ functions $\psi'_j * f$ (\Cref{lem:conv-fourier-multiplier}),
and \Cref{lem:tl-distr-eq} applies to them. (2) For the standard kernel the blocks are the blocks
$\psi_j$ of \Cref{def:lp-kernels} (\Cref{lem:lp-kernel-standard-blocks}), so both quasi-norms are
the expression of \Cref{def:triebel-lizorkin}.
\end{proof}

\begin{theorem}[Compatibility with $L_p$]
\label{thm:tl-distr-lp}
\lean{Rd.triebelLizorkinNormDistr_toTemperedDistribution,
  Rd.memTriebelLizorkinDistr_toTemperedDistribution_iff}
\leanok
\uses{lem:lp-kernel-tl-norm-basic, def:triebel-lizorkin}
For every $1 \leq p < \infty$, $s \in \R$, $1 \leq q \leq \infty$ and $f \in L_p(V)$, the
quasi-norm of the tempered distribution of $f$ (\Cref{def:tl-distr}) is the quasi-norm
$\|f\|_{F^s_{p,q}}$ of \Cref{def:triebel-lizorkin}. In particular the two definitions of
$f \in \TL{s}{p}{q}$ agree.
\end{theorem}

\begin{proof}
\leanok
This is \Cref{lem:lp-kernel-tl-norm-basic} for the standard kernel.
\end{proof}

\begin{proposition}[Monotonicity in $q$ and $s$]
\label{prop:tl-distr-mono}
\lean{Rd.LPKernel.triebelLizorkinNormDistr_mono_exponent,
  Rd.LPKernel.triebelLizorkinNormDistr_mono_smoothness,
  Rd.triebelLizorkinNormDistr_mono_exponent, Rd.triebelLizorkinNormDistr_mono_smoothness,
  Rd.MemTriebelLizorkinDistr.mono_exponent, Rd.MemTriebelLizorkinDistr.mono_smoothness}
\leanok
\uses{def:tl-distr, lem:tl-distr-eq, lem:mixed-norm-dyadic-family}
Let $K$ be an admissible kernel, $1 \leq p \leq \infty$ and $u$ a tempered distribution. For all
$1 \leq q \leq r \leq \infty$ and every real $s$, $\|u\|_{F^s_{p,r}, K} \leq \|u\|_{F^s_{p,q}, K}$,
so $\TL{s}{p}{q} \subseteq \TL{s}{p}{r}$. For all reals $s \leq t$ and every $q$,
$\|u\|_{F^s_{p,q}, K} \leq \|u\|_{F^t_{p,q}, K}$, so $\TL{t}{p}{q} \subseteq \TL{s}{p}{q}$.
\end{proposition}

\begin{proof}
\leanok
For every representing family, $\|\cdot\|_{\ell_r} \leq \|\cdot\|_{\ell_q}$ and
$2^{js} \leq 2^{jt}$ at every point.
\end{proof}

\section{Triebel--Lizorkin spaces on the torus}
\label{sec:torus-triebel-lizorkin}

The Littlewood--Paley blocks on the torus are those of \Cref{def:torus-lp-blocks}, and
$\Besov{s}{p}{q}(\T^d)$ is the approximation space of the trigonometric polynomials in
$L_p(\T^d)$, $1 \leq p < \infty$.

\begin{definition}[Triebel--Lizorkin spaces on the torus]
\label{def:torus-triebel-lizorkin}
\lean{Torus.blockFun, Torus.lpBlock_ae_eq_blockFun, Torus.aestronglyMeasurable_blockFun,
  Torus.enorm_lpBlock_eq_eLpNorm_blockFun, Torus.triebelLizorkinNorm, Torus.MemTriebelLizorkin,
  Torus.aemeasurable_lqNorm_blockFun, Torus.littlewoodPaleyNorm_eq_lqNorm_eLpNorm}
\leanok
\uses{def:torus-lp-blocks, def:torus-econv, def:lq-norm}
Let $1 \leq p < \infty$, $s \in \R$ and $1 \leq q \leq \infty$. For $f \in L_p(\T^d)$ and
$j \in \N$ the pointwise block is $(\psi_j * f)(t) = \int_{\R^d} \psi_j(y) f(t - \pi y) \, dy$
at every $t \in \T^d$; it represents the block $\Delta_j f \in L_p(\T^d)$. The
Triebel--Lizorkin quasi-norm of $f$ is
\[
  \|f\|_{F^s_{p,q}} = \Bigl\| t \mapsto \bigl\| (2^{js} |(\psi_j * f)(t)|)_{j \in \N}
  \bigr\|_{\ell_q} \Bigr\|_{L_p(\T^d)},
\]
and $\TL{s}{p}{q}(\T^d)$ consists of the $f \in L_p(\T^d)$ with $\|f\|_{F^s_{p,q}} < \infty$.
\end{definition}

\begin{proposition}[Elementary properties on the torus]
\label{prop:torus-tl-elementary}
\lean{Torus.blockFun_add_ae, Torus.blockFun_smul_ae, Torus.eLpNorm_blockFun_zero_le,
  Torus.triebelLizorkinNorm_add_le,
  Torus.triebelLizorkinNorm_smul, Torus.triebelLizorkinNorm_zero,
  Torus.triebelLizorkinNorm_zero_left, Torus.triebelLizorkinNorm_mono_exponent,
  Torus.triebelLizorkinNorm_mono_smoothness,
  Torus.MemTriebelLizorkin.mono_exponent, Torus.MemTriebelLizorkin.mono_smoothness,
  Torus.MemTriebelLizorkin.add, Torus.MemTriebelLizorkin.smul, Torus.triebelLizorkinSpace,
  Torus.mem_triebelLizorkinSpace}
\leanok
\uses{def:torus-triebel-lizorkin, lem:torus-econv-young, lem:lp-kernels-basic, lem:lq-basic,
  lem:lp-norm-homogeneity, lem:mixed-norm-dyadic-family}
For every $1 \leq p < \infty$, all $s \leq t$ in $\R$, all $1 \leq q \leq r \leq \infty$, all
$f, g \in L_p(\T^d)$ and every $c \in \C$, for almost every $x \in \T^d$ and all $j \in \N$,
\[
  (\psi_j * (f + g))(x) = (\psi_j * f)(x) + (\psi_j * g)(x), \qquad
  (\psi_j * (c f))(x) = c \, (\psi_j * f)(x),
\]
$\|\psi_0 * f\|_p \leq \|\Phi\|_1 \, \|f\|_p$, and
\[
  \|f + g\|_{F^s_{p,q}} \leq \|f\|_{F^s_{p,q}} + \|g\|_{F^s_{p,q}}, \qquad
  \|c f\|_{F^s_{p,q}} = |c| \, \|f\|_{F^s_{p,q}},
\]
\[
  \|f\|_{F^s_{p,r}} \leq \|f\|_{F^s_{p,q}}, \qquad \|f\|_{F^s_{p,q}} \leq \|f\|_{F^t_{p,q}} .
\]
In particular $\TL{s}{p}{q}(\T^d)$ is a linear subspace of $L_p(\T^d)$,
$\TL{s}{p}{q}(\T^d) \subseteq \TL{s}{p}{r}(\T^d)$ and
$\TL{t}{p}{q}(\T^d) \subseteq \TL{s}{p}{q}(\T^d)$. For every $q \in [0, \infty]$ and every
$f \in L_p(\T^d)$,
$\|f\|_{F^0_{p,q}} = \bigl\| t \mapsto \|((\psi_j * f)(t))_j\|_{\ell_q} \bigr\|_{L_p(\T^d)}$.
\end{proposition}

\begin{proof}
\leanok
The pointwise block depends on $f$ only through its class, and the convolution integral is
linear in $f$. Since $L_p(\T^d) \subseteq L_1(\T^d)$, the integrals converge absolutely at almost
every point (\Cref{lem:torus-econv-young}), which gives the identities for the pointwise blocks,
for all $j$ at once. Young's inequality on the torus (\Cref{lem:torus-econv-young}) and
$\|\psi_0\|_1 = \|\Phi\|_1$ (\Cref{lem:lp-kernels-basic}) give the bound of $\psi_0 * f$. The rest
is the proof of \Cref{prop:tl-elementary}; for $s = 0$ the weights
$2^{j \cdot 0}$ are $1$.
\end{proof}

\begin{lemma}[The weighted blocks on the torus and their tails]
\label{lem:torus-weighted-block-tails}
\lean{Torus.eLpNorm_lqNorm_ite_blockFun_le, Torus.tendsto_eLpNorm_lqNorm_blockFun_tail}
\leanok
\uses{def:torus-triebel-lizorkin, lem:mixed-norm-tails}
Let $1 \leq p < \infty$, $s \in \R$, $1 \leq q \leq \infty$ and $f \in L_p(\T^d)$. For every set
$P \subseteq \N$,
\[
  \bigl\| \|(\mathbf 1_P(j) \, 2^{js} |(\psi_j * f)(t)|)_j\|_{\ell_q} \bigr\|_p
  \leq \|f\|_{\TL{s}{p}{q}(\T^d)} .
\]
If $q < \infty$ and $f \in \TL{s}{p}{q}(\T^d)$, then
\[
  \bigl\| \|(2^{js} |(\psi_j * f)(t)|)_{j \geq N}\|_{\ell_q} \bigr\|_p \to 0 \qquad
  (N \to \infty) ,
\]
where the $\ell_q$ norm is taken of the sequence with the terms $j < N$ replaced by $0$.
\end{lemma}

\begin{proof}
\leanok
Both claims are \Cref{lem:mixed-norm-tails} for the functions
$g_j = 2^{js} |\psi_j * f|$ on $\T^d$, whose mixed norm is $\|f\|_{\TL{s}{p}{q}(\T^d)}$
(\Cref{def:torus-triebel-lizorkin}).
\end{proof}

\begin{theorem}[$\TL{s}{p}{p}(\T^d) = \Besov{s}{p}{p}(\T^d)$ and the Besov sandwich on the torus]
\label{thm:torus-tl-besov-sandwich}
\lean{Torus.triebelLizorkinNorm_eq_littlewoodPaleyNorm, Torus.memTriebelLizorkin_iff_memBesov,
  Torus.memTriebelLizorkin_iff_modulus, Torus.littlewoodPaleyNorm_max_le_triebelLizorkinNorm,
  Torus.triebelLizorkinNorm_le_littlewoodPaleyNorm_min,
  Torus.approxNormDyadic_le_triebelLizorkinNorm, Torus.triebelLizorkinNorm_le_approxNormDyadic,
  Torus.memBesov_of_memTriebelLizorkin, Torus.memTriebelLizorkin_of_memBesov}
\leanok
\uses{def:torus-triebel-lizorkin, cor:mixed-norm-sandwich, thm:torus-lp-besov, def:dyadic-norm,
  thm:torus-lp-moduli}
For every $1 \leq p < \infty$, every $s \in \R$, every $1 \leq q \leq \infty$ and every
$f \in L_p(\T^d)$,
\[
  \|f\|^{\mathrm{LP}}_{s,\max(p,q)} \leq \|f\|_{F^s_{p,q}} \leq \|f\|^{\mathrm{LP}}_{s,\min(p,q)} ,
  \qquad \|f\|_{F^s_{p,p}} = \|f\|^{\mathrm{LP}}_{s,p} .
\]
Let $A$ be the scheme of trigonometric polynomials in $L_p(\T^d)$. If $s \geq 0$, then
\[
  \|f\|_{F^s_{p,q}} \leq 2 (\#d + 1) \, 2^{(\#d + 1) s} \, \|\Phi\|_1 \,
  \|f\|^{\mathrm{dy}}_{A^s_{\min(p,q)}} .
\]
If $s > 0$, then
\[
  \|f\|^{\mathrm{dy}}_{A^s_{\max(p,q)}} \leq 2 \cdot 2^s (1 - 2^{-s})^{-1} \, \|f\|_{F^s_{p,q}} ,
\]
$f \in \TL{s}{p}{p}(\T^d)$ if and only if $f \in \Besov{s}{p}{p}(\T^d)$, for every $r \in \N$ with
$s < r$, $f \in \TL{s}{p}{p}(\T^d)$ if and only if
$\|f\|_p + \|(2^{ks} \omega_r(f, 2^{-k})_p)_k\|_{\ell_p} < \infty$, and
\[
  \Besov{s}{p}{\min(p,q)}(\T^d) \subseteq \TL{s}{p}{q}(\T^d)
  \subseteq \Besov{s}{p}{\max(p,q)}(\T^d) .
\]
\end{theorem}

\begin{proof}
\leanok
The $L_p$ norm of the pointwise block $\psi_j * f$ is $\|\Delta_j f\|_p$, so the bounds are
\Cref{cor:mixed-norm-sandwich} for the family $g_j = 2^{js} |\psi_j * f|$ on $\T^d$. The
characterisation of $\Besov{s}{p}{q}(\T^d)$ by Littlewood--Paley blocks
(\Cref{thm:torus-lp-besov}) at the exponents $p$, $\max(p,q)$ and $\min(p,q)$ gives the
bounds by the approximation norms and the statements on spaces; the description by moduli of
smoothness is \Cref{thm:torus-lp-moduli} with $q = p$.
\end{proof}

As on $\R^n$, for $q = 2$ the theorem places $\TL{s}{p}{2}(\T^d)$ between
$\Besov{s}{p}{\min(p,2)}(\T^d)$ and $\Besov{s}{p}{\max(p,2)}(\T^d)$.

\section{Kernel majorants}

The characterisations of the spaces $\TL{s}{p}{q}$ compare two kinds of blocks pointwise: a block of
one kind is a sum of blocks of the other kind, each convolved with a cross kernel. This section
bounds the cross kernels by the Peetre weights $2^{\nu n} (1 + 2^\nu \|y\|)^{-N}$ times a factor
$c_{ji}$, a power $2^{-(j-i)L}$ that decays geometrically in $j - i$ for $i \leq j$ and $1$ for
$i > j$, and such kernels act below the maximal function. As throughout this chapter, $V$ has
dimension $n$; the kernels $\Phi$, $\Theta$, $\psi_j$ are those of \Cref{def:lp-kernels}, $g$ is
a bump of radius $R$ with the kernels $m$, $k_j$, $\eta_j$ of \Cref{def:local-means}, and every
constant is explicit.

\begin{lemma}[Peetre weights]
\label{lem:peetre-weights}
\lean{ENNReal.inv_le_mul_inv_of_le, ENNReal.one_add_ne_zero, ENNReal.inv_one_add_mul_pow_le,
  Rd.inv_one_add_mul_enorm_sub_pow_le, ENNReal.inv_one_add_mul_pow_anti,
  ENNReal.inv_one_add_mul_pow_le_two_pow, ENNReal.inv_one_add_pow_mul_pow_le,
  Rd.ofReal_one_add_norm_rpow_neg, Rd.ofReal_inv_one_add_mul_norm_pow,
  Rd.two_pow_mul_enorm_le_ofReal, Rd.ofReal_one_add_two_pow_mul_norm,
  ENNReal.one_le_ofReal_mul_inv_of_le, Rd.inv_one_add_ofReal_half_pow_le,
  Rd.inv_one_add_ofReal_div_pow_le, ENNReal.inv_one_add_pow_two_mul_mul_pow}
\leanok
For every $N \in \N$, all $i, j, j', L \in \N$, all $0 \leq a \leq b < \infty$, every
$u \in [0, \infty]$ and all $y, w \in V$:
\begin{enumerate}
  \item $(1 + a \|y - w\|)^{-N} \leq (1 + a \|w\|)^N (1 + a \|y\|)^{-N}$;
  \item $(1 + b u)^{-N} \leq (1 + a u)^{-N}$ and $(1 + a u)^{-N} \leq 2^N (1 + 2 a u)^{-N}$;
  \item $(1 + a)^{-(k + r + m)} a^r \leq (1 + a)^{-k} (1 + a)^{-m}$ and
    $(1 + a)^{-2N} (1 + a)^N = (1 + a)^{-N}$ for all $k, r, m \in \N$;
  \item if $a \leq b$, then $1 \leq (1 + b)^N (1 + a)^{-N}$;
  \item $(1 + 2^i \|y\| / 2)^{-N} \leq 2^N (1 + 2^j \|y\|)^{-N}$ for all $j \leq i$;
  \item $(1 + 2^{j'} \|w\| / (2 L))^{-N} \leq (8 L)^N (1 + 2^i \|w\|)^{-N}$ for all
    $L \geq 1$ and $i \leq j' + 2$.
\end{enumerate}
\end{lemma}

\begin{proof}
\leanok
(1) follows from $\|y\| \leq \|w\| + \|y - w\|$ and
$1 + a (s + t) \leq (1 + a s)(1 + a t)$; (5) from $1 + 2^j \|y\| \leq 2 (1 + 2^i \|y\| / 2)$, and
(6) from $1 + 2^i \|w\| \leq 8 L + 4 \cdot 2^{j'} \|w\| = 8 L (1 + 2^{j'} \|w\| / (2 L))$. The
other claims are elementary. Each inequality $X^{-1} \leq c Y^{-1}$ follows from $Y \leq c X$
for $0 < c < \infty$.
\end{proof}

\begin{lemma}[Decay of Schwartz functions and of their differences]
\label{lem:schwartz-decay}
\lean{Rd.schwartzDecayConst, Rd.schwartzDecayConst_ne_top, Rd.one_add_norm_pow_mul_le_schwartz,
  Rd.one_add_norm_le_mul_of_mem_Icc, Rd.norm_iteratedFDeriv_schwartz_le_of_mem_Icc,
  Rd.enorm_higherDiff_schwartz_le, Rd.enorm_schwartz_le, Rd.enorm_two_pow_smul,
  Rd.enorm_dilate_two_pow, Rd.enorm_dilate_schwartz_le, Rd.enorm_higherDiff_dilate_schwartz_le}
\leanok
\uses{lem:differences-segment, lem:peetre-weights, lem:dilate, lem:rd-differences,
  thm:band-jackson}
For every Schwartz function $G : V \to \C$ and all $r, N \in \N$ let
\[
  C_{G,r,N} = 2^N \Bigl( \sup_x \|D^r G(x)\| + \sup_x \|x\|^N \|D^r G(x)\| \Bigr) < \infty .
\]
For all $v, y, h \in V$ and $j \in \N$,
\[
  |\Delta_v^r G(y)| \leq C_{G,r,N} \|v\|^r (1 + r \|v\|)^N (1 + \|y\|)^{-N},
\]
\[
  |\Delta_h^r G_{2^j}(y)| \leq C_{G,r,N} \|2^j h\|^r (1 + r \|2^j h\|)^N \, 2^{jn}
  (1 + 2^j \|y\|)^{-N} .
\]
\end{lemma}

\begin{proof}
\leanok
$(1 + \|x\|)^N \leq 2^N (1 + \|x\|^N)$ gives $\|D^r G(x)\| \leq C_{G,r,N} (1 + \|x\|)^{-N}$
(Mathlib's \texttt{SchwartzMap.\allowbreak one\_\allowbreak add\_\allowbreak le\_\allowbreak
sup\_\allowbreak seminorm\_\allowbreak apply} is the same kind of bound with a supremum over
seminorms), and $1 + \|y\| \leq (1 + \|y + t v\|)(1 + r \|v\|)$ for $0 \leq t \leq r$ moves the
weight from $y + t v$ to $y$; \Cref{lem:differences-segment} concludes. For the dilations,
$|G_{2^j}(y)| = 2^{jn} |G(2^j y)|$ and $\Delta_h^r G_{2^j} = (\Delta_{2^j h}^r G)_{2^j}$
(\Cref{lem:rd-differences}).
\end{proof}

With $r = 0$ the two bounds are the decay $|G(y)| \leq C_{G,0,N} (1 + \|y\|)^{-N}$ and
$|G_{2^j}(y)| \leq C_{G,0,N} \, 2^{jn} (1 + 2^j \|y\|)^{-N}$.

\begin{lemma}[Tails and changes of scale]
\label{lem:schwartz-tails}
\lean{Rd.peetreIntegral, Rd.peetreIntegral_ne_top, Rd.tailConst, Rd.tailConst_ne_top,
  Rd.enorm_mul_enorm_pow_le_of_le, Rd.lintegral_tail_schwartz_le, Rd.lintegral_dilate_peetre,
  Rd.setLIntegral_enorm_dilate_schwartz_le}
\leanok
\uses{lem:schwartz-decay, lem:peetre-weights, lem:change-of-scale}
For every $M \in \N$ let $I_M = \int_V (1 + \|u\|)^{-M} \, du$; then $I_M < \infty$ for $M > n$.
For every Schwartz function $G$ and all $r, N \in \N$ let
$T_{G,r,N} = C_{G,0,N+r+n+1} \, I_{n+1} < \infty$, and write $t_+ = \max(t, 0)$. For all
$\rho, a \in \R$ and $j \in \N$,
\[
  \int_{\{\|z\| \geq \rho\}} |G(z)| \, \|z\|^r \, dz \leq T_{G,r,N} (1 + \rho_+)^{-N},
\]
\[
  \int_{\{\|z\| \geq a\}} |G_{2^j}(z)| \, dz \leq T_{G,0,N} (1 + (2^j a)_+)^{-N},
\]
and $\int_V 2^{jn} (1 + 2^j \|u\|)^{-N} \, du = I_N$.
\end{lemma}

\begin{proof}
\leanok
$I_M < \infty$ is Mathlib's \texttt{finite\_integral\_one\_add\_norm}. On $\{\|z\| \geq \rho\}$,
\Cref{lem:schwartz-decay} and part (3) of \Cref{lem:peetre-weights} give
$|G(z)| \|z\|^r \leq C_{G,0,N+r+n+1} (1 + \rho_+)^{-N} (1 + \|z\|)^{-(n+1)}$. The change of scale
with $c = 2^j$ (\Cref{lem:change-of-scale}) turns the tail of $G_{2^j}$ beyond $a$ into the tail
of $G$ beyond $2^j a$.
\end{proof}

\begin{lemma}[Two Peetre bumps]
\label{lem:two-decay-bumps}
\lean{Rd.peetre_mul_peetre_le, Rd.lintegral_peetre_mul_le}
\leanok
\uses{lem:peetre-weights, lem:schwartz-tails}
For all $i \leq j$ in $\N$, every $N \in \N$ and every $y \in V$,
\[
  \int_V 2^{jn} (1 + 2^j \|y - w\|)^{-2N} \, 2^{in} (1 + 2^i \|w\|)^{-N} \, dw
  \leq I_N \, 2^{in} (1 + 2^i \|y\|)^{-N} .
\]
\end{lemma}

\begin{proof}
\leanok
By part (1) of \Cref{lem:peetre-weights} and $2^i \leq 2^j$,
$(1 + 2^i \|w\|)^{-N} \leq (1 + 2^j \|y - w\|)^N (1 + 2^i \|y\|)^{-N}$, which leaves
$\int_V 2^{jn} (1 + 2^j \|y - w\|)^{-N} \, dw = I_N$ by the substitution $u = 2^j (y - w)$.
\end{proof}

The bound holds for every $N$; it is useful for $N > n$, where $I_N < \infty$.

\begin{lemma}[Dyadic weights of integer order]
\label{lem:dyadic-weights-integer}
\lean{ENNReal.dyadicWeight_natCast, ENNReal.dyadicWeight_neg_natCast,
  ENNReal.dyadicWeight_neg_sub_le, ENNReal.one_le_two_pow_mul_dyadicWeight_neg,
  ENNReal.dyadicWeight_mul_dyadicWeight_neg_eq}
\leanok
\uses{def:lq-norm}
For all $L, i, j', k, m \in \N$: $2^{-mL} = (2^m)^{-L}$; if $i \leq j' \leq k$, then
$2^{-(k-i)L} \leq 2^L \, 2^{-(j'+1-i)L}$; if $m \leq 1$, then $1 \leq 2^L 2^{-mL}$; and if
$i \leq j' + 2$, then $2^{iL} 2^{-j'L} = 4^L \, 2^{-(j'+2-i)L}$.
\end{lemma}

\begin{proof}
\leanok
The exponents add: $2^{-(j'+2-i)L} 2^{-iL} = 2^{-j'L} 2^{-2L}$.
\end{proof}

\begin{lemma}[Convolution in the second argument and averaged differences]
\label{lem:conv-second-argument}
\lean{Rd.conv_sub_ae, Rd.conv_finsetSum_ae, Rd.ae_enorm_sub_conv_dilate_momentKernel_le}
\leanok
\uses{lem:conv-assoc, lem:averaged-differences}
Let $1 \leq p < \infty$ and $K \in L_1$. For all $f, f' \in L_p$ and every finite family
$(f_i)_{i \in S}$ in $L_p$, almost everywhere on $V$,
\[
  K * (f - f') = K * f - K * f', \qquad K * \sum_{i \in S} f_i = \sum_{i \in S} K * f_i .
\]
For every $L \in \N$, every $g \in L_1$ with $\int_V g = 1$ and its moment kernel $m^g$ of
order $L$, every $c > 0$, every $f \in L_p$ and almost every $x \in V$,
\[
  |f(x) - (m^g_c * f)(x)| \leq \int_V |g(z)| \, |\Delta^L_{-c^{-1} z} f(x)| \, dz .
\]
\end{lemma}

\begin{proof}
\leanok
The first two identities follow from the additivity in \Cref{lem:conv-assoc}, by induction over
$S$ for the sum. The bound is \Cref{lem:averaged-differences} with the triangle inequality for
the integral.
\end{proof}

\begin{lemma}[Bounds of the local means kernels]
\label{lem:local-means-kernel-bounds}
\lean{Rd.LocalMeansBump.kernelSupConst, Rd.LocalMeansBump.kernelSupConst_nonneg,
  Rd.LocalMeansBump.norm_scaled_le, Rd.LocalMeansBump.norm_kernel_le,
  Rd.LocalMeansBump.kernel_eq_zero_of_lt, Rd.LocalMeansBump.enorm_kernel_le}
\leanok
\uses{def:local-means, lem:dual-kernel-bounds}
Let $L \in \N$ and $A_L = 2 \sup_V |m|$. For all $j \in \N$ and $y \in V$,
$|m_{2^j}(y)| \leq 2^{jn} \sup_V |m|$ and $|k_j(y)| \leq A_L 2^{jn}$, and $k_j(y) = 0$ if
$\|y\| > 2 L R \, 2^{-j}$.
\end{lemma}

\begin{proof}
\leanok
$m_{2^j}(y) = 2^{jn} m(2^j y)$, and $k_j$ is a difference of two such terms, each supported in
$\overline B(0, L R \, 2^{-(j-1)})$ by \Cref{lem:dual-kernel-bounds}.
\end{proof}

\begin{lemma}[Local means kernels from one scale]
\label{lem:local-means-one-scale}
\lean{Rd.LocalMeansBump.ae_enorm_conv_kernel_le_of_scaled}
\leanok
\uses{def:local-means, lem:conv-assoc, lem:dyadic-weights-integer}
For every $L \in \N$, all $j < i$, every $f \in L_1$, every $C \in [0, \infty]$ and every
$u : V \to [0, \infty]$ such that for every $k \geq j$ almost everywhere
$|f(x) - (m_{2^k} * f)(x)| \leq C \, 2^{-(k-j)L} u(x)$, almost everywhere
\[
  |(k_i * f)(x)| \leq 2 \cdot 2^L \, C \, 2^{-(i-j)L} u(x) .
\]
\end{lemma}

\begin{proof}
\leanok
Since $i \geq 1$, $k_i = m_{2^i} - m_{2^{i-1}}$, so
$k_i * f = (f - m_{2^{i-1}} * f) - (f - m_{2^i} * f)$ almost everywhere (\Cref{lem:conv-assoc}),
and $2^{-(k-j)L} \leq 2^L \, 2^{-(i-j)L}$ for $k \in \{i - 1, i\}$
(\Cref{lem:dyadic-weights-integer}).
\end{proof}

\begin{theorem}[The cross kernel of local means and band-limiting kernels]
\label{thm:cross-kernel-theta}
\lean{Rd.LocalMeansBump.enorm_mul_higherDiff_Theta_le,
  Rd.LocalMeansBump.ae_enorm_Theta_sub_conv_scaled_le, Rd.LocalMeansBump.kernelThetaMomentConst,
  Rd.LocalMeansBump.kernelThetaMomentConst_ne_top,
  Rd.LocalMeansBump.ae_enorm_conv_kernel_Theta_le_of_lt,
  Rd.LocalMeansBump.kernelThetaSupportConst, Rd.LocalMeansBump.kernelThetaSupportConst_ne_top,
  Rd.LocalMeansBump.lintegral_Theta_mul_kernel_le_near,
  Rd.LocalMeansBump.lintegral_Theta_mul_kernel_le_far,
  Rd.LocalMeansBump.enorm_conv_kernel_Theta_le_of_le, Rd.LocalMeansBump.kernelThetaConst,
  Rd.LocalMeansBump.kernelThetaConst_ne_top, Rd.LocalMeansBump.enorm_conv_kernel_Theta_le}
\leanok
\uses{lem:local-means-kernel-bounds, lem:local-means-one-scale, lem:conv-second-argument,
  lem:schwartz-decay, lem:schwartz-tails, lem:peetre-weights, lem:dyadic-weights-integer,
  lem:lp-kernels-basic, lem:conv-assoc, lem:weighted-jensen}
Let $L, N \in \N$, let
\[
  C^{\mathrm m}_{L,N} = 2 \|g\|_1 (2R)^L (1 + 2LR)^N C_{\Theta,L,N},
\]
\[
  C^{\mathrm s}_{L,N} = A_L \bigl( \|\Theta\|_1 (1 + 4LR)^N + 2^N T_{\Theta,0,N} \bigr),
\]
and $C^\Theta_{L,N} = C^{\mathrm m}_{L,N} + C^{\mathrm s}_{L,N} < \infty$.
For all $i, j \in \N$, with $\nu = \min(i, j)$, $c_{ji} = 2^{-(j-i)L}$ for $i \leq j$ and
$c_{ji} = 1$ for $i > j$, and for almost every $y \in V$,
\[
  |(k_j * \Theta_i)(y)| \leq C^\Theta_{L,N} \, c_{ji} \, 2^{\nu n} (1 + 2^\nu \|y\|)^{-N} .
\]
\end{theorem}

\begin{proof}
\leanok
\emph{Moment case $i < j$.} By \Cref{lem:local-means-one-scale} it suffices to bound
$|\Theta_i - m_c * \Theta_i|$ for $c = 2^k$, $k \geq i$. \Cref{lem:conv-second-argument} bounds
$|\Theta_i - m_c * \Theta_i|(y)$ by $\int |g(z)| |\Delta^L_{-2^{-k} z} \Theta_i(y)| \, dz$. On the
support of $g$ the step $2^i \cdot 2^{-k} z$ of the dilated difference has length at most
$2^{-(k-i)} R \leq R$, and \Cref{lem:schwartz-decay} for $G = \Theta$ gives
$|\Delta^L_{-2^{-k} z} \Theta_i(y)| \leq C_{\Theta,L,N} R^L 2^{-(k-i)L} (1 + LR)^N 2^{in}
(1 + 2^i \|y\|)^{-N}$; \Cref{lem:dyadic-weights-integer} replaces $2^{-(k-i)L}$ by
$2^L 2^{-(j-i)L}$.

\emph{Support case $i \geq j$.} By commutativity,
$|(k_j * \Theta_i)(y)| \leq \int |\Theta_i(z)| |k_j(y - z)| \, dz$, and
$|k_j| \leq A_L 2^{jn}$ vanishes outside $\overline B(0, 2LR \, 2^{-j})$
(\Cref{lem:local-means-kernel-bounds}). For $\|y\| \leq 4LR \, 2^{-j}$ the integral is at most
$A_L 2^{jn} \|\Theta\|_1$, and $1 \leq (1 + 4LR)^N (1 + 2^j \|y\|)^{-N}$. For
$\|y\| > 4LR \, 2^{-j}$ the integrand vanishes unless $\|z\| \geq \|y\|/2$, and
\Cref{lem:schwartz-tails} bounds the tail of $\Theta_i$ by
$T_{\Theta,0,N} (1 + 2^i \|y\|/2)^{-N} \leq 2^N T_{\Theta,0,N} (1 + 2^j \|y\|)^{-N}$ (part (5)
of \Cref{lem:peetre-weights}). For $i = j$ the weight is $c_{jj} = 1$.
\end{proof}

The moment case uses only $\int g = 1$, through the averaged differences, so the bound holds
for every bump $g$ and every order $L$.

\begin{lemma}[The moment kernels of the Littlewood--Paley kernel]
\label{lem:lp-moment-kernels}
\lean{Rd.fourier_dilate_momentKernel, Rd.continuous_momentKernel, Rd.bddAbove_momentKernel,
  Rd.lpMoment, Rd.lpMomentScaled, Rd.integrable_lpMomentScaled, Rd.continuous_lpMomentScaled,
  Rd.bddAbove_lpMomentScaled, Rd.fourier_lpMomentScaled_eq_zero,
  Rd.conv_lpBlockKernel_lpMomentScaled}
\leanok
\uses{def:moment-kernel, lem:fourier-dilate, lem:lp-kernels-basic, lem:lp-kernels-frequency,
  lem:conv-of-fourier}
For every integrable $h : V \to \C$, every $L \in \N$, every $c > 0$ and every $\xi \in V$, the
moment kernel $m^h$ of order $L$ of $h$ satisfies
\[
  \widehat{(m^h)_c}(\xi) = \sum_{l=1}^L (-1)^{l+1} \binom{L}{l} \, \widehat h(l c^{-1} \xi) .
\]
Let $\mu$ be the moment kernel of order $L$ of $\Phi$. For all $j \in \N$ and $\xi \in V$ with
$\|\xi\| \geq 2^{j+1}$, $\widehat{\mu_{2^j}}(\xi) = 0$, and $\psi_{j+2} * \mu_{2^j} = 0$.
\end{lemma}

\begin{proof}
\leanok
The formula follows from \Cref{def:moment-kernel} and \Cref{lem:fourier-dilate}, term by term.
For $\|\xi\| \geq 2^{j+1}$ and $l \geq 1$, $\|l \, 2^{-j} \xi\| \geq 2$, where
$\widehat\Phi = \chi_{1,2}$ vanishes. Since $\widehat{\psi_{j+2}}$ vanishes on the ball of
radius $2^{j+1}$ (\Cref{lem:lp-kernels-frequency}), $\widehat{\psi_{j+2}} \widehat{\mu_{2^j}} = 0$,
and \Cref{lem:conv-of-fourier} applies to the bounded continuous kernel $\mu_{2^j}$.
\end{proof}

\begin{lemma}[Decay of the Littlewood--Paley blocks]
\label{lem:lp-blocks-decay}
\lean{Rd.psiDecayConst, Rd.psiDecayConst_ne_top, Rd.enorm_lpBlockKernel_le}
\leanok
\uses{lem:schwartz-decay, lem:peetre-weights, def:lp-kernels}
For every $N \in \N$ let $P_N = (1 + 2^N) C_{\Phi,0,N} < \infty$. For every $j \in \N$ and every
$y \in V$, $|\psi_j(y)| \leq P_N 2^{jn} (1 + 2^j \|y\|)^{-N}$.
\end{lemma}

\begin{proof}
\leanok
$\psi_j = \Phi_j - \Phi_{j-1}$ for $j \geq 1$; \Cref{lem:schwartz-decay} bounds both dilations,
and $(1 + 2^{j-1} \|y\|)^{-N} \leq 2^N (1 + 2^j \|y\|)^{-N}$.
\end{proof}

\begin{lemma}[Kernels of small support against Peetre weights]
\label{lem:small-support-decay}
\lean{Rd.lintegral_mul_peetre_le_of_support}
\leanok
\uses{lem:peetre-weights}
Let $k, N \in \N$, $\rho \geq 0$, $C \in [0, \infty]$, and let $\eta, \psi : V \to \C$ with
$\eta$ almost everywhere strongly measurable, $\eta = 0$ outside $\overline B(0, \rho 2^{-k})$
and $|\psi(z)| \leq C 2^{kn} (1 + 2^k \|z\|)^{-N}$ for all $z$. For every $y \in V$,
\[
  \int_V |\eta(w)| \, |\psi(y - w)| \, dw \leq \|\eta\|_1 \, C (1 + \rho)^N 2^{kn}
  (1 + 2^k \|y\|)^{-N} .
\]
\end{lemma}

\begin{proof}
\leanok
On the support of $\eta$, $2^k \|w\| \leq \rho$, and part (1) of \Cref{lem:peetre-weights} gives
$(1 + 2^k \|y - w\|)^{-N} \leq (1 + \rho)^N (1 + 2^k \|y\|)^{-N}$.
\end{proof}

\begin{lemma}[Averaged differences of the dual kernels]
\label{lem:dual-averaged-differences}
\lean{Rd.LocalMeansBump.enorm_higherDiff_dual_le, Rd.LocalMeansBump.higherDiff_dual_eq_zero,
  Rd.LocalMeansBump.lintegral_lpKernel_mul_higherDiff_dual_le, Rd.LocalMeansBump.dualTailConst,
  Rd.LocalMeansBump.dualTailConst_ne_top, Rd.LocalMeansBump.higherDiff_dual_eq_zero_of_norm_lt,
  Rd.LocalMeansBump.ae_enorm_dual_sub_conv_lpMomentScaled_le}
\leanok
\uses{lem:dual-kernel-derivatives, lem:dual-kernel-bounds, lem:smooth-differences,
  lem:conv-second-argument, lem:schwartz-tails, lem:peetre-weights, lem:dyadic-weights-integer,
  lem:lp-moment-kernels, lem:lp-kernels-basic}
Let $L, N \in \N$, $L' \geq 1$, let $\mu$ be the moment kernel of order $L'$ of $\Phi$, let
$A_{L,L'}$ be the constant of \Cref{lem:dual-kernel-derivatives} for the order $K = L'$, and
\[
  D_{L,L',N} = A_{L,L'} \, 4^{L'} (8 L')^N (1 + 4LR)^N \, T_{\Phi,L',N} < \infty .
\]
For all $i, j' \in \N$ with $i \leq j' + 2$ and almost every $w \in V$,
\[
  |\eta_i(w) - (\mu_{2^{j'}} * \eta_i)(w)| \leq D_{L,L',N} \, 2^{-(j'+2-i)L'} \, 2^{in}
  (1 + 2^i \|w\|)^{-N} .
\]
\end{lemma}

\begin{proof}
\leanok
By \Cref{lem:conv-second-argument} the left side is at most
$\int |\Phi(z)| |\Delta^{L'}_{-2^{-j'} z} \eta_i(w)| \, dz$. By \Cref{lem:smooth-differences} and
\Cref{lem:dual-kernel-derivatives},
$|\Delta^{L'}_{-2^{-j'} z} \eta_i(w)| \leq A_{L,L'} 2^{in} 2^{iL'} 2^{-j'L'} \|z\|^{L'}$, and
the difference vanishes if $\|w\| > 2LR \, 2^{-i} + L' 2^{-j'} \|z\|$, since $\eta_i$ vanishes
outside $\overline B(0, 2LR \, 2^{-i})$. For $\|w\| \leq 4LR \, 2^{-i}$, \Cref{lem:schwartz-tails}
with $\rho = 0$ bounds the integral by $T_{\Phi,L',N}$, and
$1 \leq (1 + 4LR)^N (1 + 2^i \|w\|)^{-N}$. For $\|w\| > 4LR \, 2^{-i}$ only
$\|z\| \geq 2^{j'} \|w\| / (2L')$ contributes, and the tail bound with part (6) of
\Cref{lem:peetre-weights} gives the factor $(8L')^N (1 + 2^i \|w\|)^{-N}$. Finally
$2^{iL'} 2^{-j'L'} = 4^{L'} 2^{-(j'+2-i)L'}$ (\Cref{lem:dyadic-weights-integer}).
\end{proof}

\begin{theorem}[The cross kernel of blocks and dual kernels]
\label{thm:cross-kernel-dual}
\lean{Rd.LocalMeansBump.psiDualSupportConst, Rd.LocalMeansBump.psiDualSupportConst_ne_top,
  Rd.LocalMeansBump.lintegral_enorm_dual_le,
  Rd.LocalMeansBump.enorm_conv_lpBlockKernel_dualKernel_le_of_lt,
  Rd.LocalMeansBump.enorm_conv_lpBlockKernel_dualKernel_le_of_le_one,
  Rd.LocalMeansBump.conv_lpBlockKernel_dualKernel_ae_eq,
  Rd.LocalMeansBump.ae_enorm_conv_lpBlockKernel_dualKernel_le_of_two_le,
  Rd.LocalMeansBump.psiDualConst, Rd.LocalMeansBump.psiDualConst_ne_top,
  Rd.LocalMeansBump.enorm_conv_lpBlockKernel_dualKernel_le}
\leanok
\uses{lem:lp-blocks-decay, lem:small-support-decay, lem:dual-averaged-differences,
  lem:lp-moment-kernels, lem:two-decay-bumps, lem:dual-kernel-bounds, lem:dual-kernel-moduli,
  lem:dyadic-weights-integer, lem:conv-assoc, lem:conv-second-argument, lem:weighted-jensen}
Let $L, N \in \N$ and $L' \geq 1$, and let
\[
  S_{L,N} = 2 P_N (1 + 2LR)^N \|m\|_1,
\]
\[
  C_{L,L',N} = (1 + 2^{n + L'}) S_{L,N} + D_{L,L',N} \, P_{2N} \, I_N ,
\]
with $C_{L,L',N} < \infty$ for $N > n$. For all $i, j \in \N$, with $\nu = \min(i, j)$,
$c_{ji} = 2^{-(j-i)L'}$ for $i \leq j$ and $c_{ji} = 1$ for $i > j$, and for almost every
$y \in V$,
\[
  |(\psi_j * \eta_i)(y)| \leq C_{L,L',N} \, c_{ji} \, 2^{\nu n} (1 + 2^\nu \|y\|)^{-N} .
\]
\end{theorem}

\begin{proof}
\leanok
By commutativity, $|(\psi_j * \eta_i)(y)| \leq \int |\eta_i(w)| |\psi_j(y - w)| \, dw$, and
$\|\eta_i\|_1 \leq 2 \|m\|_1$ (\Cref{lem:dual-kernel-moduli}).

\emph{Case $i > j$.} $\eta_i$ vanishes outside $\overline B(0, 2LR \, 2^{-i})$, which lies in
$\overline B(0, 2LR \, 2^{-j})$; \Cref{lem:small-support-decay} with $k = j$ and
\Cref{lem:lp-blocks-decay} give the bound $S_{L,N} 2^{jn} (1 + 2^j \|y\|)^{-N}$.

\emph{Case $i \leq j \leq 1$.} The same with $k = i$, since
$|\psi_j(z)| \leq 2^n P_N 2^{in} (1 + 2^i \|z\|)^{-N}$, and $1 \leq 2^{L'} 2^{-(j-i)L'}$.

\emph{Case $i \leq j$, $j \geq 2$.} By \Cref{lem:lp-moment-kernels} and associativity,
$\psi_j * (\mu_{2^{j-2}} * \eta_i) = (\psi_j * \mu_{2^{j-2}}) * \eta_i = 0$ almost everywhere,
so $\psi_j * \eta_i = \psi_j * (\eta_i - \mu_{2^{j-2}} * \eta_i)$ almost everywhere.
\Cref{lem:dual-averaged-differences} with $j' = j - 2$, \Cref{lem:lp-blocks-decay} with the
order $2N$ and \Cref{lem:two-decay-bumps} give the bound
$D_{L,L',N} P_{2N} I_N \, 2^{-(j-i)L'} 2^{in} (1 + 2^i \|y\|)^{-N}$.
\end{proof}

The order $L'$ of the moment kernel of $\Phi$ is free, so the decay $2^{-(j-i)L'}$ is available
for every $L'$, independently of the order $L$ of the local means.

\begin{theorem}[Convolutions under Peetre majorants]
\label{thm:decay-kernels-maximal}
\lean{Rd.enorm_conv_le_maximalFunction_of_decay}
\leanok
\uses{lem:maximal-function-kernels, def:conv, lem:weighted-jensen}
Let $N > n$ and $\nu \geq 0$ be integers, $C \in [0, \infty)$, and let $K : V \to \C$
satisfy $|K(y)| \leq C 2^{\nu n} (1 + 2^\nu \|y\|)^{-N}$ for almost every $y$. For every
$h : V \to \C$ and every $x \in V$,
\[
  |(K * h)(x)| \leq C_{n,N} \, C \, M|h|(x),
\]
with the constant $C_{n,N}$ of \Cref{lem:maximal-function-kernels}.
\end{theorem}

\begin{proof}
\leanok
$|(K * h)(x)| \leq \int |K(y)| |h(x - y)| \, dy$; the majorant holds almost everywhere, so it
bounds the integral from above, and part (3) of \Cref{lem:maximal-function-kernels} applies to
it.
\end{proof}

The bound holds for every function $h$, periodic functions included.

\section{The local means characterisation}
\label{sec:tl-local-means}

This section describes $\TL{s}{p}{q}$ by the local means of \Cref{def:local-means}, in the range
$1 < p < \infty$, $1 < q \leq \infty$, $0 < s < L$. As in the previous section, $g$ is a bump of
radius $R$, $m$ its moment kernel of order $L$, and $k_j$, $\eta_j$ are the local means kernels
and the dual kernels; $M$ is the maximal function of \Cref{def:maximal-function}. For a real
$L$ the weights are $c^L_{ji} = 2^{-(j-i)L}$ for $i \leq j$ and $c^L_{ji} = 1$ for $i > j$.
Both directions of the characterisation are pointwise bounds by the cross kernels of the
previous section, summed in $\ell_q$ by the almost diagonal estimate and integrated by the
Fefferman--Stein inequality. Rychkov~\cite{Rychkov} describes $\TL{s}{p}{q}$ by local means for
all $0 < p < \infty$ and $0 < q \leq \infty$, with the Peetre maximal function in place of $M$.

\begin{definition}[Local means of Triebel--Lizorkin type]
\label{def:tl-local-means}
\lean{Rd.LocalMeansBump.tlLocalMeansNorm, Rd.LocalMeansBump.localMean_ae_eq,
  Rd.LocalMeansBump.continuous_conv_kernel,
  Rd.LocalMeansBump.measurable_dyadicWeight_mul_conv_kernel,
  Rd.LocalMeansBump.aemeasurable_lqNorm_conv_kernel}
\leanok
\uses{def:local-means, def:lq-norm, lem:reproducing-pointwise, lem:conv-op,
  lem:lp-norm-nonnegative}
Let $1 \leq p < \infty$, $s \in \R$, $1 \leq q \leq \infty$ and $L \in \N$. For $f \in L_p$ the
local means $k_j * f$ are continuous functions on $V$, defined at every point, and represent the
local means $k_j * f \in L_p$. The local means quasi-norm of Triebel--Lizorkin type is
\[
  \|f\|^{g,L}_{F^s_{p,q}} = \Bigl\| x \mapsto \bigl\| (2^{js} |(k_j * f)(x)|)_{j \in \N}
  \bigr\|_{\ell_q} \Bigr\|_{L_p} .
\]
\end{definition}

The local means quasi-norm of the Besov spaces (\Cref{def:local-means}) takes the two norms in
the other order.

\begin{lemma}[Almost diagonal sums at one point]
\label{lem:almost-diagonal-pointwise}
\lean{Rd.ite_dyadicWeight_ne_top, Rd.dyadicWeight_mul_maximalFunction}
\leanok
\uses{lem:almost-diagonal, lem:maximal-function-basic}
For every $L \in \R$ and all $i, j \in \N$, $c^L_{ji} < \infty$. For every $s \in \R$, every
$h : V \to [0, \infty]$, every $i \in \N$ and every $x \in V$,
$2^{is} M h(x) = M(2^{is} h)(x)$.
\end{lemma}

With the constant $A_{s,L}$ of \Cref{lem:almost-diagonal}, a family with
$a_j(x) \leq C \sum_i c^L_{ji} M h_i(x)$ for every $j$ therefore satisfies
$\|(2^{js} a_j(x))_j\|_{\ell_q} \leq C A_{s,L} \|(M(2^{is} h_i)(x))_i\|_{\ell_q}$ at every point
$x$.

\begin{proof}
\leanok
The weights are powers of $2$ or $1$. The second claim is part (3) of
\Cref{lem:maximal-function-basic}.
\end{proof}

\begin{lemma}[From pointwise maximal bounds to $L_p(\ell_q)$]
\label{lem:maximal-bounds-mixed-norm}
\lean{Rd.eLpNorm_le_of_le_lqNorm_maximalFunction}
\leanok
\uses{thm:fefferman-stein, lem:lp-norm-nonnegative, lem:lp-norm-homogeneity,
  lem:maximal-function-measurable}
Let $1 < p < \infty$, $1 < q \leq \infty$ and $C \in [0, \infty]$. For every almost everywhere
measurable $F : V \to [0, \infty]$ and every family of almost everywhere measurable
$h_i : V \to [0, \infty]$, $i \in \N$, with $F(x) \leq C \|(M h_i(x))_i\|_{\ell_q}$ for almost
every $x$,
\[
  \|F\|_p \leq C \, C_{n,p,q} \Bigl\| x \mapsto \|(h_i(x))_i\|_{\ell_q} \Bigr\|_p ,
\]
with the constant $C_{n,p,q}$ of \Cref{thm:fefferman-stein}.
\end{lemma}

\begin{proof}
\leanok
Monotonicity and homogeneity of the $L_p$ norm, then \Cref{thm:fefferman-stein}.
\end{proof}

\begin{theorem}[Blocks by local means]
\label{thm:blocks-by-local-means}
\lean{Rd.LocalMeansBump.blockByLocalMeansConst, Rd.LocalMeansBump.blockByLocalMeansConst_ne_top,
  Rd.LocalMeansBump.enorm_conv_conv_lpBlockKernel_dualKernel_le,
  Rd.LocalMeansBump.lpBlock_sum_dual_localMean_ae, Rd.LocalMeansBump.ae_enorm_lpBlock_sum_dual_le,
  Rd.LocalMeansBump.maximalFunction_enorm_localMean,
  Rd.LocalMeansBump.enorm_blockFun_le_tsum_maximalFunction}
\leanok
\uses{thm:cross-kernel-dual, thm:decay-kernels-maximal, lem:maximal-function-kernels,
  lem:reproducing-formula, lem:lp-limits-pointwise, lem:conv-op, lem:almost-diagonal-pointwise,
  def:triebel-lizorkin, def:tl-local-means}
Let $L \in \N$ and $L' \geq 1$, and let
$B_{L,L'} = C_{L,L',n+1} \, C_{n,n+1} < \infty$ with the constants of
\Cref{thm:cross-kernel-dual} and \Cref{lem:maximal-function-kernels}. For every
$h : V \to \C$, all $i, j \in \N$ and every $x \in V$,
\[
  |((\psi_j * \eta_i) * h)(x)| \leq B_{L,L'} \, c^{L'}_{ji} \, M|h|(x) .
\]
If $L \geq 1$ and $1 \leq p < \infty$, then for every $f \in L_p$, almost every $x \in V$ and
every $j \in \N$,
\[
  |(\psi_j * f)(x)| \leq B_{L,L'} \sum_{i \in \N} c^{L'}_{ji} \, M|k_i * f|(x) .
\]
\end{theorem}

\begin{proof}
\leanok
The first bound is \Cref{thm:cross-kernel-dual} with $N = n + 1$ and
\Cref{thm:decay-kernels-maximal}, with the finite weights $c^{L'}_{ji}$ of
\Cref{lem:almost-diagonal-pointwise}. Let
$G_J = \sum_{i \leq J} \eta_i * k_i * f$. By associativity (\Cref{lem:conv-op}),
$\psi_j * G_J = \sum_{i \leq J} (\psi_j * \eta_i) * (k_i * f)$ almost everywhere, so
$|\psi_j * G_J| \leq B_{L,L'} \sum_i c^{L'}_{ji} M|k_i * f|$ almost everywhere. Since
$G_J \to f$ in $L_p$ (\Cref{lem:reproducing-formula}) and $f \mapsto \psi_j * f$ is bounded on
$L_p$, $\psi_j * G_J \to \psi_j * f$ in $L_p$, and \Cref{lem:lp-limits-pointwise} passes the
bound to $f$.
\end{proof}

\begin{theorem}[Local means by blocks]
\label{thm:local-means-by-blocks}
\lean{Rd.LocalMeansBump.localMeansByBlockConst, Rd.LocalMeansBump.localMeansByBlockConst_ne_top,
  Rd.LocalMeansBump.enorm_conv_conv_kernel_Theta_le, Rd.LocalMeansBump.localMean_lpBlock,
  Rd.LocalMeansBump.localMean_sum_lpBlock_ae, Rd.LocalMeansBump.ae_enorm_localMean_sum_lpBlock_le,
  Rd.LocalMeansBump.enorm_localMean_le_tsum_maximalFunction}
\leanok
\uses{thm:cross-kernel-theta, thm:decay-kernels-maximal, lem:maximal-function-kernels,
  lem:lp-blocks, thm:kernel-identities, lem:lp-limits-pointwise, lem:conv-op,
  lem:almost-diagonal-pointwise, def:triebel-lizorkin, def:tl-local-means}
Let $L \in \N$ and $1 \leq p < \infty$, and let
$B^\Theta_L = C^\Theta_{L,n+1} \, C_{n,n+1} < \infty$ with the constants of
\Cref{thm:cross-kernel-theta} and \Cref{lem:maximal-function-kernels}. For every
$h : V \to \C$, all $i, j \in \N$ and every $x \in V$,
\[
  |((k_j * \Theta_i) * h)(x)| \leq B^\Theta_L \, c^L_{ji} \, M|h|(x) .
\]
For every $f \in L_p$, almost every $x \in V$ and every $j \in \N$,
\[
  |(k_j * f)(x)| \leq B^\Theta_L \sum_{i \in \N} c^L_{ji} \, M|\psi_i * f|(x) .
\]
\end{theorem}

\begin{proof}
\leanok
The first bound is \Cref{thm:cross-kernel-theta} with $N = n + 1$ and
\Cref{thm:decay-kernels-maximal}, with the finite weights $c^L_{ji}$ of
\Cref{lem:almost-diagonal-pointwise}. Since
$\Theta_i * \psi_i = \psi_i$ (\Cref{thm:kernel-identities}),
$k_j * \Delta_i f = (k_j * \Theta_i) * \Delta_i f$ in $L_p$, so
$k_j * \sum_{i < J} \Delta_i f = \sum_{i < J} (k_j * \Theta_i) * (\psi_i * f)$ almost
everywhere. The partial sums converge to $f$ in $L_p$ (\Cref{lem:lp-blocks}), so their local
means converge to $k_j * f$, and \Cref{lem:lp-limits-pointwise} passes the bound to $f$.
\end{proof}

The theorem holds for every order $L \in \N$: the cross kernel $k_j * \Theta_i$ uses only
$\int g = 1$.

\begin{lemma}[The almost diagonal step]
\label{lem:tl-local-means-lq}
\lean{Rd.LocalMeansBump.lqNorm_blockFun_le_lqNorm_maximalFunction,
  Rd.LocalMeansBump.lqNorm_localMean_le_lqNorm_maximalFunction}
\leanok
\uses{thm:blocks-by-local-means, thm:local-means-by-blocks, lem:almost-diagonal-pointwise}
Let $1 \leq p < \infty$, $1 \leq q \leq \infty$, $s \in \R$ and $L, L' \in \N$. For every
$f \in L_p$ and almost every $x \in V$,
\[
  \bigl\| (2^{js} |(k_j * f)(x)|)_j \bigr\|_{\ell_q} \leq B^\Theta_L A_{s,L}
  \bigl\| (M(2^{is} |\psi_i * f|)(x))_i \bigr\|_{\ell_q},
\]
and if $L \geq 1$ and $L' \geq 1$, then
\[
  \bigl\| (2^{js} |(\psi_j * f)(x)|)_j \bigr\|_{\ell_q} \leq B_{L,L'} A_{s,L'}
  \bigl\| (M(2^{is} |k_i * f|)(x))_i \bigr\|_{\ell_q} .
\]
\end{lemma}

\begin{proof}
\leanok
\Cref{thm:blocks-by-local-means} and \Cref{thm:local-means-by-blocks} at a point, with
\Cref{lem:almost-diagonal-pointwise}.
\end{proof}

\begin{theorem}[Local means characterisation of Triebel--Lizorkin spaces]
\label{thm:tl-local-means}
\lean{Rd.LocalMeansBump.tlUpperConst, Rd.LocalMeansBump.tlUpperConst_ne_top,
  Rd.LocalMeansBump.tlLowerConst, Rd.LocalMeansBump.tlLowerConst_ne_top,
  Rd.LocalMeansBump.triebelLizorkinNorm_le_tlLocalMeansNorm,
  Rd.LocalMeansBump.tlLocalMeansNorm_le_triebelLizorkinNorm,
  Rd.LocalMeansBump.memTriebelLizorkin_iff_tlLocalMeans}
\leanok
\uses{lem:tl-local-means-lq, lem:maximal-bounds-mixed-norm, lem:almost-diagonal-pointwise,
  thm:blocks-by-local-means, thm:local-means-by-blocks, thm:fefferman-stein,
  def:triebel-lizorkin, def:tl-local-means}
Let $1 < p < \infty$, $1 < q \leq \infty$, $s \in \R$ and $L \in \N$, let
$L_s = \lceil s \rceil_+ + 1$, where $\lceil s \rceil_+ = \max(\lceil s \rceil, 0)$ is the least
natural number at least $s$, and let
\[
  U_{L,p,q,s} = B_{L,L_s} A_{s,L_s} C_{n,p,q}, \qquad
  D_{L,p,q,s} = B^\Theta_L A_{s,L} C_{n,p,q} .
\]
Then $U_{L,p,q,s} < \infty$ for $s > 0$, and $D_{L,p,q,s} < \infty$ for $0 < s < L$. For every
$f \in L_p$,
\[
  \|f\|^{g,L}_{F^s_{p,q}} \leq D_{L,p,q,s} \, \|f\|_{F^s_{p,q}},
\]
and if $L \geq 1$, then $\|f\|_{F^s_{p,q}} \leq U_{L,p,q,s} \, \|f\|^{g,L}_{F^s_{p,q}}$. In
particular, for $0 < s < L$, $f \in \TL{s}{p}{q}$ if and only if
$\|f\|^{g,L}_{F^s_{p,q}} < \infty$.
\end{theorem}

\begin{proof}
\leanok
Both bounds are \Cref{lem:tl-local-means-lq}, with $L' = L_s$ for the second, followed by
\Cref{lem:maximal-bounds-mixed-norm} for the families $2^{is} |\psi_i * f|$ and
$2^{is} |k_i * f|$. The constants $B_{L,L_s}$, $B^\Theta_L$ and $C_{n,p,q}$ are finite, and
since $L_s > s$, \Cref{lem:almost-diagonal} gives $A_{s,L_s} < \infty$ for $s > 0$
and $A_{s,L} < \infty$ for $0 < s < L$. For $0 < s < L$ the order satisfies $L \geq 1$.
\end{proof}

\begin{theorem}[Synthesis in Triebel--Lizorkin spaces]
\label{thm:tl-synthesis-dual}
\lean{Rd.LocalMeansBump.lpBlock_sum_dual_ae,
  Rd.LocalMeansBump.ae_enorm_lpBlock_sum_dualKernel_le_tsum_maximalFunction,
  Rd.LocalMeansBump.enorm_blockFun_le_tsum_maximalFunction_of_tendsto,
  Rd.LocalMeansBump.triebelLizorkinNorm_le_of_tendsto_sum_dual}
\leanok
\uses{thm:blocks-by-local-means, thm:tl-local-means, lem:lp-limits-pointwise, lem:conv-op,
  lem:almost-diagonal-pointwise, lem:maximal-bounds-mixed-norm}
Let $g$ be a bump with the dual kernels $\eta_i$ of order $L$, let $1 < p < \infty$,
$1 < q \leq \infty$ and $s \in \R$. For every family $h_i \in L_p$ and every $F \in L_p$ with
$\sum_{i \leq J} \eta_i * h_i \to F$ in $L_p$,
\[
  \|F\|_{F^s_{p,q}} \leq C \, \bigl\| x \mapsto \|(2^{is} |h_i(x)|)_i\|_{\ell_q} \bigr\|_p
\]
with the constant $C$ of \Cref{thm:tl-local-means}, finite for $s > 0$.
\end{theorem}

\begin{proof}
\leanok
The proof of \Cref{thm:blocks-by-local-means} with $h_i$ in place of $k_i * f$ gives
$|(\psi_j * F)(x)| \leq B_{L,L'} \sum_i c^{L'}_{ji} M|h_i|(x)$ almost everywhere, for
$L' = \lceil s \rceil + 1$; the almost diagonal estimate (\Cref{lem:almost-diagonal-pointwise})
and the Fefferman--Stein inequality (\Cref{lem:maximal-bounds-mixed-norm}) finish as in
\Cref{thm:tl-local-means}.
\end{proof}

\begin{corollary}[Independence of the bump and the order]
\label{cor:tl-local-means-independence}
\lean{Rd.LocalMeansBump.tlLocalMeansNorm_le_tlLocalMeansNorm,
  Rd.LocalMeansBump.tlLocalMeansNorm_lt_top_iff}
\leanok
\uses{thm:tl-local-means}
Let $1 < p < \infty$, $1 < q \leq \infty$ and $s \in \R$, and let $g, g'$ be two bumps with
orders $L, L' \in \N$. If $L \geq 1$, then for every $f \in L_p$,
\[
  \|f\|^{g',L'}_{F^s_{p,q}} \leq D'_{L',p,q,s} \, U_{L,p,q,s} \, \|f\|^{g,L}_{F^s_{p,q}},
\]
with $D'$ the constant of \Cref{thm:tl-local-means} for $g'$. If
$0 < s < \min(L, L')$, then $\|f\|^{g,L}_{F^s_{p,q}} < \infty$ if and only if
$\|f\|^{g',L'}_{F^s_{p,q}} < \infty$.
\end{corollary}

\begin{proof}
\leanok
Both quasi-norms are compared with $\|f\|_{F^s_{p,q}}$ by \Cref{thm:tl-local-means}.
\end{proof}

Every bump and every order $L > s$ describe the same space $\TL{s}{p}{q}$, the space defined by
the Littlewood--Paley blocks.

\section{Independence of the Littlewood--Paley family}
\label{sec:lp-family-independence}

The spaces $\TL{s}{p}{q}$ of this chapter use the standard kernel $\Phi$ of
\Cref{def:lp-kernels}. An admissible kernel (\Cref{def:lp-kernel-admissible}) is a Schwartz
function $\Phi'$ with $\widehat{\Phi'} = 1$ on a ball and $\widehat{\Phi'} = 0$ outside a larger
ball; its blocks are generated by dilation, as in Triebel's standard example of a system in his
class $\Phi(\R^n)$~\cite[Section~2.3.1]{Triebel}, and it defines the quasi-norm
$\|u\|_{F^s_{p,q}, K}$ of \Cref{def:tl-distr}. Triebel proves independence for his whole class
$\Phi(\R^n)$~\cite[Section~2.3.2]{Triebel}. This section proves it for the families generated by
one admissible kernel, with arbitrary radii $0 < a < b$: any two admissible kernels give
equivalent Triebel--Lizorkin quasi-norms, on tempered distributions and on $L_p$, for every real
$s$, $1 < p < \infty$ and $1 < q \leq \infty$, with the explicit constant $C_{K,N,c,s}$ times the
Fefferman--Stein constant (\Cref{lem:lp-kernel-comparison-blocks}), for every window $c \in \N$
with $b \leq 2^c a'$ and $b' \leq 2^c a$; such a window exists
(\Cref{lem:lp-kernel-finite-identity}). The Littlewood--Paley description of the Besov spaces in
\Cref{ch:rd} uses only the standard kernel; the Besov spaces themselves are described by
differences (\Cref{thm:rd-besov}). The proof uses only finite sums: on the Fourier side the
blocks of one kernel are finite sums of the blocks of the other.

\begin{definition}[The seminorm bounds of an admissible kernel]
\label{def:lp-kernel-base-bound}
\lean{Rd.LPKernel.baseBound}
\leanok
\uses{def:lp-kernel-admissible}
Let $K$ be an admissible kernel with $\Phi'$. For all $k, l \in \N$ let $A^{K}_{k,l}$ be the
larger of the seminorms $\sup_x \|x\|^k \|D^l \Phi'(x)\|$ and
$\sup_x \|x\|^k \|D^l (\Phi' - \delta_{1/2} \Phi')(x)\|$.
\end{definition}

\begin{lemma}[The blocks of an admissible kernel]
\label{lem:lp-kernel-blocks}
\lean{Rd.fourier_schwartzDilate_apply, Rd.LPKernel.isUniformSchwartz_base,
  Rd.LPKernel.fourier_psiS_zero, Rd.LPKernel.fourier_psiS_succ,
  Rd.LPKernel.sum_range_fourier_psiS, Rd.LPKernel.fourier_kernel_inv_smul_eq_one,
  Rd.LPKernel.fourier_kernel_inv_smul_eq_zero, Rd.LPKernel.norm_lt_of_fourier_psiS_ne_zero,
  Rd.LPKernel.lt_norm_of_fourier_psiS_succ_ne_zero}
\leanok
\uses{def:lp-kernel-admissible, def:lp-kernel-base-bound, lem:fourier-dilate,
  def:uniform-schwartz}
Let $K$ be an admissible kernel with $\Phi'$, $a$, $b$. The bases $(G'_j)$ are uniformly Schwartz
with bounds $A^K$. For every $\xi \in V$, $\widehat{\psi'_0}(\xi) = \widehat{\Phi'}(\xi)$,
\[
  \widehat{\psi'_{j+1}}(\xi) = \widehat{\Phi'}(2^{-j-1} \xi) - \widehat{\Phi'}(2^{-j} \xi)
  \quad \text{and} \quad
  \sum_{j \leq J} \widehat{\psi'_j}(\xi) = \widehat{\Phi'}(2^{-J} \xi)
\]
for all $j, J \in \N$. If $\widehat{\psi'_j}(\xi) \neq 0$, then $\|\xi\| < 2^j b$, and if
$\widehat{\psi'_{j+1}}(\xi) \neq 0$, then $2^j a < \|\xi\|$.
\end{lemma}

\begin{proof}
\leanok
$\widehat{G_c}(\xi) = \widehat G(c^{-1} \xi)$ for $c > 0$ (\Cref{lem:fourier-dilate}), and the
bases take only two values. The sum telescopes. If $\|\xi\| \geq 2^j b$, both terms of
$\widehat{\psi'_j}(\xi)$ vanish; if $\|\xi\| \leq 2^j a$, both terms of
$\widehat{\psi'_{j+1}}(\xi)$ are $1$.
\end{proof}

\begin{lemma}[The finite identity]
\label{lem:lp-kernel-finite-identity}
\lean{Rd.LPKernel.exists_window, Rd.LPKernel.sum_Icc_fourier_psiS_eq_one,
  Rd.LPKernel.fourier_psiS_mul_sum_Icc, Rd.LPKernel.distribBlock_eq_sum,
  Rd.LPKernel.distribBlock_eq_sum_convLp}
\leanok
\uses{lem:lp-kernel-blocks, lem:conv-fourier-multiplier}
Let $K$, $K'$ be admissible kernels with $a, b$ and $a', b'$. There is $c \in \N$ with
$b \leq 2^c a'$ and $b' \leq 2^c a$. For every such $c$, every $j \in \N$ and every $\xi \in V$
with $\widehat{\psi_j}(\xi) \neq 0$,
\[
  \sum_{k = j - c}^{j + c} \widehat{\psi'_k}(\xi) = 1 ,
\]
so $\widehat{\psi_j} \sum_{k = j - c}^{j + c} \widehat{\psi'_k} = \widehat{\psi_j}$, with the
blocks $\psi_j$ of $K$ and $\psi'_k$ of $K'$. For every tempered distribution $u$ and every $j$,
\[
  \Delta_j u = \sum_{k = j - c}^{j + c} (\widehat{\psi_j})(D) \, \Delta'_k u ,
\]
and if $1 \leq p < \infty$ and $\Delta'_k u = g_k \in L_p(V)$ for every $k$, then $\Delta_j u$ is
the tempered distribution of $\sum_{k = j - c}^{j + c} \psi_j * g_k \in L_p(V)$.
\end{lemma}

\begin{proof}
\leanok
Take $c$ with $2^c \geq \max(b/a', b'/a)$. By \Cref{lem:lp-kernel-blocks} the sum is
$\widehat{\Phi'}(2^{-j-c} \xi) - \widehat{\Phi'}(2^{-(j-c-1)} \xi)$, the second term absent for
$j \leq c$. Since $\|\xi\| < 2^j b \leq 2^{j+c} a'$, the first term is $1$; for $j > c$,
$\|\xi\| > 2^{j-1} a \geq 2^{j-c-1} b'$, so the second term is $0$. Both sides of the identity for
$u$ are Fourier multipliers of $u$, with the symbols $\widehat{\psi_j}$ and
$\sum_k \widehat{\psi'_k} \, \widehat{\psi_j}$ (Mathlib's
\texttt{TemperedDistribution.\allowbreak fourierMultiplierCLM\_\allowbreak
fourierMultiplierCLM\_\allowbreak apply} and \texttt{fourierMultiplierCLM\_\allowbreak sum}).
The last claim follows with \Cref{lem:conv-fourier-multiplier}.
\end{proof}

\begin{lemma}[The pointwise bound]
\label{lem:lp-kernel-pointwise}
\lean{Rd.LPKernel.dyadicWeight_mul_enorm_sum_conv_le, Rd.LPKernel.comparisonConst,
  Rd.LPKernel.comparisonConst_ne_top}
\leanok
\uses{lem:lp-kernel-blocks, thm:uniform-family-maximal, lem:window-sums,
  lem:almost-diagonal-pointwise}
Let $K$ be an admissible kernel, $N > n$, $c, j \in \N$ and $s$ real. For every family
$g_k \colon V \to \C$ and every $x \in V$,
\[
  2^{js} \Bigl| \sum_{k = j - c}^{j + c} (\psi_j * g_k)(x) \Bigr|
  \leq 2^{c|s|} \, C_{n,N} \, C_{A^K,0,N} \sum_{k = j - c}^{j + c} M\bigl(2^{ks} |g_k|\bigr)(x) .
\]
The comparison constant $C_{K,N,c,s} = 2^{c|s|} \, C_{n,N} \, C_{A^K,0,N} \, (2c + 1)$ is
finite.
\end{lemma}

\begin{proof}
\leanok
By \Cref{thm:uniform-family-maximal} and \Cref{lem:lp-kernel-blocks},
\[
  |(\psi_j * g_k)(x)| \leq C_{n,N} \, C_{A^K,0,N} \, M|g_k|(x) .
\]
In the window
$2^{js} \leq 2^{c|s|} 2^{ks}$ (\Cref{lem:window-sums}), and $2^{ks} M|g_k| = M(2^{ks} |g_k|)$
(\Cref{lem:almost-diagonal-pointwise}).
\end{proof}

\begin{lemma}[The comparison of two admissible kernels]
\label{lem:lp-kernel-comparison-blocks}
\lean{Rd.LPKernel.triebelLizorkinNormDistr_le_of_blocks,
  Rd.LPKernel.triebelLizorkinNormDistr_le_mul}
\leanok
\uses{lem:lp-kernel-finite-identity, lem:tl-distr-eq, lem:lp-kernel-pointwise,
  lem:window-sums, lem:maximal-bounds-mixed-norm, thm:fefferman-stein, lem:conv-fourier-multiplier}
Let $K$, $K'$ be admissible kernels with $a, b$ and $a', b'$, $1 < p < \infty$,
$1 < q \leq \infty$, $s$ real, $N > n$ and $c \in \N$ with $b \leq 2^c a'$ and $b' \leq 2^c a$, and
let $C^{\mathrm{FS}}_{p,q}$ be the Fefferman--Stein constant of \Cref{thm:fefferman-stein}.
\begin{enumerate}
  \item For every tempered distribution $u$ and every family $(g_k)$ in $L_p(V)$ with
    $\Delta'_k u = g_k$ for every $k$,
    \[
      \|u\|_{F^s_{p,q}, K} \leq C_{K,N,c,s} \, C^{\mathrm{FS}}_{p,q}
      \Bigl\| x \mapsto \bigl\| (2^{ks} |g_k(x)|)_k \bigr\|_{\ell_q} \Bigr\|_{L_p} .
    \]
  \item For every tempered distribution $u$,
    \[
      \|u\|_{F^s_{p,q}, K} \leq \max\bigl(1, C_{K,N,c,s} \, C^{\mathrm{FS}}_{p,q}\bigr) \,
      \|u\|_{F^s_{p,q}, K'} .
    \]
\end{enumerate}
\end{lemma}

\begin{proof}
\leanok
(1) By \Cref{lem:lp-kernel-finite-identity} the blocks of $u$ for $K$ are the $L_p$ functions
$h_j = \sum_{k = j - c}^{j + c} \psi_j * g_k$. At almost every $x$, \Cref{lem:lp-kernel-pointwise}
and \Cref{lem:window-sums} give
\[
  \bigl\| (2^{js} |h_j(x)|)_j \bigr\|_{\ell_q}
  \leq C_{K,N,c,s} \, \bigl\| \bigl(M(2^{ks} |g_k|)(x)\bigr)_k \bigr\|_{\ell_q} ,
\]
and \Cref{lem:maximal-bounds-mixed-norm} gives the bound. (2) follows from (1) by
\Cref{lem:tl-distr-eq}, also when the blocks of $u$ for $K'$ are no $L_p$ functions.
\end{proof}

\begin{theorem}[Independence of the Littlewood--Paley family]
\label{thm:lp-family-independence}
\lean{Rd.LPKernel.exists_triebelLizorkinNormDistr_le, Rd.LPKernel.exists_triebelLizorkinNorm_le}
\leanok
\uses{lem:lp-kernel-comparison-blocks, lem:lp-kernel-finite-identity, lem:lp-kernel-pointwise,
  lem:lp-kernel-tl-norm-basic, thm:fefferman-stein}
Let $K$, $K'$ be admissible kernels, $1 < p < \infty$, $1 < q \leq \infty$ and $s$ real. There is
$C < \infty$ such that for every tempered distribution $u$ on $V$ and every $f \in L_p(V)$,
\[
  \|u\|_{F^s_{p,q}, K} \leq C \, \|u\|_{F^s_{p,q}, K'}
  \quad \text{and} \quad
  \|f\|_{F^s_{p,q}, K} \leq C \, \|f\|_{F^s_{p,q}, K'} .
\]
\end{theorem}

\begin{proof}
\leanok
Take a window $c$ from \Cref{lem:lp-kernel-finite-identity}, $N = n + 1$ and
$C = \max(1, C_{K,N,c,s} C^{\mathrm{FS}}_{p,q})$, which is finite (\Cref{lem:lp-kernel-pointwise},
\Cref{thm:fefferman-stein}). Part (2) of \Cref{lem:lp-kernel-comparison-blocks} gives the first
bound, and for $f \in L_p(V)$ the two quasi-norms of $f$ agree
(\Cref{lem:lp-kernel-tl-norm-basic}).
\end{proof}

\begin{corollary}[Every admissible kernel gives the same spaces]
\label{cor:lp-family-same-spaces}
\lean{Rd.LPKernel.memTriebelLizorkinDistr_iff, Rd.LPKernel.memTriebelLizorkin_iff}
\leanok
\uses{thm:lp-family-independence, lem:lp-kernel-tl-norm-basic, thm:tl-distr-lp}
Let $K$ be an admissible kernel, $1 < p < \infty$, $1 < q \leq \infty$ and $s$ real. A tempered
distribution $u$ has $\|u\|_{F^s_{p,q}, K} < \infty$ if and only if $u \in \TL{s}{p}{q}$, and
$f \in L_p(V)$ has $\|f\|_{F^s_{p,q}, K} < \infty$ if and only if $f \in \TL{s}{p}{q}$.
\end{corollary}

\begin{proof}
\leanok
\Cref{thm:lp-family-independence} for $K$ and the standard kernel, in both orders, and
\Cref{thm:tl-distr-lp}.
\end{proof}

\section{Local means on the torus}
\label{sec:torus-tl-local-means}

This section transfers the local means characterisation to the torus $\T^d$. The kernels $k_j$,
$\eta_j$, $\psi_j$ and $\Theta_j$ on $\R^d$ act on $L_p(\T^d)$ by convolution
(\Cref{def:torus-econv}), and the operator identities between them hold on the torus as
identities of Fourier multipliers (\Cref{lem:torus-econv-identities}). For the pointwise bounds,
a function $h$ on $\T^d$ is lifted to the $\Z^d$-periodic function $h \circ \pi$ on $\R^d$, with
the quotient map $\pi$ of \Cref{lem:torus-proj}: since $(K * h)(\pi x) = (K * (h \circ \pi))(x)$,
the bounds of the cross kernels on $\R^d$ apply to lifts, and the Fefferman--Stein inequality
transfers by localisation. Throughout, $\#d$ is the dimension of the torus, $|\cdot|$ the
Lebesgue measure on $\R^d$ and $M$ the maximal function of \Cref{def:maximal-function} on $\R^d$.
The weights $c^L_{ji}$ are those of the previous section (\Cref{lem:almost-diagonal-pointwise})
and the constants $A_{s,L}$ those of \Cref{lem:almost-diagonal}.

\begin{definition}[Period boxes]
\label{def:period-box}
\lean{Torus.periodBox, Torus.tripleBox, Torus.measurableSet_periodBox,
  Torus.measurableSet_tripleBox, Torus.mem_periodBox_intCast_iff, Torus.mem_periodBox_ceil,
  Torus.pairwise_disjoint_periodBox, Torus.tripleBox_eq_biUnion}
\leanok
For $a \in \R^d$ the period box is
$P_a = \{x \in \R^d : a_i < x_i \leq a_i + 1 \text{ for every } i\}$, and
$T = \{y \in \R^d : -1 < y_i \leq 2 \text{ for every } i\}$ is the box of side $3$. Both are
measurable. For every $m \in \Z^d$ and every $x \in \R^d$, $x \in P_m$ if and only if
$\lceil x_i \rceil = m_i + 1$ for every $i$; in particular $x \in P_m$ for
$m_i = \lceil x_i \rceil - 1$, the boxes $P_m$, $m \in \Z^d$, are pairwise disjoint, and $T$ is
the union of the $P_m$ with $m \in \{-1, 0, 1\}^d$.
\end{definition}

\begin{lemma}[Period boxes are fundamental domains]
\label{lem:period-box-measure}
\lean{Torus.lintegral_periodBox, Torus.measurePreserving_proj_periodBox, Torus.volume_periodBox,
  Torus.eLpNorm_comp_proj_periodBox, Torus.quasiMeasurePreserving_proj, Torus.ae_comp_proj}
\leanok
\uses{def:period-box, lem:torus-proj}
For every $a \in \R^d$ and every $u : \T^d \to [0, \infty]$,
\[
  \int_{P_a} u(\pi x) \, dx = \int_{\T^d} u ,
\]
so $\pi$ maps the Lebesgue measure on $P_a$ to the Haar measure of $\T^d$, $|P_a| = 1$, and
$\|F \circ \pi\|_{L_p(P_a)} = \|F\|_{L_p(\T^d)}$ for every $p \in [0, \infty]$ and every almost
everywhere measurable $F$ on $\T^d$. For every null set $N \subseteq \T^d$, $\pi^{-1}(N)$ is a
null set; hence for every property $P$ that holds at almost every point of $\T^d$, $P(\pi x)$
holds for almost every $x \in \R^d$.
\end{lemma}

\begin{proof}
\leanok
The integral identity is Mathlib's \texttt{UnitAddTorus.\allowbreak lintegral\_\allowbreak
preimage} on $\R^d$ with the product measure, carried to the Euclidean space by the
volume-preserving identification of the two; for indicator functions it is the statement on
measures. The boxes $P_m$, $m \in \Z^d$, cover $\R^d$ (\Cref{def:period-box}), and
$|\pi^{-1}(N) \cap P_m| = |N| = 0$ for each of the countably many $m$.
\end{proof}

\begin{lemma}[Counting boxes]
\label{lem:period-box-counting}
\lean{Torus.ball_subset_iUnion_periodBox, Torus.card_Icc_ceil_le,
  Torus.setLIntegral_ball_comp_proj_le, Torus.locallyIntegrable_comp_proj,
  Torus.ball_subset_tripleBox, Torus.setLIntegral_tripleBox}
\leanok
\uses{lem:period-box-measure, def:period-box}
For every $u : \T^d \to [0, \infty]$, every $x \in \R^d$ and every $r > 0$,
\[
  \int_{B(x, r)} u(\pi y) \, dy \leq (2r + 2)^{\#d} \int_{\T^d} u, \qquad
  \int_T u(\pi y) \, dy = 3^{\#d} \int_{\T^d} u .
\]
For every $h \in L_1(\T^d)$, the lift $h \circ \pi$ is locally integrable on $\R^d$. For every
$x \in P_0$ and every $r \leq 1$, $B(x, r) \subseteq T$.
\end{lemma}

\begin{proof}
\leanok
Every point $y$ of the ball lies in $P_m$ with $m_i = \lceil y_i \rceil - 1$, and
$\lceil x_i - r \rceil - 1 \leq m_i \leq \lceil x_i + r \rceil - 1$ since $|y_i - x_i| < r$. These
are at most $\lceil x_i + r \rceil - \lceil x_i - r \rceil + 1 \leq 2r + 2$ integers in each
coordinate, and each box contributes $\int_{\T^d} u$ (\Cref{lem:period-box-measure}). The box $T$
is the disjoint union of $3^{\#d}$ period boxes. Local integrability follows from the first
bound for $u = |h|$ and $r = 1$. For $x \in P_0$ and $y \in B(x, r)$ with $r \leq 1$,
$0 < x_i \leq 1$ and $|y_i - x_i| < 1$ give $-1 < y_i < 2$.
\end{proof}
For $r \leq 0$ the ball is empty, and the Lean statement of the first bound covers every real $r$.

\begin{lemma}[$L_p$ norms on the box of side $3$]
\label{lem:triple-box-norm}
\lean{Torus.eLpNorm_indicator_tripleBox_comp_proj, Torus.eLpNorm_indicator_tripleBox,
  Torus.memLp_indicator_tripleBox_comp_proj}
\leanok
\uses{lem:period-box-counting, def:period-box}
For every $0 < p < \infty$ and every almost everywhere strongly measurable $F$ on $\T^d$ with
values in a space with an extended norm (for example $\C$ or $[0, \infty]$),
\[
  \|\mathbf 1_T (F \circ \pi)\|_{L_p(\R^d)} = 3^{\#d/p} \|F\|_{L_p(\T^d)} .
\]
For every $0 < p < \infty$ and every $f \in L_p(\T^d)$, $\mathbf 1_T (f \circ \pi) \in L_p(\R^d)$.
\end{lemma}

\begin{proof}
\leanok
Both sides are $p$-th roots of integrals of $|F|^p$, and
$\int_T u(\pi y) \, dy = 3^{\#d} \int_{\T^d} u$ (\Cref{lem:period-box-counting}) for
$u = |F|^p$.
\end{proof}

\begin{lemma}[Convolutions and differences of lifts]
\label{lem:torus-lift-conv}
\lean{Torus.econv_proj, Torus.higherDiff_comp_proj}
\leanok
\uses{def:torus-econv, lem:torus-proj, def:iter-diff}
For every $\Phi : \R^d \to \C$, every $h : \T^d \to \C$, every $L \in \N$ and all
$x, z \in \R^d$,
\[
  (\Phi * h)(\pi x) = (\Phi * (h \circ \pi))(x), \qquad
  \Delta^L_{\pi z} h(\pi x) = \Delta^L_z (h \circ \pi)(x),
\]
with the convolution on $\R^d$ on the right of the first identity.
\end{lemma}

\begin{proof}
\leanok
$\pi$ is additive, so $h(\pi x - \pi y) = (h \circ \pi)(x - y)$ and
$h(\pi x + l \pi z) = (h \circ \pi)(x + l z)$.
\end{proof}

\begin{lemma}[Localisation of the maximal function]
\label{lem:torus-maximal-localisation}
\lean{Torus.ofReal_two_mul_add_two_pow_le_volume_ball_mul,
  Torus.maximalFunction_comp_quotientProj_le, Torus.lqNorm_maximalFunction_comp_quotientProj_le,
  Torus.maximalFunction_enorm_comp_quotientProj_congr_ae}
\leanok
\uses{def:maximal-function, lem:period-box-counting, def:period-box, lem:lq-basic,
  lem:period-box-measure, lem:maximal-function-basic}
For every $u : \T^d \to [0, \infty]$ and every $x \in P_0$,
\[
  M(u \circ \pi)(x) \leq M(\mathbf 1_T (u \circ \pi))(x) + 4^{\#d} |B(0,1)|^{-1}
  \int_{\T^d} u .
\]
For every $1 \leq q \leq \infty$, every family $u_j : \T^d \to [0, \infty]$, $j \in \N$, and every
$x \in P_0$,
\begin{multline*}
  \|(M(u_j \circ \pi)(x))_j\|_{\ell_q} \leq \|(M(\mathbf 1_T (u_j \circ \pi))(x))_j\|_{\ell_q} \\
  + 4^{\#d} |B(0,1)|^{-1} \Bigl\| \Bigl(\int_{\T^d} u_j\Bigr)_j \Bigr\|_{\ell_q} .
\end{multline*}
For all $v, w : \T^d \to \C$ with $v = w$ almost everywhere,
$M(|v| \circ \pi) = M(|w| \circ \pi)$.
\end{lemma}

\begin{proof}
\leanok
For $r \leq 1$ the ball $B(x, r)$ lies in $T$ (\Cref{lem:period-box-counting}), so its mean of
$u \circ \pi$ is a mean of $\mathbf 1_T (u \circ \pi)$. For $r > 1$,
\Cref{lem:period-box-counting} and $|B(0, r)| = r^{\#d} |B(0, 1)|$ with $2r + 2 \leq 4r$ bound
the mean by $4^{\#d} |B(0,1)|^{-1} \int_{\T^d} u$. The second bound follows by the triangle
inequality and homogeneity in $\ell_q$. If $v = w$ almost everywhere, the lifts agree almost
everywhere (\Cref{lem:period-box-measure}), and the maximal function does not see the difference
(\Cref{lem:maximal-function-basic}).
\end{proof}

\begin{theorem}[The Fefferman--Stein inequality for periodic lifts]
\label{thm:fefferman-stein-periodic}
\lean{Torus.periodicFSConst, Torus.periodicFSConst_ne_top, Torus.feffermanStein_periodic}
\leanok
\uses{thm:fefferman-stein, lem:torus-maximal-localisation, lem:triple-box-norm,
  lem:period-box-measure, lem:minkowski-integral-lq, lem:maximal-function-measurable,
  lem:lp-norm-nonnegative, lem:fefferman-stein-split}
Let $1 < p < \infty$ and $1 < q \leq \infty$, and let
$C^{\T}_{p,q} = 3^{\#d/p} C_{\#d,p,q} + 4^{\#d} |B(0,1)|^{-1}$ with the constant $C_{\#d,p,q}$
of \Cref{thm:fefferman-stein} on $\R^d$. Then $C^{\T}_{p,q} < \infty$, and for every family of
almost everywhere measurable $u_j : \T^d \to [0, \infty]$, $j \in \N$,
\[
  \Bigl\| x \mapsto \|(M(u_j \circ \pi)(x))_j\|_{\ell_q} \Bigr\|_{L_p(P_0)}
  \leq C^{\T}_{p,q} \Bigl\| t \mapsto \|(u_j(t))_j\|_{\ell_q} \Bigr\|_{L_p(\T^d)} .
\]
\end{theorem}

\begin{proof}
\leanok
Write $G(t) = \|(u_j(t))_j\|_{\ell_q}$. By \Cref{lem:torus-maximal-localisation} the left side is
at most the sum of two terms. The first is at most its norm in $L_p(\R^d)$, which
\Cref{thm:fefferman-stein} bounds by $C_{\#d,p,q} \|\mathbf 1_T (G \circ \pi)\|_{L_p(\R^d)}
= C_{\#d,p,q} \, 3^{\#d/p} \|G\|_{L_p(\T^d)}$ (\Cref{lem:triple-box-norm}). The second is a
constant on a set of measure $1$, and Minkowski's integral inequality
(\Cref{lem:minkowski-integral-lq}) on the probability space $\T^d$
gives $\|(\int u_j)_j\|_{\ell_q} \leq \|G\|_{L_p(\T^d)}$.
\end{proof}

\begin{lemma}[From maximal bounds on the lift to $L_p(\ell_q)$]
\label{lem:torus-maximal-bounds-mixed-norm}
\lean{Torus.eLpNorm_le_of_le_lqNorm_maximalFunction_quotientProj}
\leanok
\uses{thm:fefferman-stein-periodic, lem:period-box-measure, lem:lp-norm-homogeneity,
  lem:maximal-function-measurable}
Let $1 < p < \infty$, $1 < q \leq \infty$ and $C \in [0, \infty]$. For every almost everywhere
measurable $F : \T^d \to [0, \infty]$ and every family of almost everywhere measurable
$h_i : \T^d \to [0, \infty]$, $i \in \N$, with
$F(\pi x) \leq C \|(M(h_i \circ \pi)(x))_i\|_{\ell_q}$ for almost every $x \in P_0$,
\[
  \|F\|_{L_p(\T^d)} \leq C \, C^{\T}_{p,q} \Bigl\| t \mapsto \|(h_i(t))_i\|_{\ell_q}
  \Bigr\|_{L_p(\T^d)} .
\]
\end{lemma}

\begin{proof}
\leanok
$\|F\|_{L_p(\T^d)} = \|F \circ \pi\|_{L_p(P_0)}$ (\Cref{lem:period-box-measure}), then
monotonicity and homogeneity of the $L_p$ norm and \Cref{thm:fefferman-stein-periodic}.
\end{proof}

\begin{definition}[Local means of Triebel--Lizorkin type on the torus]
\label{def:torus-tl-local-means}
\lean{Torus.tlLocalMeansNorm, Torus.aemeasurable_dyadicWeight_mul_econv_kernel,
  Torus.aemeasurable_lqNorm_econv_kernel}
\leanok
\uses{def:tl-local-means, def:torus-econv, def:local-means, lem:lp-norm-nonnegative}
Let $1 \leq p < \infty$, $s \in \R$, $1 \leq q \leq \infty$ and $L \in \N$, and let $k_j$ be the
local means kernels on $\R^d$ of a bump $g$ and the order $L$. For $f \in L_p(\T^d)$ the local
means $(k_j * f)(t) = \int_{\R^d} k_j(y) f(t - \pi y) \, dy$ are almost everywhere measurable,
and the local means quasi-norm of Triebel--Lizorkin type on the torus is
\[
  \|f\|^{g,L}_{F^s_{p,q}(\T^d)} = \Bigl\| t \mapsto \bigl\| (2^{js} |(k_j * f)(t)|)_{j \in \N}
  \bigr\|_{\ell_q} \Bigr\|_{L_p(\T^d)} .
\]
\end{definition}

\begin{lemma}[The reproducing formula on the torus]
\label{lem:torus-reproducing-formula}
\lean{Torus.sum_econvLp_dual_kernel, Torus.tendsto_sum_econvLp_dual_kernel,
  Torus.econvLp_kernel_lpBlock}
\leanok
\uses{lem:torus-econv-identities, thm:torus-approx-identity, def:local-means,
  lem:reproducing-formula, thm:kernel-identities, def:torus-lp-blocks}
Let $1 \leq p < \infty$, $L \in \N$ and let $m$, $k_j$, $\eta_j$ be the moment kernel, the local
means kernels and the dual kernels on $\R^d$ of a bump $g$ and the order $L$. For every
$f \in L_p(\T^d)$ and every $J \in \N$,
\[
  \sum_{j \leq J} \eta_j * k_j * f = m_{2^J} * m_{2^J} * f ,
\]
and if $L \geq 1$, then $\sum_{j \leq J} \eta_j * k_j * f \to f$ in $L_p(\T^d)$ as $J \to \infty$.
For all $i, j \in \N$, $k_j * \Delta_i f = (k_j * \Theta_i) * \Delta_i f$.
\end{lemma}

\begin{proof}
\leanok
As on $V$ (\Cref{lem:reproducing-formula}): the operators commute and are additive in the
kernel (\Cref{lem:torus-econv-identities}), so the partial sums telescope by
$a^2 - b^2 = (a - b)(a + b)$. Since $\int m = 1$ for $L \geq 1$, $m_{2^J} * f \to f$
(\Cref{thm:torus-approx-identity}), and the operators $m_{2^J} *$ are bounded by $\|m\|_1$. The
last identity is $\Theta_i * \psi_i = \psi_i$ (\Cref{thm:kernel-identities}) and associativity.
\end{proof}

\begin{theorem}[Blocks by local means on the torus]
\label{thm:torus-blocks-by-local-means}
\lean{Torus.lpBlock_sum_econvLp_dual_ae,
  Torus.maximalFunction_periodicEconvLp_meansKernel_quotientProj,
  Torus.ae_enorm_lpBlock_sum_proj_le, Torus.ae_enorm_comp_proj_le_of_tendsto_Lp,
  Torus.enorm_blockFun_proj_le_tsum_maximal}
\leanok
\uses{thm:blocks-by-local-means, lem:torus-reproducing-formula, lem:torus-lift-conv,
  lem:period-box-measure, lem:torus-econv-identities, lem:lp-limits-pointwise,
  lem:torus-maximal-localisation, def:torus-triebel-lizorkin, def:torus-tl-local-means,
  lem:torus-econv-young}
Let $L, L' \in \N$ with $L, L' \geq 1$ and $1 \leq p < \infty$, and let $B_{L,L'}$ be the
constant of \Cref{thm:blocks-by-local-means} on $\R^d$. For every $f \in L_p(\T^d)$, almost every
$x \in P_0$ and every $j \in \N$,
\[
  |(\psi_j * f)(\pi x)| \leq B_{L,L'} \sum_{i \in \N} c^{L'}_{ji} \,
  M\bigl(|k_i * f| \circ \pi\bigr)(x) .
\]
\end{theorem}

\begin{proof}
\leanok
Let $G_J = \sum_{i \leq J} \eta_i * k_i * f$. By \Cref{lem:torus-econv-identities},
$\psi_j * G_J = \sum_{i \leq J} (\psi_j * \eta_i) * (k_i * f)$ almost everywhere on $\T^d$, hence
at $\pi x$ for almost every $x \in \R^d$ (\Cref{lem:period-box-measure}). There every term is the
convolution on $\R^d$ of $\psi_j * \eta_i$ with the lift of $k_i * f$ (\Cref{lem:torus-lift-conv}),
and the first bound of \Cref{thm:blocks-by-local-means}, which holds for every function, bounds
it by $B_{L,L'} c^{L'}_{ji} M(|k_i * f| \circ \pi)(x)$; the maximal function does not see the
choice of representative (\Cref{lem:torus-maximal-localisation}). Composition with $\pi$ is an
isometry from $L_p(\T^d)$ to $L_p(P_0)$, $\psi_j * G_J \to \psi_j * f$ by
\Cref{lem:torus-reproducing-formula}, and \Cref{lem:lp-limits-pointwise} on $P_0$ passes the
bound to $f$.
\end{proof}

\begin{theorem}[Local means by blocks on the torus]
\label{thm:torus-local-means-by-blocks}
\lean{Torus.econvLp_kernel_sum_lpBlock_ae, Torus.maximalFunction_periodicLpBlock_quotientProj,
  Torus.ae_enorm_econvLp_kernel_sum_proj_le, Torus.enorm_econv_kernel_proj_le_tsum_maximal}
\leanok
\uses{thm:local-means-by-blocks, lem:torus-reproducing-formula, lem:torus-lift-conv,
  lem:period-box-measure, lem:torus-lp-blocks, lem:lp-limits-pointwise,
  lem:torus-maximal-localisation, def:torus-triebel-lizorkin, def:torus-tl-local-means,
  lem:torus-econv-young, thm:torus-blocks-by-local-means}
Let $L \in \N$ and $1 \leq p < \infty$, and let $B^\Theta_L$ be the constant of
\Cref{thm:local-means-by-blocks} on $\R^d$. For every $f \in L_p(\T^d)$, almost every
$x \in P_0$ and every $j \in \N$,
\[
  |(k_j * f)(\pi x)| \leq B^\Theta_L \sum_{i \in \N} c^L_{ji} \,
  M\bigl(|\psi_i * f| \circ \pi\bigr)(x) .
\]
\end{theorem}

\begin{proof}
\leanok
By \Cref{lem:torus-reproducing-formula},
$k_j * \sum_{i < J} \Delta_i f = \sum_{i < J} (k_j * \Theta_i) * (\psi_i * f)$ almost everywhere,
and the first bound of \Cref{thm:local-means-by-blocks} applies to the lifts at $\pi x$ as in
\Cref{thm:torus-blocks-by-local-means}. The partial sums converge to $f$ in $L_p(\T^d)$
(\Cref{lem:torus-lp-blocks}), and \Cref{lem:lp-limits-pointwise} on $P_0$ passes the bound to
$f$.
\end{proof}

\begin{lemma}[The almost diagonal step on the torus]
\label{lem:torus-tl-local-means-lq}
\lean{Torus.lqNorm_blockFun_proj_le_lqNorm_maximal,
  Torus.lqNorm_econv_kernel_proj_le_lqNorm_maximal}
\leanok
\uses{thm:torus-blocks-by-local-means, thm:torus-local-means-by-blocks,
  lem:almost-diagonal-pointwise}
Let $1 \leq p < \infty$, $1 \leq q \leq \infty$, $s \in \R$ and $L, L' \in \N$. For every
$f \in L_p(\T^d)$ and almost every $x \in P_0$,
\[
  \bigl\| (2^{js} |(k_j * f)(\pi x)|)_j \bigr\|_{\ell_q} \leq B^\Theta_L A_{s,L}
  \bigl\| (M(2^{is} |\psi_i * f| \circ \pi)(x))_i \bigr\|_{\ell_q},
\]
and if $L \geq 1$ and $L' \geq 1$, then
\[
  \bigl\| (2^{js} |(\psi_j * f)(\pi x)|)_j \bigr\|_{\ell_q} \leq B_{L,L'} A_{s,L'}
  \bigl\| (M(2^{is} |k_i * f| \circ \pi)(x))_i \bigr\|_{\ell_q} .
\]
\end{lemma}

\begin{proof}
\leanok
\Cref{thm:torus-blocks-by-local-means} and \Cref{thm:torus-local-means-by-blocks} at a point,
with \Cref{lem:almost-diagonal-pointwise}.
\end{proof}

\begin{theorem}[Local means characterisation of Triebel--Lizorkin spaces on the torus]
\label{thm:torus-tl-local-means}
\lean{Torus.tlUpperConst, Torus.tlUpperConst_ne_top, Torus.tlLowerConst,
  Torus.tlLowerConst_ne_top, Torus.triebelLizorkinNorm_le_tlLocalMeansNorm,
  Torus.tlLocalMeansNorm_le_triebelLizorkinNorm, Torus.memTriebelLizorkin_iff_tlLocalMeans}
\leanok
\uses{lem:torus-tl-local-means-lq, lem:torus-maximal-bounds-mixed-norm,
  lem:almost-diagonal-pointwise, thm:fefferman-stein-periodic, def:torus-triebel-lizorkin,
  def:torus-tl-local-means, thm:tl-local-means}
Let $1 < p < \infty$, $1 < q \leq \infty$, $s \in \R$ and $L \in \N$, let
$L_s = \lceil s \rceil_+ + 1$, and let
\[
  U^{\T}_{L,p,q,s} = B_{L,L_s} A_{s,L_s} C^{\T}_{p,q}, \qquad
  D^{\T}_{L,p,q,s} = B^\Theta_L A_{s,L} C^{\T}_{p,q} .
\]
Then $U^{\T}_{L,p,q,s} < \infty$ for $s > 0$, and $D^{\T}_{L,p,q,s} < \infty$ for $0 < s < L$.
For every $f \in L_p(\T^d)$,
\[
  \|f\|^{g,L}_{F^s_{p,q}(\T^d)} \leq D^{\T}_{L,p,q,s} \, \|f\|_{F^s_{p,q}(\T^d)},
\]
and if $L \geq 1$, then
$\|f\|_{F^s_{p,q}(\T^d)} \leq U^{\T}_{L,p,q,s} \, \|f\|^{g,L}_{F^s_{p,q}(\T^d)}$. In particular,
for $0 < s < L$, $f \in \TL{s}{p}{q}(\T^d)$ if and only if
$\|f\|^{g,L}_{F^s_{p,q}(\T^d)} < \infty$.
\end{theorem}

\begin{proof}
\leanok
Both bounds are \Cref{lem:torus-tl-local-means-lq}, with $L' = L_s$ for the second, followed by
\Cref{lem:torus-maximal-bounds-mixed-norm} for the families $2^{is} |\psi_i * f|$ and
$2^{is} |k_i * f|$; the constants are finite as in \Cref{thm:tl-local-means}, and
$C^{\T}_{p,q} < \infty$ by \Cref{thm:fefferman-stein-periodic}.
\end{proof}

\section{Ball means of differences}
\label{sec:tl-ball-means}

This section describes $\TL{s}{p}{q}$ on $V$ by ball means of differences, in the range
$1 < p < \infty$, $1 < q \leq \infty$, $0 < s < L$. As in the previous sections, $g$ is a bump of
radius $R$, $m$ its moment kernel of order $L$, $k_j$ and $\eta_j$ are the local means kernels and
the dual kernels, $A_{L,K}$ are the constants of \Cref{lem:dual-kernel-derivatives}, and
$c^L_{ji} = 2^{-(j-i)L}$ for $i \leq j$, $c^L_{ji} = 1$ for $i > j$. Let $r = \lceil R \rceil$,
the least natural number at least $R \geq 0$, so that $R < 2^r$, and $B_j = B(0, 2^{-j})$. Local
means are differences of averaged differences, so they are bounded by ball means; conversely the
differences of $\eta_i * u$ are bounded by the maximal function of $u$, and the reproducing
formula, the almost diagonal estimate and the Fefferman--Stein inequality finish as for local
means.

\begin{definition}[Ball means of differences]
\label{def:ball-means-differences}
\lean{Rd.ballMeanDiff, Rd.one_div_two_pow_pos, Rd.two_pow_mul_volume_ball_one_div,
  Rd.setLIntegral_eq_mul_ballMeanDiff, Rd.tlDiffNorm, Rd.measurable_lqNorm_ballMeanDiff}
\leanok
\uses{def:rd-modulus, def:lq-norm, lem:lq-norm-means}
For $f : V \to \C$, $L, j \in \N$ and $x \in V$ the ball mean of differences of order $L$ at
level $j$ is
\[
  d_j f(x) = \frac{1}{|B_j|} \int_{B_j} |\Delta^L_h f(x)| \, dh \in [0, \infty] .
\]
For every $c > 0$ and $j \in \N$, $|B(0, c 2^{-j})| = c^n |B_j|$ (Mathlib's
\texttt{MeasureTheory.Measure.addHaar\_ball\_mul\_of\_pos}) and $2^{jn} |B_j| = |B(0,1)|$.
Let $1 \leq p < \infty$, $s \in \R$ and $1 \leq q \leq \infty$. For $f \in L_p$ the ball means
quasi-norm of Triebel--Lizorkin type is
\[
  \|f\|^{L}_{F^s_{p,q}, \Delta} = \|f\|_p + \Bigl\| x \mapsto \bigl\| (2^{js} d_j f(x))_{j \in \N}
  \bigr\|_{\ell_q} \Bigr\|_{L_p},
\]
and $x \mapsto \|(2^{js} d_j f(x))_j\|_{\ell_q}$ is measurable.
\end{definition}

\begin{lemma}[Elementary properties of ball means]
\label{lem:ball-means-basic}
\lean{Rd.ballMeanDiff_le_of_forall, Rd.ballMeanDiff_level_add_le,
  Rd.measurable_enorm_higherDiff_step,
  Rd.measurable_ballMeanDiff, Rd.ballMeanDiff_congr_ae, Rd.ballMeanDiff_add_le,
  Rd.ballMeanDiff_zero, Rd.ballMeanDiff_finsetSum_le}
\leanok
\uses{def:ball-means-differences, lem:rd-differences}
For all $f, f' : V \to \C$, all $L, j, k \in \N$, every $x \in V$ and every $c \in [0, \infty]$:
\begin{enumerate}
  \item if $|\Delta^L_h f(x)| \leq c$ for every $h \in B_j$, then $d_j f(x) \leq c$;
  \item $d_{j+k} f(x) \leq 2^{kn} d_j f(x)$;
  \item if $f$ is measurable, then $h \mapsto |\Delta^L_h f(x)|$ and $d_j f$ are measurable;
  \item if $f = f'$ almost everywhere, then $d_j f = d_j f'$ almost everywhere;
  \item if $f$ is measurable, then $d_j(f + f')(x) \leq d_j f(x) + d_j f'(x)$, and
    $d_j \bigl(\sum_{i \in S} f_i\bigr)(x) \leq \sum_{i \in S} d_j f_i(x)$ for every finite
    family of measurable $f_i : V \to \C$; $d_j 0 = 0$.
\end{enumerate}
\end{lemma}

\begin{proof}
\leanok
(1) is the mean of a bounded function, and (2) holds since $B_{j+k} \subseteq B_j$ and
$|B_j| = 2^{kn} |B_{j+k}|$. (3) The function $(x, h) \mapsto |\Delta^L_h f(x)|$ is jointly
measurable (\Cref{lem:rd-differences}), and Tonelli's theorem gives the measurability of
$d_j f$. (4) At a point $x$ with
$f(x) = f'(x)$ the values $f(x + l h)$ and $f'(x + l h)$, $1 \leq l \leq L$, agree for almost
every $h$, since $h \mapsto x + l h$ maps null sets to null sets. (5) The triangle inequality
and the additivity of the integral, which needs one measurable summand; finite sums by
induction.
\end{proof}

\begin{lemma}[Ball means in $L_p$]
\label{lem:ball-means-lp}
\lean{Rd.ballMeanDiff_rpow_le, Rd.eLpNorm_ballMeanDiff_le_of_measurable,
  Rd.eLpNorm_ballMeanDiff_le}
\leanok
\uses{lem:ball-means-basic, lem:weighted-jensen, lem:rd-differences, lem:lp-norm-nonnegative}
For every measurable $f : V \to \C$, all $L, j \in \N$, every $P \geq 1$ and every $x \in V$,
$(d_j f(x))^P \leq |B_j|^{-1} \int_{B_j} |\Delta^L_h f(x)|^P \, dh$. For all $L, j \in \N$,
every $1 \leq p < \infty$ and every almost everywhere strongly measurable $f : V \to \C$,
$\|d_j f\|_p \leq 2^L \|f\|_p$.
\end{lemma}

\begin{proof}
\leanok
The first bound is \Cref{lem:weighted-jensen} for the probability measure
$|B_j|^{-1} \mathbf 1_{B_j}(h) \, dh$. Integrating over $x$ and exchanging the integrals,
$\int (d_j f)^p \leq |B_j|^{-1} \int_{B_j} \|\Delta^L_h f\|_p^p \, dh \leq (2^L \|f\|_p)^p$ by
\Cref{lem:rd-differences}. A general $f$ agrees almost everywhere with a measurable function, and
\Cref{lem:ball-means-basic} (4) transfers the bound.
\end{proof}

\begin{theorem}[Local means by ball means]
\label{thm:local-means-by-ball-means}
\lean{Rd.LocalMeansBump.iSup_enorm_lt_top, Rd.LocalMeansBump.ballMeanConst,
  Rd.LocalMeansBump.ballMeanConst_ne_top, Rd.ofReal_abs_inv_neg_inv_two_pow_pow,
  Rd.LocalMeansBump.smul_closedBall_subset_ball,
  Rd.LocalMeansBump.lintegral_enorm_mul_higherDiff_le,
  Rd.LocalMeansBump.enorm_localMean_le_ballMeanDiff}
\leanok
\uses{def:local-means, def:ball-means-differences, lem:conv-second-argument, lem:conv-assoc,
  lem:change-of-scale}
Let
\[
  C^\Delta = \sup |g| \, |B(0,1)| \, 2^{rn} (1 + 2^n);
\]
then $C^\Delta < \infty$. For every $L \in \N$, every $1 \leq p < \infty$, every $f \in L_p$,
almost every $x \in V$ and every $j \in \N$,
\[
  |(k_{j+r+1} * f)(x)| \leq C^\Delta \, d_j f(x) .
\]
\end{theorem}

\begin{proof}
\leanok
Let $J = j + r$. Since $k_{J+1} = m_{2^{J+1}} - m_{2^J}$,
$k_{J+1} * f = (f - m_{2^J} * f) - (f - m_{2^{J+1}} * f)$ almost everywhere, and by
\Cref{lem:conv-second-argument}, for $c = 2^{J}$ and $c = 2^{J+1}$,
$|f(x) - (m_c * f)(x)| \leq \int |g(z)| \, |\Delta^L_{-z/c} f(x)| \, dz$ almost everywhere. The
bump vanishes outside $\bar B(0, R)$, and the substitution $h = -z/c$
(\Cref{lem:change-of-scale}) gives at most $\sup |g| \, c^n \int_{\bar B(0, R/c)}
|\Delta^L_h f(x)| \, dh$. Since $R < 2^r$, $\bar B(0, R/c) \subseteq B_j$, and
$c^n |B_j|$ is $2^{rn} |B(0,1)|$ for $c = 2^J$ and $2^{(r+1)n} |B(0,1)|$ for $c = 2^{J+1}$. The
sum of the two bounds is $C^\Delta d_j f(x)$.
\end{proof}

Young's inequality bounds the first $r + 1$ local means by $\|f\|_p$
(\Cref{thm:tl-ball-means}).

\begin{lemma}[Tools for ball means of convolutions]
\label{lem:moving-balls}
\lean{Rd.enorm_le_inv_two_pow_of_mem_ball, Rd.dyadicWeight_mul_dyadicWeight_neg_of_le,
  Rd.volume_setOf_dist_add_smul_lt_le, Rd.indicator_ball_add_smul_le,
  Rd.setLIntegral_indicator_ball_add_smul_le, Rd.setLIntegral_setLIntegral_ball_le}
\leanok
\uses{def:ball-means-differences}
For every $j \in \N$, every $h \in B_j$ has $|h| \leq 2^{-j}$. For all $i, j \in \N$ with
$i \leq j$ and every $L \in \R$, $2^{iL} 2^{-jL} = 2^{-(j-i)L}$. For every real $l \geq 1$, all
$a, \rho \in \R$, all $x, w \in V$ and every $G : V \to [0, \infty]$,
$|\{v : |x + l v - w| < \rho\}| \leq |B(0, \rho)|$ and
\[
  \int_{B(0,a)} 1_{B(x + l v, \rho)}(w) \, G(w) \, dv \leq |B(0, \rho)| \,
  1_{B(x, \rho + l a)}(w) \, G(w) ;
\]
if $G$ is almost everywhere measurable, then
\[
  \int_{B(0,a)} \int_{B(x + l v, \rho)} G(w) \, dw \, dv \leq |B(0, \rho)|
  \int_{B(x, \rho + l a)} G .
\]
\end{lemma}

\begin{proof}
\leanok
The set $\{v : |x + l v - w| < \rho\}$ is the ball $B((w - x)/l, \rho/l)$, of measure
$l^{-n} |B(0,\rho)|$. A point $w$ of $B(x + l v, \rho)$ with $|v| < a$ satisfies
$|w - x| < \rho + l a$ by the triangle inequality, and $v$ lies in this set, which gives the
single integral. For the double integral, exchange the integrals and apply the single integral at
every $w$.
\end{proof}

\begin{theorem}[Ball means of the dual kernels]
\label{thm:ball-means-dual}
\lean{Rd.LocalMeansBump.diffByLocalMeansConst, Rd.LocalMeansBump.diffByLocalMeansConst_ne_top,
  Rd.LocalMeansBump.enorm_higherDiff_dual_le_mul,
  Rd.LocalMeansBump.enorm_higherDiff_dual_le_indicator,
  Rd.LocalMeansBump.enorm_higherDiff_conv_dual_le_of_le,
  Rd.LocalMeansBump.ballMeanDiff_conv_dual_le_of_le, Rd.LocalMeansBump.enorm_dual_le_indicator,
  Rd.LocalMeansBump.enorm_conv_dual_le, Rd.LocalMeansBump.enorm_conv_dual_le_maximal,
  Rd.LocalMeansBump.laverage_enorm_conv_dual_add_smul_le_of_one_le,
  Rd.LocalMeansBump.laverage_enorm_conv_dual_add_smul_le,
  Rd.LocalMeansBump.ballMeanDiff_conv_dual_le_of_lt, Rd.LocalMeansBump.ballMeanDiff_conv_dual_le}
\leanok
\uses{lem:dual-kernel-derivatives, lem:dual-kernel-bounds, lem:dual-averaged-differences,
  lem:smooth-differences, lem:conv-pointwise, lem:moving-balls, lem:maximal-function-kernels,
  def:maximal-function, lem:ball-means-basic}
Let $L \in \N$ and
\[
  E_L = |B(0,1)| \max\bigl( A_{L,L} (2LR + L + 1)^n, \;
  2^L A_{L,0} (2LR + 1)^n (2LR + L + 1)^n \bigr);
\]
then $E_L < \infty$. For every locally integrable $u : V \to \C$, all $i, j \in \N$ and every
$x \in V$,
\[
  d_j(\eta_i * u)(x) \leq E_L \, c^L_{ji} \, M|u|(x) .
\]
\end{theorem}

\begin{proof}
\leanok
Case $i \leq j$. For $h \in B_j$, $\Delta^L_h (\eta_i * u) = (\Delta^L_h \eta_i) * u$
(\Cref{lem:conv-pointwise}), and by \Cref{lem:smooth-differences} and
\Cref{lem:dual-kernel-derivatives},
$|\Delta^L_h \eta_i| \leq A_{L,L} 2^{in} 2^{iL} |h|^L \leq A_{L,L} 2^{in} 2^{-(j-i)L}$;
the difference vanishes outside $B(0, (2LR + L + 1) 2^{-i})$
(\Cref{lem:dual-averaged-differences}). \Cref{lem:maximal-function-kernels} (1) bounds
$|\Delta^L_h (\eta_i * u)(x)|$ by $A_{L,L} (2LR + L + 1)^n |B(0,1)| 2^{-(j-i)L} M|u|(x)$ for
every $h$, hence also the mean $d_j(\eta_i * u)(x)$.

Case $i > j$. By \Cref{lem:conv-pointwise},
$|\Delta^L_h (\eta_i * u)(x)| \leq \sum_{l \leq L} \binom{L}{l} |(\eta_i * u)(x + l h)|$, and
$|\eta_i| \leq A_{L,0} 2^{in} \mathbf 1_{B(0, \rho_i)}$ with $\rho_i = (2LR + 1) 2^{-i}$
(\Cref{lem:dual-kernel-bounds}), so
$|(\eta_i * u)(v)| \leq A_{L,0} 2^{in} \int_{B(v, \rho_i)} |u|$ for every $v \in V$
(\Cref{lem:maximal-function-kernels} (1)). For $l = 0$ this is at most
$A_{L,0} (2LR + 1)^n |B(0,1)| M|u|(x)$. For $l \geq 1$ the mean over $h \in B_j$ is at most
$|B_j|^{-1} A_{L,0} 2^{in} |B(0, \rho_i)| \int_{B(x, (2LR + L + 1) 2^{-j})} |u|$ by
\Cref{lem:moving-balls}, since $\rho_i + l 2^{-j} \leq (2LR + L + 1) 2^{-j}$, and this is at most
$A_{L,0} (2LR + 1)^n (2LR + L + 1)^n |B(0,1)| M|u|(x)$. The binomial weights sum to $2^L$.
\end{proof}

For a single step $h$ the ball $B(x + l h, \rho_i)$ lies in $B(x, (2LR + L + 1) 2^{-j})$, which
costs the factor $2^{(i-j)n}$; the average over the steps $h \in B_j$ removes it.

\begin{theorem}[Ball means by local means]
\label{thm:ball-means-by-local-means}
\lean{Rd.LocalMeansBump.ballMeanDiff_partialSum_le, Rd.LocalMeansBump.ballMeanDiff_le_tsum_add,
  Rd.LocalMeansBump.tendsto_eLpNorm_ballMeanDiff_sub,
  Rd.LocalMeansBump.ballMeanDiff_le_tsum_maximalFunction}
\leanok
\uses{thm:ball-means-dual, lem:reproducing-formula, lem:reproducing-pointwise,
  lem:ball-means-basic, lem:ball-means-lp, lem:limits-of-pointwise-bounds,
  lem:maximal-function-basic, lem:maximal-function-measurable, def:tl-local-means, lem:conv-op}
Let $L \in \N$ with $L \geq 1$ and $1 \leq p < \infty$. For every $f \in L_p$, almost every
$x \in V$ and every $j \in \N$,
\[
  d_j f(x) \leq E_L \sum_{i \in \N} c^L_{ji} \, M|k_i * f|(x) .
\]
\end{theorem}

\begin{proof}
\leanok
Let $F_J = m_{2^J} * m_{2^J} * f = \sum_{i \leq J} \eta_i * (k_i * f)$, an identity of continuous
functions (\Cref{lem:reproducing-pointwise}). By \Cref{lem:ball-means-basic} (5) and
\Cref{thm:ball-means-dual}, with $M|k_i * f|$ unchanged when $k_i * f$ changes on a null set,
$d_j F_J(x) \leq E_L \sum_{i \leq J} c^L_{ji} M|k_i * f|(x)$ for every $x$, and
$d_j f \leq d_j F_J + d_j(f - F_J)$ everywhere. Since $F_J \to f$ in $L_p$
(\Cref{lem:reproducing-formula}), $\|d_j(f - F_J)\|_p \leq 2^L \|f - F_J\|_p \to 0$
(\Cref{lem:ball-means-lp}), and \Cref{lem:limits-of-pointwise-bounds} gives the bound almost
everywhere.
\end{proof}

\begin{theorem}[Ball means characterisation of Triebel--Lizorkin spaces]
\label{thm:tl-ball-means}
\lean{Rd.LocalMeansBump.enorm_le_tlLocalMeansNorm, Rd.LocalMeansBump.diffUpperConst,
  Rd.LocalMeansBump.diffUpperConst_ne_top, Rd.LocalMeansBump.tlDiffNorm_le_tlLocalMeansNorm,
  Rd.LocalMeansBump.diffLowerConst, Rd.LocalMeansBump.diffLowerConst_ne_top,
  Rd.LocalMeansBump.eLpNorm_dyadicWeight_mul_conv_kernel_le,
  Rd.LocalMeansBump.lqNorm_localMean_shift_le, Rd.LocalMeansBump.tlLocalMeansNorm_le_tlDiffNorm}
\leanok
\uses{thm:ball-means-by-local-means, thm:local-means-by-ball-means,
  lem:almost-diagonal-pointwise, lem:maximal-bounds-mixed-norm, lem:reproducing-formula,
  lem:dual-kernel-moduli, lem:conv-op, cor:mixed-norm-sandwich, lem:lq-basic,
  lem:local-means-synthesis, def:tl-local-means, def:ball-means-differences,
  lem:lp-norm-homogeneity}
Let $s \in \R$, $L \in \N$, $1 \leq p < \infty$ and $1 \leq q \leq \infty$, and let
\[
  D^\Delta_{L,s} = 2 \|m\|_1 \sum_{j \leq r} 2^{js} + 2^{(r+1)s} C^\Delta .
\]
Then $D^\Delta_{L,s} < \infty$, and for every $f \in L_p$,
\[
  \|f\|^{g,L}_{F^s_{p,q}} \leq D^\Delta_{L,s} \, \|f\|^{L}_{F^s_{p,q}, \Delta} .
\]
If $L \geq 1$, then $\|f\|_p \leq 2 \|m\|_1 (1 - 2^{-s})^{-1} \|f\|^{g,L}_{F^s_{p,q}}$ for every
$f \in L_p$. If $L \geq 1$, $1 < p$ and $1 < q$, let
\[
  U^\Delta_{L,p,q,s} = 2 \|m\|_1 (1 - 2^{-s})^{-1} + E_L A_{s,L} C_{n,p,q},
\]
with the constants $A_{s,L}$ of \Cref{lem:almost-diagonal} and $C_{n,p,q}$ of
\Cref{thm:fefferman-stein}; then $U^\Delta_{L,p,q,s} < \infty$ for $0 < s < L$, and for every
$f \in L_p$,
\[
  \|f\|^{L}_{F^s_{p,q}, \Delta} \leq U^\Delta_{L,p,q,s} \, \|f\|^{g,L}_{F^s_{p,q}} .
\]
\end{theorem}

\begin{proof}
\leanok
Lower bound: pointwise
$\|(2^{js} |k_j * f(x)|)_j\|_{\ell_q} \leq \sum_{j \leq r} 2^{js} |k_j * f(x)| +
\|(2^{(j+r+1)s} |k_{j+r+1} * f(x)|)_j\|_{\ell_q}$ (\Cref{lem:lq-basic}); the head is bounded by
Young's inequality $\|k_j * f\|_p \leq 2 \|m\|_1 \|f\|_p$ (\Cref{lem:conv-op},
\Cref{lem:dual-kernel-moduli}), the tail by \Cref{thm:local-means-by-ball-means}, almost
everywhere, after $2^{(j+r+1)s} = 2^{(r+1)s} 2^{js}$.

Upper bound: the reproducing formula $f = \sum_j \eta_j * k_j * f$ (\Cref{lem:reproducing-formula})
with $\|\eta_j\|_1 \leq 2 \|m\|_1$ gives
$\|f\|_p \leq 2 \|m\|_1 (1 - 2^{-s})^{-1} \sup_j 2^{js} \|k_j * f\|_p$ (the series bound of
\Cref{lem:local-means-synthesis}), and $\sup_j 2^{js} \|k_j * f\|_p \leq \|f\|^{g,L}_{F^s_{p,q}}$
by \Cref{cor:mixed-norm-sandwich} at $q = \infty$ and $\ell_q \subseteq \ell_\infty$.
\Cref{thm:ball-means-by-local-means} and \Cref{lem:almost-diagonal-pointwise} bound
$\|(2^{js} d_j f(x))_j\|_{\ell_q}$ by $E_L A_{s,L} \|(M(2^{is} |k_i * f|)(x))_i\|_{\ell_q}$
almost everywhere, and \Cref{lem:maximal-bounds-mixed-norm} finishes.
\end{proof}

\begin{corollary}[Ball means describe $\TL{s}{p}{q}$]
\label{cor:tl-ball-means}
\lean{Rd.memTriebelLizorkin_iff_tlDiffNorm_lt_top}
\leanok
\uses{thm:tl-ball-means, thm:tl-local-means, thm:smooth-atomic-decomposition}
For every $1 < p < \infty$, every $1 < q \leq \infty$, every $L \in \N$, every $0 < s < L$ and
every $f \in L_p$, $f \in \TL{s}{p}{q}$ if and only if $\|f\|^{L}_{F^s_{p,q}, \Delta} < \infty$.
\end{corollary}

\begin{proof}
\leanok
$L \geq 1$ since $0 < s < L$. Both conditions are equivalent to the finiteness of the local means
quasi-norm of the standard bump (\Cref{thm:smooth-atomic-decomposition}), by
\Cref{thm:tl-local-means} and \Cref{thm:tl-ball-means}.
\end{proof}

\section{The atomic characterisation}
\label{sec:tl-atomic}

This section describes $\TL{s}{p}{q}$ by atomic decompositions, in the range $1 < p < \infty$,
$1 < q \leq \infty$, $0 < s < L$. As in the previous sections, $g$ is a bump of radius $R$,
$m^g$ its moment kernel of order $L$, $A_L = 2 \sup_V |m^g|$ the constant of
\Cref{lem:local-means-kernel-bounds}, $k_j$ and $\eta_j$ are the local means kernels and the dual
kernels, and $c^L_{ij} = 2^{-(i-j)L}$ for $j \leq i$, $c^L_{ij} = 1$ for $j > i$. The cells
$P_{j,m}$ and their centres $c_{j,m}$ are those of \Cref{def:dyadic-cells}, $|B_1|$ is the
volume of the unit ball, and $A_+ = \max(A, 0)$ for $A \in \R$. The atoms are the smooth atoms
of \Cref{def:smooth-atoms}, normalised in $L_\infty$; their constants are nonnegative
(\Cref{lem:smooth-atoms-basic}), so $A_+ = A$ for them. The coefficient of a cell in the
decomposition of $F$ is the mean of $|k_j * F|$ over the cell, at most a multiple of
$M(k_j * F)$ at every point of the cell. The Besov spaces use atoms normalised in $L_p$
(\Cref{thm:atomic-characterisation}); here the maximal function bounds means in $L_1$, so the
coefficients are $L_1$ means over the cells. Triebel~\cite{TriebelII} uses $(s,p)$-atoms
normalised in $L_p$; his coefficients are $2^{j(s - n/p)} \lambda_{j,m}$ for the coefficients
$\lambda_{j,m}$ here.

\begin{definition}[The sequence space $f^s_{p,q}$]
\label{def:tl-sequence-space}
\lean{Rd.cellCoefFun, Rd.cellCoefFun_of_mem, Rd.measurable_cellCoefFun, Rd.tlCoeffNorm,
  Rd.aemeasurable_lqNorm_cellCoefFun}
\leanok
\uses{def:dyadic-cells, def:lq-norm, lem:lp-norm-nonnegative}
Let $\lambda = (\lambda_{j,m})_{j \in \N, m \in \Z^n}$ be complex coefficients. For $j \in \N$ and
$x \in V$ let $\lambda_j(x) = |\lambda_{j,m}|$ for the index $m$ of the cell $P_{j,m}$ of level
$j$ that contains $x$. The function $\lambda_j : V \to [0, \infty)$ is measurable. For
$1 \leq p < \infty$, $s \in \R$ and $1 \leq q \leq \infty$,
\[
  \|\lambda\|_{f^s_{p,q}} = \Bigl\| x \mapsto \bigl\| (2^{js} \lambda_j(x))_{j \in \N}
  \bigr\|_{\ell_q} \Bigr\|_{L_p} ,
\]
and $x \mapsto \|(2^{js} \lambda_j(x))_j\|_{\ell_q}$ is measurable.
\end{definition}

\begin{lemma}[Coefficients and the maximal function]
\label{lem:coefficients-maximal}
\lean{Rd.dyadicWeight_mul_setLIntegral_ball_le, Rd.cellCoefFun_le_maximalFunction,
  Rd.enorm_coef_le_maximalFunction_cellCoefFun, Rd.dyadicWeight_mul_tsum_ite_le_maximalFunction,
  ENNReal.tsum_mul_le_of_le}
\leanok
\uses{def:tl-sequence-space, def:maximal-function, lem:dyadic-ball-volume, lem:dyadic-cubes,
  lem:dyadic-cells, lem:cells-partition, lem:product-modulus}
Let $\lambda$ be complex coefficients and $j \in \N$.
\begin{enumerate}
  \item For every $m \in \Z^n$, $|P_{j,m}| = 2^{-jn}$; for every $\rho > 0$, every
    $u : V \to [0, \infty]$ and every $x \in V$,
    $2^{jn} \int_{B(x, \rho 2^{-j})} u \leq \rho^n |B_1| \, M u(x)$.
  \item For every $u : V \to [0, \infty]$ with $|\lambda_{j,m}| \leq 2^{jn} \int_{P_{j,m}} u$
    for every $m \in \Z^n$, and every $x \in V$,
    $\lambda_j(x) \leq (2n + 1)^n |B_1| \, M u(x)$.
  \item For every $\rho \in \R$, every $m \in \Z^n$ and every $x \in V$ with
    $\|x - c_{j,m}\| \leq \rho 2^{-j}$, $|\lambda_{j,m}| \leq (2(\rho + 1))^n M \lambda_j(x)$.
  \item For every $i \leq j$, every $\rho \geq 0$ and every $x \in V$,
    \[
      2^{-jn} \sum_{m : \|x - c_{j,m}\| \leq \rho 2^{-i}} |\lambda_{j,m}| \leq
      (\rho + n + 1)^n |B_1| \, 2^{-in} M \lambda_j(x) .
    \]
  \item For every family $(\mu_m)$ in $[0, \infty]$, every predicate $P$ on the indices, every
    $B \in [0, \infty]$ and all $u_m \in [0, \infty]$ with $u_m \leq B \, \mathbf 1_{P(m)}$,
    $\sum_m \mu_m u_m \leq B \sum_{m : P(m)} \mu_m$.
\end{enumerate}
\end{lemma}

\begin{proof}
\leanok
(1) The cell is a cube of side $2^{-j}$ (\Cref{lem:dyadic-cubes}), the volume of balls scales
with the $n$-th power of the radius (\Cref{lem:dyadic-ball-volume}),
$(2^{-j})^n = 2^{-jn}$ (\Cref{lem:product-modulus}), and
\Cref{def:maximal-function}. (2) The cell of $x$ lies in
$B(x, (2n + 1) 2^{-j})$, since all its points have distance at most $n 2^{-j}$ from its centre
(\Cref{lem:dyadic-cells}); apply (1). (3) $\rho \geq 0$ since $\|x - c_{j,m}\| \geq 0$. The ball
$B(c_{j,m}, 2^{-j}/2)$ lies in $P_{j,m}$
(\Cref{lem:dyadic-cubes}), where $\lambda_j = |\lambda_{j,m}|$, and in
$B(x, (\rho + 1) 2^{-j})$, whose volume is $(2(\rho + 1))^n$ times as large. (4) The cells in the
sum lie in $B(x, (\rho + n + 1) 2^{-i})$, and $2^{-jn} |\lambda_{j,m}| = \int_{P_{j,m}} \lambda_j$;
the cells are disjoint (\Cref{lem:cells-partition}), and (1) at level $i$ finishes. (5) is
termwise.
\end{proof}

\begin{lemma}[Smooth atoms on $V$]
\label{lem:smooth-atoms-basic}
\lean{Rd.IsSmoothAtom.norm_le, Rd.IsSmoothAtom.const_nonneg, Rd.IsSmoothAtom.continuous,
  Rd.IsSmoothAtom.hasCompactSupport, Rd.IsSmoothAtom.integrable,
  Rd.IsSmoothAtom.lintegral_enorm_le}
\leanok
\uses{def:smooth-atoms, lem:pointwise-atoms}
Let $a$ be a smooth atom on a set $\Omega$ of level $j$ with centre $c$, radius factor $d$ and
constant $A$. Then $|a| \leq A$ on $\Omega$, and $a$ has compact support. If $\Omega = V$, then
$A \geq 0$, $a$ is continuous and integrable, and for $d \geq 0$,
$\int_V |a| \leq A \, d^n |B_1| \, 2^{-jn}$.
\end{lemma}

\begin{proof}
\leanok
The bound is the case $l = 0$ of \Cref{def:smooth-atoms}; for $\Omega = V$ it gives
$0 \leq |a(0)| \leq A$. The atom vanishes outside $\overline B(c, d 2^{-j})$, whose volume is
$d^n |B_1| 2^{-jn}$ (\Cref{lem:pointwise-atoms}).
\end{proof}

\begin{lemma}[Convolution with a series of atoms]
\label{lem:conv-atom-series}
\lean{Rd.conv_tsum_atom, Rd.enorm_conv_le_tsum_atom}
\leanok
\uses{lem:smooth-atoms-basic, lem:cells-partition, lem:bounded-overlap, lem:cell-decomposition,
  def:conv}
Let $K$ be continuous with $K = 0$ outside $\overline B(0, \rho)$, $\rho \geq 0$, let $d \geq 0$,
and for every $m \in \Z^n$ let $a_m$ be a smooth atom on $V$ of level $j$ with centre $c_{j,m}$,
radius factor $d$ and constant $A$. For all complex $\lambda_m$ and every $x \in V$,
\[
  \Bigl(K * \sum_m \lambda_m a_m\Bigr)(x) = \sum_m \lambda_m (K * a_m)(x) ,
\]
and for every $f : V \to \C$ with $f = \sum_m \lambda_m a_m$ almost everywhere,
$|(K * f)(x)| \leq \sum_m |\lambda_m| \, |(K * a_m)(x)|$ for every $x \in V$.
\end{lemma}

\begin{proof}
\leanok
If $K(y) \neq 0$ and $a_m(x - y) \neq 0$, then $\|x - c_{j,m}\| \leq \rho + d 2^{-j}$; only
finitely many centres satisfy this (\Cref{lem:cells-partition}, \Cref{lem:bounded-overlap}), so
\Cref{lem:cell-decomposition} (3) applies to $u_m = \lambda_m a_m$. The bound follows since
$K * f = K * \sum_m \lambda_m a_m$ (\Cref{def:conv}).
\end{proof}

\begin{lemma}[Coarse local means of atoms]
\label{lem:coarse-local-means-atoms}
\lean{Rd.LocalMeansBump.lintegral_enorm_kernel_mul_atom_le,
  Rd.LocalMeansBump.tlAtomCoarseConst, Rd.LocalMeansBump.tlAtomCoarseConst_ne_top,
  Rd.LocalMeansBump.enorm_conv_kernel_level_le_of_le}
\leanok
\uses{lem:local-means-kernel-bounds, lem:reproducing-pointwise, lem:dyadic-weights-integer,
  lem:smooth-atoms-basic, lem:conv-atom-series, lem:coefficients-maximal, lem:weighted-jensen}
Let $L \in \N$, $d \geq 0$ and $A \in \R$, and let
\[
  C^{\mathrm c}_{L,d,A} = A_L \, A_+ \, d^n |B_1| \, (2LR + d + n + 1)^n |B_1| ;
\]
then $C^{\mathrm c}_{L,d,A} < \infty$. For all $i \leq j$, every smooth atom $a$ on $V$ of level
$j$ with centre $c$, radius factor $d$ and constant $A$, and every $x \in V$,
\[
  \int_V |k_i(y)| \, |a(x - y)| \, dy \leq A_L \, A_+ \, d^n |B_1| \, 2^{in} 2^{-jn}
  \mathbf 1_{\|x - c\| \leq (2LR + d) 2^{-i}} .
\]
For all $i \leq j$, all coefficients $\lambda$, smooth atoms $a_{j,m}$ on $V$ of level $j$ with
centres $c_{j,m}$, radius factor $d$ and constant $A$, every $f : V \to \C$ with
$f = \sum_m \lambda_{j,m} a_{j,m}$ almost everywhere and every $x \in V$,
\[
  |(k_i * f)(x)| \leq C^{\mathrm c}_{L,d,A} \, M \lambda_j(x) .
\]
\end{lemma}

\begin{proof}
\leanok
The kernel $k_i$ is continuous (\Cref{lem:reproducing-pointwise}), and
$|k_i| \leq A_L 2^{in}$ vanishes outside $\overline B(0, 2LR 2^{-i})$
(\Cref{lem:local-means-kernel-bounds}, \Cref{lem:dyadic-weights-integer}). Moreover
$\int |a| \leq A d^n |B_1| 2^{-jn}$ with $A = A_+$ (\Cref{lem:smooth-atoms-basic}); if
$\|x - c\| > (2LR + d) 2^{-i}$, then $\|x - y - c\| > d 2^{-i} \geq d 2^{-j}$ wherever
$k_i(y) \neq 0$. By \Cref{lem:conv-atom-series} and \Cref{lem:weighted-jensen},
$|(k_i * f)(x)| \leq A_L A_+ d^n |B_1| 2^{in} \cdot 2^{-jn}
\sum_{m : \|x - c_{j,m}\| \leq (2LR + d) 2^{-i}} |\lambda_{j,m}|$, and
\Cref{lem:coefficients-maximal} (4) bounds the last factor by
$(2LR + d + n + 1)^n |B_1| 2^{-in} M \lambda_j(x)$.
\end{proof}

\begin{lemma}[Fine local means of atoms]
\label{lem:fine-local-means-atoms}
\lean{Rd.IsSmoothAtom.higherDiff_eq_zero_of_lt, Rd.LocalMeansBump.enorm_mul_higherDiff_atom_le,
  Rd.LocalMeansBump.ae_enorm_atom_sub_conv_scaled_le, Rd.LocalMeansBump.momentAtomConst,
  Rd.LocalMeansBump.momentAtomConst_ne_top,
  Rd.LocalMeansBump.ae_enorm_conv_kernel_atom_le_of_lt, Rd.LocalMeansBump.tlAtomFineConst,
  Rd.LocalMeansBump.tlAtomFineConst_ne_top,
  Rd.LocalMeansBump.ae_enorm_conv_kernel_level_le_of_lt}
\leanok
\uses{lem:smooth-atoms-regular, lem:local-means-one-scale, lem:conv-second-argument,
  def:local-means, lem:local-means-kernel-bounds, lem:reproducing-pointwise, def:rd-modulus,
  def:rd-modulus-on, lem:smooth-atoms-basic, lem:conv-atom-series, lem:coefficients-maximal}
Let $L \in \N$ with $L \geq 1$, $d \geq 0$ and $A \in \R$, and let
\[
  D_{L,A} = 2 \|g\|_1 \, 2^L L^{L+1} A_+ \, (2R)^L , \qquad
  C^{\mathrm f}_{L,d,A} = D_{L,A} \, (d + LR + n + 1)^n |B_1| ;
\]
then both are finite. For all $j < i$ and every smooth atom $a$ on $V$ of level $j$ and order
$L$ with centre $c$, radius factor $d$ and constant $A$, almost everywhere
\[
  |(k_i * a)(x)| \leq D_{L,A} \, 2^{-(i-j)L} \, \mathbf 1_{\|x - c\| \leq (d + LR) 2^{-j}} .
\]
For all $j < i$, all coefficients $\lambda$, smooth atoms $a_{j,m}$ on $V$ of level $j$ and order
$L$ with centres $c_{j,m}$, radius factor $d$ and constant $A$, and every $f : V \to \C$ with
$f = \sum_m \lambda_{j,m} a_{j,m}$ almost everywhere, almost everywhere
\[
  |(k_i * f)(x)| \leq C^{\mathrm f}_{L,d,A} \, 2^{-(i-j)L} \, M \lambda_j(x) .
\]
\end{lemma}

\begin{proof}
\leanok
By \Cref{lem:local-means-one-scale} it suffices to show, for every $k \geq j$ and almost every
$x$,
\[
  |a(x) - (m^g_{2^k} * a)(x)| \leq \|g\|_1 \, 2^L L^{L+1} A_+ R^L \, 2^{-(k-j)L}
  \mathbf 1_{\|x - c\| \leq (d + LR) 2^{-j}} ,
\]
since $2 \cdot 2^L$ times this constant is $D_{L,A}$. By \Cref{lem:conv-second-argument},
$|a(x) - (m^g_{2^k} * a)(x)| \leq \int |g(z)| \, |\Delta^L_{-2^{-k} z} a(x)| \, dz$ almost
everywhere. On the support of $g$ the step $h = -2^{-k} z$ has $\|h\| \leq R 2^{-k}$, so
$|\Delta^L_h a(x)| \leq 2^L L^{L+1} A (2^j \|h\|)^L \leq 2^L L^{L+1} A R^L 2^{-(k-j)L}$
(\Cref{lem:smooth-atoms-regular}, with the steps of \Cref{def:rd-modulus-on}). The difference
is a sum of values $a(x + l h)$, $l \leq L$ (\Cref{def:rd-modulus}), which vanish when
$\|x - c\| > d 2^{-j} + L \|h\|$; since $\|h\| \leq R 2^{-j}$, they vanish unless
$\|x - c\| \leq (d + LR) 2^{-j}$. For the level, \Cref{lem:conv-atom-series} and
\Cref{lem:coefficients-maximal} (5) give
$|(k_i * f)(x)| \leq D_{L,A} 2^{-(i-j)L} \sum_{m : \|x - c_{j,m}\| \leq (d + LR) 2^{-j}}
|\lambda_{j,m}|$, and \Cref{lem:coefficients-maximal} (4) with $i = j$ bounds the sum by
$(d + LR + n + 1)^n |B_1| M \lambda_j(x)$.
\end{proof}

The fine local means need no moment conditions on the atoms: the moments of the kernel $k_i$ and
the smoothness of the atoms give the decay $2^{-(i-j)L}$.

\begin{theorem}[Synthesis from atoms]
\label{thm:tl-atomic-synthesis}
\lean{Rd.LocalMeansBump.tlAtomSynthesisConst, Rd.LocalMeansBump.tlAtomSynthesisConst_ne_top,
  Rd.LocalMeansBump.ae_enorm_conv_kernel_level_le,
  Rd.LocalMeansBump.enorm_conv_kernel_le_tsum_level,
  Rd.LocalMeansBump.enorm_localMean_le_tsum_maximalFunction_cellCoefFun,
  Rd.LocalMeansBump.tlLocalMeansNorm_le_tlCoeffNorm}
\leanok
\uses{lem:coarse-local-means-atoms, lem:fine-local-means-atoms, lem:lp-limits-pointwise,
  def:tl-local-means, lem:conv-op, lem:almost-diagonal-pointwise, lem:maximal-bounds-mixed-norm,
  def:tl-sequence-space}
Let $L \in \N$ with $L \geq 1$, $d \geq 0$, $A \in \R$, and
$C^{\mathrm a}_{L,d,A} = \max(C^{\mathrm c}_{L,d,A}, C^{\mathrm f}_{L,d,A}) < \infty$. Let
$1 < p < \infty$, $1 < q \leq \infty$ and $s \in \R$. For all coefficients $\lambda$, all smooth
atoms $a_{j,m}$ on $V$ of level $j$ and order $L$ with centres $c_{j,m}$, radius factor $d$ and
constant $A$, all $f_j \in L_p$ with $f_j = \sum_m \lambda_{j,m} a_{j,m}$ almost everywhere, and
every $F \in L_p$ with $\sum_{j \leq J} f_j \to F$ in $L_p$, almost everywhere and for all
$i \in \N$
\[
  |(k_i * F)(x)| \leq C^{\mathrm a}_{L,d,A} \sum_{j \in \N} c^L_{ij} \, M \lambda_j(x) ,
\]
and
\[
  \|F\|^{g,L}_{F^s_{p,q}} \leq C^{\mathrm a}_{L,d,A} \, A_{s,L} \, C_{n,p,q} \,
  \|\lambda\|_{f^s_{p,q}} ,
\]
with the constants $A_{s,L}$ of \Cref{lem:almost-diagonal-pointwise}, finite for $0 < s < L$,
and $C_{n,p,q}$ of \Cref{thm:fefferman-stein}.
\end{theorem}

\begin{proof}
\leanok
\Cref{lem:coarse-local-means-atoms} for $j \geq i$ and \Cref{lem:fine-local-means-atoms} for
$j < i$ give $|(k_i * f_j)(x)| \leq C^{\mathrm a}_{L,d,A} c^L_{ij} M \lambda_j(x)$ almost
everywhere. The local mean $k_i *$ is continuous on $L_p$ (\Cref{lem:conv-op}), so
$k_i * \sum_{j \leq J} f_j \to k_i * F$, and the bound
$|k_i * \sum_{j \leq J} f_j| \leq \sum_j |k_i * f_j|$ of the partial sums passes to $F$
(\Cref{lem:lp-limits-pointwise}). \Cref{lem:almost-diagonal-pointwise} with the local means index
$i$ and the level $j$ and \Cref{lem:maximal-bounds-mixed-norm} for the functions
$2^{js} \lambda_j$ finish.
\end{proof}

\begin{lemma}[Coefficients with exponent one]
\label{lem:cell-coefficients-exponent-one}
\lean{Rd.enorm_cellCoef_one_le, Rd.lqNorm_cellCoefFun_le_lqNorm_maximalFunction}
\leanok
\uses{def:cell-atoms, lem:coefficients-maximal, lem:almost-diagonal-pointwise, lem:lq-basic}
For every almost everywhere strongly measurable $u : V \to \C$, every $j \in \N$ and every
$m \in \Z^n$, the coefficient $\lambda_{j,m}$ of $u$ with exponent $p = 1$
(\Cref{def:cell-atoms}) satisfies $|\lambda_{j,m}| \leq 2^{jn} \int_{P_{j,m}} |u|$. For all
almost everywhere strongly measurable $u_j : V \to \C$, $j \in \N$, the coefficients
$\lambda_{j,m}$ of $u_j$ with exponent one satisfy, for every $q$, every
$s \in \R$ and every $x \in V$,
\[
  \|(2^{js} \lambda_j(x))_j\|_{\ell_q} \leq (2n + 1)^n |B_1| \,
  \|(M(2^{js} |u_j|)(x))_j\|_{\ell_q} .
\]
\end{lemma}

\begin{proof}
\leanok
$\lambda_{j,m} = 2^{jn} \|1_{P_{j,m}} u\|_1$ when this is finite, and $0$ otherwise. The second
bound is \Cref{lem:coefficients-maximal} (2) for every level, with the weights moved inside the
maximal function (\Cref{lem:almost-diagonal-pointwise}).
\end{proof}

\begin{theorem}[Atomic decomposition of Triebel--Lizorkin type]
\label{thm:tl-atomic-decomposition}
\lean{Rd.LocalMeansBump.tl_atomic_decomposition, Rd.LocalMeansBump.tlCoeffNorm_cellCoef_le}
\leanok
\uses{lem:cell-atoms-smooth, lem:cell-decomposition, def:cell-atoms, lem:reproducing-formula,
  def:tl-local-means, lem:cell-coefficients-exponent-one, lem:maximal-bounds-mixed-norm,
  thm:young, def:local-means}
Let $L \in \N$ with $L \geq 1$, $K \in \N$, $1 \leq p < \infty$ and $F \in L_p$, and let
$u_j = k_j * F$. Then $u_j \in L_p$. For $j \in \N$ and $m \in \Z^n$ let
\[
  \lambda_{j,m} = 2^{jn} \int_{P_{j,m}} |u_j| , \qquad
  a_{j,m} = \lambda_{j,m}^{-1} \, \eta_j * (1_{P_{j,m}} u_j) ,
\]
with $a_{j,m} = 0$ if $\lambda_{j,m} = 0$: the coefficients and cell atoms with exponent $p = 1$
of \Cref{def:cell-atoms} for the kernel $\eta_j$ and $g = u_j$. Then:
\begin{enumerate}
  \item every $a_{j,m}$ is a smooth atom on $V$ of level $j$ and order $K$ with centre $c_{j,m}$,
    radius factor $2LR + n$ and constant $A_{L,K}$ of \Cref{lem:dual-kernel-derivatives};
  \item $(\eta_j * u_j)(x) = \sum_m \lambda_{j,m} a_{j,m}(x)$ for every $j \in \N$ and every
    $x \in V$;
  \item $\sum_{j \leq J} \eta_j * u_j \to F$ in $L_p$;
  \item if $p > 1$, then for every $1 < q \leq \infty$ and every $s \in \R$,
    \[
      \|\lambda\|_{f^s_{p,q}} \leq (2n + 1)^n |B_1| \, C_{n,p,q} \, \|F\|^{g,L}_{F^s_{p,q}} .
    \]
\end{enumerate}
\end{theorem}

\begin{proof}
\leanok
$u_j \in L_p$ by Young's inequality (\Cref{thm:young}), so $1_{P_{j,m}} u_j \in L_1$
(\Cref{lem:cell-decomposition}) and $\lambda_{j,m}$ is the finite integral. (1) is
\Cref{lem:cell-atoms-smooth} at $p = 1$, where the factor $((2LR)^n |B_1|)^{1 - 1/p}$ is $1$,
applied to $1_{P_{j,m}} u_j$; the atom of a cell depends only on $u_j$ on that cell
(\Cref{def:cell-atoms}). (2) is \Cref{lem:cell-decomposition} (4), and (3) is the reproducing
formula (\Cref{lem:reproducing-formula}). (4) is \Cref{lem:cell-coefficients-exponent-one} for
the continuous $u_j$ and \Cref{lem:maximal-bounds-mixed-norm} for the functions $2^{js} |u_j|$.
\end{proof}

The bound (4) uses neither $L \geq 1$ nor the order $K$.

\begin{corollary}[Atoms describe $\TL{s}{p}{q}$]
\label{cor:tl-atomic}
\lean{Rd.memTriebelLizorkin_iff_atomic}
\leanok
\uses{thm:tl-atomic-synthesis, thm:tl-atomic-decomposition, thm:tl-local-means,
  thm:smooth-atomic-decomposition}
For every $1 < p < \infty$, every $1 < q \leq \infty$, every $L \in \N$, every $0 < s < L$ and
every $F \in L_p$, $F \in \TL{s}{p}{q}$ if and only if there are $\rho \geq 0$, $A \in \R$,
complex coefficients $\lambda_{j,m}$, smooth atoms $a_{j,m}$ on $V$ of level $j$ and order $L$
with centres $c_{j,m}$, radius factor $\rho$ and constant $A$, and $f_j \in L_p$ with
$f_j = \sum_m \lambda_{j,m} a_{j,m}$ almost everywhere, $\sum_{j \leq J} f_j \to F$ in $L_p$ and
$\|\lambda\|_{f^s_{p,q}} < \infty$.
\end{corollary}

\begin{proof}
\leanok
$L \geq 1$ since $0 < s < L$, and $F \in \TL{s}{p}{q}$ iff the local means quasi-norm of a bump is
finite (\Cref{thm:tl-local-means}, with the bump of \Cref{thm:smooth-atomic-decomposition}).
Forward: \Cref{thm:tl-atomic-decomposition} with $K = L$, $\rho = 2LR + n$, $A = A_{L,L}$ and
$f_j = \eta_j * (k_j * F)$. Backward: \Cref{thm:tl-atomic-synthesis}, since
$A_{s,L} < \infty$.
\end{proof}

The Besov analogue is \Cref{thm:atomic-characterisation}. Here the atoms are indexed by the cells
and centred at their centres, which takes the place of the bounded overlap condition there, and
the atoms of the decomposition are smooth of every order.

\section{Triebel--Lizorkin spaces as interpolation spaces}
\label{sec:tl-interpolation}

The Besov sandwich (\Cref{thm:tl-besov-sandwich}) places the quasi-norm of $\TL{s}{p}{q}$
between two Besov quasi-norms of the same smoothness, so the sandwich principle
(\Cref{cor:approximation-space-sandwich}) identifies the real interpolation spaces of two
Triebel--Lizorkin spaces, and of $L_p$ and one Triebel--Lizorkin space, with Besov spaces; see
Triebel~\cite[Section~2.4]{TriebelInterpolation}. \Cref{sec:bessel-potential-interpolation}
draws the same conclusion for the Bessel potential spaces in their own norms. In this section
$1 \leq p < \infty$, and $\Besov{s}{p}{q}$ is the approximation space of the band-limited scheme
in $L_p = L_p(V)$.

\begin{lemma}[The upper constant of the Besov sandwich is finite]
\label{lem:tl-sandwich-constants}
\lean{Rd.eight_mul_dyadicWeight_mul_eLpNorm_lpKernel_ne_top}
\leanok
\uses{thm:tl-besov-sandwich, lem:lp-kernels-basic}
For every $s \in \R$, $8 \cdot 8^s \, \|\Phi\|_1 < \infty$.
\end{lemma}

\begin{proof}
\leanok
$\Phi$ is integrable (\Cref{lem:lp-kernels-basic}).
\end{proof}

The lower constant $2 (1 - 2^{-s})^{-1}$ is finite for every $s > 0$, since $2^{-s} < 1$
(\Cref{def:lq-norm}).

\begin{definition}[Triebel--Lizorkin couples]
\label{def:tl-couples}
\lean{Rd.triebelLizorkinCouple, Rd.triebelLizorkinCouple_norm0, Rd.triebelLizorkinCouple_norm1,
  Rd.lpTriebelLizorkinCouple, Rd.lpTriebelLizorkinCouple_norm0,
  Rd.lpTriebelLizorkinCouple_norm1}
\leanok
\uses{def:couple, def:triebel-lizorkin}
For all $s_0, s_1 \in \R$ and $q_0, q_1 \in [1, \infty]$, the couple
$(\TL{s_0}{p}{q_0}, \TL{s_1}{p}{q_1})$ on $L_p$ is
$(\|\cdot\|_{F^{s_0}_{p,q_0}}, \|\cdot\|_{F^{s_1}_{p,q_1}})$, and the couple
$(L_p, \TL{s_1}{p}{q_1})$ is $(\|\cdot\|_p, \|\cdot\|_{F^{s_1}_{p,q_1}})$.
\end{definition}

\begin{theorem}[Triebel--Lizorkin spaces interpolate to Besov spaces]
\label{thm:tl-interpolate}
\lean{Rd.memKMethod_triebelLizorkinCouple_iff, Rd.memKMethod_lpTriebelLizorkinCouple_iff,
  Rd.memKMethod_iff_memApprox_of_triebelLizorkin_equiv,
  Rd.memKMethod_iff_memApprox_of_triebelLizorkin_equiv_left}
\leanok
\uses{def:tl-couples, thm:tl-besov-sandwich, lem:tl-sandwich-constants, def:lq-norm,
  cor:approximation-space-sandwich, cor:comparable-couples, def:couple-k-method}
Let $1 \leq p < \infty$, $0 < \eta < 1$, $q_0, q_1, q \in [1, \infty]$ and $f \in L_p$.
\begin{enumerate}
  \item For every $0 < s_0 < s_1$, $f \in (\TL{s_0}{p}{q_0}, \TL{s_1}{p}{q_1})_{\eta, q}$ if and
    only if $f \in \Besov{(1 - \eta) s_0 + \eta s_1}{p}{q}$.
  \item For every $s_1 > 0$, $f \in (L_p, \TL{s_1}{p}{q_1})_{\eta, q}$ if and only if
    $f \in \Besov{\eta s_1}{p}{q}$.
  \item For every $0 < s_0 < s_1$ and all $N_0, N_1 : L_p \to [0, \infty]$ and
    $c_0, c_0', c_1, c_1' \in [0, \infty)$ with $N_i \leq c_i \|\cdot\|_{F^{s_i}_{p,q_i}}$ and
    $\|\cdot\|_{F^{s_i}_{p,q_i}} \leq c_i' N_i$ for $i = 0, 1$, $f$ lies in the space of the
    continuous $K$-method with $(\eta, q)$ of the couple $(N_0, N_1)$ if and only if
    $f \in \Besov{(1 - \eta) s_0 + \eta s_1}{p}{q}$.
  \item For every $s_1 > 0$ and all $N_1 : L_p \to [0, \infty]$ and $c_1, c_1' \in [0, \infty)$
    with $N_1 \leq c_1 \|\cdot\|_{F^{s_1}_{p,q_1}}$ and
    $\|\cdot\|_{F^{s_1}_{p,q_1}} \leq c_1' N_1$, $f$ lies in the space of the continuous
    $K$-method with $(\eta, q)$ of the couple $(\|\cdot\|_p, N_1)$ if and only if
    $f \in \Besov{\eta s_1}{p}{q}$.
\end{enumerate}
\end{theorem}

\begin{proof}
\leanok
By \Cref{thm:tl-besov-sandwich}, for $s > 0$ the quasi-norm of $\TL{s}{p}{q_i}$ is at most
$8 \cdot 8^s \|\Phi\|_1$ times the dyadic quasi-norm of $A^s_{\min(p,q_i)}$ and at least
$(2 (1 - 2^{-s})^{-1})^{-1}$ times that of $A^s_{\max(p,q_i)}$, with finite constants
(\Cref{lem:tl-sandwich-constants} and \Cref{def:lq-norm}). Parts (1) and (2) are the two parts of
\Cref{cor:approximation-space-sandwich}. Parts (3) and (4) follow, since equivalent couples have
the same spaces (\Cref{cor:comparable-couples}).
\end{proof}

\section{Triebel--Lizorkin spaces on the torus as interpolation spaces}
\label{sec:torus-tl-interpolation}

The Besov sandwich on the torus (\Cref{thm:torus-tl-besov-sandwich}) places the quasi-norm of
$\TL{s}{p}{q}(\T^d)$ between two Besov quasi-norms of the trigonometric polynomials, and the
sandwich principle (\Cref{cor:approximation-space-sandwich}) gives the torus form of
\Cref{thm:tl-interpolate}.

\begin{definition}[The constants of the Besov sandwich on the torus]
\label{def:torus-tl-sandwich-consts}
\lean{Torus.triebelLizorkinBesovConst, Torus.besovTriebelLizorkinConst}
\leanok
\uses{thm:torus-tl-besov-sandwich}
For every $s \in \R$,
\[
  K_{d,s} = 2 (\#d + 1) \, 2^{(\#d + 1) s} \, \|\Phi\|_1
  \quad \text{and} \quad
  K'_s = 2 \cdot 2^s (1 - 2^{-s})^{-1} ,
\]
the constants of the two bounds of \Cref{thm:torus-tl-besov-sandwich}.
\end{definition}

\begin{lemma}[The constants on the torus are finite]
\label{lem:torus-tl-sandwich-consts-finite}
\lean{Torus.triebelLizorkinBesovConst_ne_top, Torus.besovTriebelLizorkinConst_ne_top}
\leanok
\uses{def:torus-tl-sandwich-consts, lem:lp-kernels-basic}
For every $s \in \R$, $K_{d,s} < \infty$, and for every $s > 0$, $K'_s < \infty$.
\end{lemma}

\begin{proof}
\leanok
$\Phi$ is integrable (\Cref{lem:lp-kernels-basic}), and $2^{-s} < 1$ for $s > 0$.
\end{proof}

\begin{definition}[Triebel--Lizorkin couples on the torus]
\label{def:torus-tl-couples}
\lean{Torus.triebelLizorkinCouple, Torus.triebelLizorkinCouple_norm0,
  Torus.triebelLizorkinCouple_norm1, Torus.lpTriebelLizorkinCouple,
  Torus.lpTriebelLizorkinCouple_norm0, Torus.lpTriebelLizorkinCouple_norm1}
\leanok
\uses{def:couple, def:torus-triebel-lizorkin}
For every $1 \leq p < \infty$, all $s_0, s_1 \in \R$ and $q_0, q_1 \in [1, \infty]$, the couple
$(\TL{s_0}{p}{q_0}(\T^d), \TL{s_1}{p}{q_1}(\T^d))$ on $L_p(\T^d)$ is
$(\|\cdot\|_{F^{s_0}_{p,q_0}}, \|\cdot\|_{F^{s_1}_{p,q_1}})$, and the couple
$(L_p(\T^d), \TL{s_1}{p}{q_1}(\T^d))$ is $(\|\cdot\|_p, \|\cdot\|_{F^{s_1}_{p,q_1}})$.
\end{definition}

\begin{theorem}[Triebel--Lizorkin spaces on the torus interpolate to Besov spaces]
\label{thm:torus-tl-interpolate}
\lean{Torus.memKMethod_triebelLizorkinCouple_iff, Torus.memKMethod_lpTriebelLizorkinCouple_iff}
\leanok
\uses{def:torus-tl-couples, def:couple, thm:torus-tl-besov-sandwich,
  lem:torus-tl-sandwich-consts-finite, cor:approximation-space-sandwich, def:couple-k-method}
Let $1 \leq p < \infty$, $0 < \eta < 1$, $q_0, q_1, q \in [1, \infty]$ and $f \in L_p(\T^d)$.
\begin{enumerate}
  \item For every $0 < s_0 < s_1$,
    $f \in (\TL{s_0}{p}{q_0}(\T^d), \TL{s_1}{p}{q_1}(\T^d))_{\eta, q}$ if and only if
    $f \in \Besov{(1 - \eta) s_0 + \eta s_1}{p}{q}(\T^d)$.
  \item For every $s_1 > 0$, $f \in (L_p(\T^d), \TL{s_1}{p}{q_1}(\T^d))_{\eta, q}$ if and only
    if $f \in \Besov{\eta s_1}{p}{q}(\T^d)$.
\end{enumerate}
\end{theorem}

\begin{proof}
\leanok
By \Cref{thm:torus-tl-besov-sandwich}, for $s > 0$ the quasi-norm of $\TL{s}{p}{q_i}(\T^d)$ is at
most $K_{d,s}$ times the dyadic quasi-norm of $A^s_{\min(p,q_i)}$ and at least $(K'_s)^{-1}$
times that of $A^s_{\max(p,q_i)}$, for the scheme $A$ of trigonometric polynomials, with finite
constants (\Cref{lem:torus-tl-sandwich-consts-finite}). Apply
\Cref{cor:approximation-space-sandwich}.
\end{proof}

\chapter{Vector-valued singular integrals}
\label{ch:singular-integrals}

This chapter works on a finite-dimensional real inner product space $V$ of dimension $n$ with
Lebesgue measure. It proves the strong type $(p,p)$, $1 < p < \infty$, of the analysis operator
$f \mapsto (\theta_j * f)_j$ and of the synthesis operators
$(g_j)_j \mapsto \sum_{j \in s} \theta_j * g_j$, $s \subseteq \N$ finite, of a
Calder\'on--Zygmund family $(\theta_j)_j$, with explicit constants. For the weak and strong
type, $n \geq 1$. \Cref{sec:dilation-families} introduces these families and proves the
H\"ormander condition for the dilates of uniform Schwartz families;
\Cref{sec:convolution-adjoint} moves convolution kernels and their derivatives to the other side
of an integral; \Cref{sec:plancherel-bounds} gives the bounds on $L_2(V)$;
\Cref{sec:vector-valued-weak-type} proves the weak type $(1,1)$ with the signed
Calder\'on--Zygmund decomposition of \Cref{sec:calderon-zygmund-signed}; and
\Cref{sec:vector-valued-strong-type} interpolates with the theorems of
\Cref{ch:operator-interpolation}. \Cref{sec:square-function} applies the strong type to the
Littlewood--Paley blocks and proves the square function theorem $\TL{0}{p}{2} = L_p(V)$.
\Cref{sec:torus-singular-integrals} transfers the strong type to the torus by localising the
periodic lift, and \Cref{sec:torus-square-function} proves
$\TL{0}{p}{2}(\T^d) = L_p(\T^d)$. The same
families give $\Sob{m}{p}(V) = \TL{m}{p}{2}$ (\Cref{ch:sobolev-rd}),
$H^{s,p}(V) = \TL{s}{p}{2}$ (\Cref{ch:bessel-potential-spaces}) and the extension theorem for
$\Sob{1}{p}$ (\Cref{ch:sobolev-domains}).

\section{Dilation families and the H\"ormander condition}
\label{sec:dilation-families}

The Sobolev spaces $\Sob{m}{p}(V)$ and the Bessel potential spaces for $1 < p < \infty$ rest on
vector-valued singular integrals whose kernels are families $\theta_j = (G_j)_{2^j}$,
$j \in \N$, of dilated Schwartz functions $G_j$ with Schwartz seminorms bounded uniformly in
$j$. Two numbers control such a family: the Fourier overlap $B$ with
$\sum_j |\widehat{\theta_j}(\xi)|^2 \leq B$ for every $\xi$, which gives the $L_2$ bounds by
Plancherel's theorem, and the H\"ormander constant $H$ with
$\int_{\{\|x\| > 2\|y\|\}} \|(\theta_j(x - y) - \theta_j(x))_j\|_{\ell_2} \, dx \leq H$ for every
$y$, which gives the weak type $(1,1)$ through the Calder\'on--Zygmund decomposition. This
section proves the H\"ormander condition for every uniform family, with an explicit constant,
and shows that the Littlewood--Paley blocks $\psi_j$ and the synthesis kernels of the square
function theorem are such families. The H\"ormander condition for kernels with values in a
Hilbert space, here $\ell_2$, is the hypothesis of the theory of vector-valued singular
integrals of Benedek, Calder\'on and Panzone~\cite{BenedekCalderonPanzone}, which extends
H\"ormander's condition for scalar kernels~\cite{Hormander1960}; see also
Stein~\cite{SteinSingular}. Throughout, $C_{G,r,N}$, $T_{G,r,N}$ and
$I_M$ are the constants of \Cref{lem:schwartz-decay} and \Cref{lem:schwartz-tails},
$\Delta_v F(x) = F(x + v) - F(x)$, and $\delta_c F = F_c$ denotes the dilation by $c$, which
keeps it apart from the kernels $\Phi_k = \delta_{2^k} \Phi$. Inverses and minima of the
quantities $2^j \|y\|$ are taken in $[0, \infty]$, with $0^{-1} = \infty$ and
$\infty^{-1} = 0$; so at $y = 0$ the bounds below read $0 \leq 0$ or $0 \leq \infty$.

\begin{definition}[Dilation, reflection and conjugation of Schwartz functions]
\label{def:schwartz-dilate}
\lean{Rd.schwartzDilate, Rd.coe_schwartzDilate, Rd.schwartzReflect, Rd.coe_schwartzReflect,
  Rd.schwartzConj, Rd.coe_schwartzConj}
\leanok
\uses{lem:dilate}
For every Schwartz function $G$ on $V$ and every real $c \neq 0$, the functions
$G_c(x) = c^n G(c x)$, $x \mapsto G(-x)$ and $x \mapsto \overline{G(x)}$ are Schwartz functions.
\end{definition}

The first two are compositions of $G$ with linear equivalences of $V$ (Mathlib's
\texttt{SchwartzMap.\allowbreak compCLMOfContinuousLinearEquiv}), the third the composition of
the real-linear conjugation with $G$ (Mathlib's \texttt{SchwartzMap.\allowbreak postcompCLM}).

\begin{lemma}[Derivatives of dilations]
\label{lem:derivative-dilate}
\lean{Rd.schwartzDilate_sum, Rd.lineDerivOp_schwartzDilate,
  Rd.iteratedLineDerivOp_schwartzDilate}
\leanok
\uses{def:schwartz-dilate}
Let $c \neq 0$ be real. For every finite family $(G_i)$ of Schwartz functions,
$(\sum_i G_i)_c = \sum_i (G_i)_c$. For every Schwartz function $G$, every $v \in V$, every
$k \in \N$ and all $w \in V^k$,
\[
  \partial_v (G_c) = c \, (\partial_v G)_c , \qquad \partial^w (G_c) = c^k (\partial^w G)_c .
\]
\end{lemma}

\begin{proof}
\leanok
$G_c$ is $c^n$ times the composition of $G$ with the linear map $x \mapsto c x$, so the chain rule
gives $\partial_v (G_c) = c^n (\partial_{c v} G)(c \, \cdot) = c \, (\partial_v G)_c$ (Mathlib's
\texttt{SchwartzMap.\allowbreak lineDerivOp\_\allowbreak compCLMOfContinuousLinearEquiv}).
Induction on $k$ gives the second identity; the first holds since $G \mapsto G_c$ is linear.
\end{proof}

\begin{lemma}[Dilations are linear]
\label{lem:dilate-sub-l1}
\lean{Rd.dilate_sub}
\leanok
\uses{lem:dilate}
For all functions $F, G : V \to \C$ and every real $c$, $(F - G)_c = F_c - G_c$.
\end{lemma}

\begin{proof}
\leanok
The identity holds pointwise.
\end{proof}

\begin{definition}[Uniform Schwartz families]
\label{def:uniform-schwartz}
\lean{Rd.IsUniformSchwartz, Rd.familyDecayConst, Rd.familyDecayConst_ne_top,
  Rd.familyTailConst, Rd.familyTailConst_ne_top}
\leanok
\uses{lem:schwartz-decay, lem:schwartz-tails}
A family $(G_j)_{j \in \N}$ of Schwartz functions is uniformly Schwartz with bounds
$A = (A_{k,l})_{k, l \in \N}$ if $\sup_x \|x\|^k \|D^l G_j(x)\| \leq A_{k,l}$ for all
$j, k, l \in \N$. For all $r, N \in \N$ let $C_{A,r,N} = 2^N (A_{0,r} + A_{N,r})$ and
$T_{A,r,N} = C_{A,0,N+r+n+1} \, I_{n+1}$; both are finite.
\end{definition}

\begin{lemma}[Constants of uniform families]
\label{lem:uniform-schwartz-constants}
\lean{Rd.IsUniformSchwartz.const, Rd.IsUniformSchwartz.of_forall_mem,
  Rd.IsUniformSchwartz.of_forall_eq_or_eq, Rd.schwartzDecayConst_le_familyDecayConst,
  Rd.tailConst_le_familyTailConst}
\leanok
\uses{def:uniform-schwartz}
For every uniform family $(G_j)$ with bounds $A$ and all $j, r, N \in \N$,
$C_{G_j,r,N} \leq C_{A,r,N}$ and $T_{G_j,r,N} \leq T_{A,r,N}$. For every Schwartz function $G$,
the constant family $(G)_j$ is uniformly Schwartz with bounds
$A_{k,l} = \sup_x \|x\|^k \|D^l G(x)\|$. For every finite nonempty set $S$
of Schwartz functions and every family $(G_j)$ with $G_j \in S$ for every $j$, $(G_j)$ is
uniformly Schwartz with bounds $A_{k,l} = \max_{g \in S} \sup_x \|x\|^k \|D^l g(x)\|$; in
particular for $S = \{G^0, G^1\}$ with
$A_{k,l} = \max(\sup_x \|x\|^k \|D^l G^0(x)\|, \sup_x \|x\|^k \|D^l G^1(x)\|)$.
\end{lemma}

\begin{proof}
\leanok
The constants increase with the seminorms.
\end{proof}

\begin{lemma}[Bounds of uniform families]
\label{lem:uniform-schwartz-exceptions}
\lean{Rd.IsUniformSchwartz.nonneg, Rd.IsUniformSchwartz.of_forall_ge_eq}
\leanok
\uses{def:uniform-schwartz}
For every uniform family $(G_j)$ with bounds $A$, $A_{k,l} \geq 0$ for all $k, l \in \N$. For
every $N \in \N$ and every family $(G'_j)$ of Schwartz functions with $G'_j = G_j$ for every
$j \geq N$, $(G'_j)$ is uniformly Schwartz with bounds
$\max \bigl( A_{k,l}, \sum_{i < N} \sup_x \|x\|^k \|D^l G'_i(x)\| \bigr)$.
\end{lemma}

\begin{proof}
\leanok
$A_{k,l}$ dominates the seminorm $\sup_x \|x\|^k \|D^l G_0(x)\| \geq 0$. For $j \geq N$,
$G'_j = G_j$ has seminorms at most $A_{k,l}$, and for $j < N$ the seminorm of $G'_j$ is a term of
the sum.
\end{proof}

\begin{definition}[Gradient mass]
\label{def:gradient-mass}
\lean{Rd.gradMassConst, Rd.gradMassConst_ne_top, Rd.familyGradMass, Rd.familyGradMass_ne_top}
\leanok
\uses{def:uniform-schwartz}
For every Schwartz function $G$ let $M_G = C_{G,1,n+1} \, I_{n+1}$, and for every family of
bounds $A$ let $M_A = C_{A,1,n+1} \, I_{n+1}$; both are finite.
\end{definition}

\begin{lemma}[$L_1$ norms of directional derivatives]
\label{lem:gradient-mass}
\lean{Rd.enorm_iteratedFDeriv_schwartz_le, Rd.lintegral_enorm_lineDerivOp_le,
  Rd.gradMassConst_le_familyGradMass}
\leanok
\uses{def:gradient-mass, lem:schwartz-decay, lem:schwartz-differences,
  lem:uniform-schwartz-constants}
For every Schwartz function $G$, all $r, N \in \N$ and all $x, v \in V$,
$\|D^r G(x)\| \leq C_{G,r,N} (1 + \|x\|)^{-N}$ and $\int |\partial_v G| \leq \|v\| \, M_G$. For
every uniform family $(G_j)$ with bounds $A$ and every $j \in \N$, $M_{G_j} \leq M_A$.
\end{lemma}

\begin{proof}
\leanok
The first bound is $(1 + \|x\|)^N \|D^r G(x)\| \leq C_{G,r,N}$ (\Cref{lem:schwartz-decay}).
With $|\partial_v G(x)| \leq \|v\| \, \|DG(x)\|$ (\Cref{lem:schwartz-differences}) and $r = 1$,
$N = n + 1$, integration gives the second. The third follows from
\Cref{lem:uniform-schwartz-constants}.
\end{proof}

\begin{lemma}[Short steps]
\label{lem:hormander-short}
\lean{Rd.lintegral_enorm_dilate_sub_le_mul}
\leanok
\uses{lem:gradient-mass, lem:dilate, lem:rd-differences, lem:schwartz-differences-lp}
For every Schwartz function $G$, every $j \in \N$ and every $y \in V$,
\[
  \int_V |G_{2^j}(x - y) - G_{2^j}(x)| \, dx \leq M_G \, 2^j \|y\| .
\]
\end{lemma}

\begin{proof}
\leanok
$x \mapsto G_{2^j}(x - y) - G_{2^j}(x)$ is $\Delta_{-y} G_{2^j} = (\Delta_{-2^j y} G)_{2^j}$
(\Cref{lem:rd-differences}), whose $L_1$ norm is $\|\Delta_{-2^j y} G\|_1$
(\Cref{lem:dilate}). A difference of a Schwartz function is at most its directional
derivative in $L_1$ (\Cref{lem:schwartz-differences-lp}), and
$\|\partial_{-2^j y} G\|_1 \leq 2^j \|y\| \, M_G$ (\Cref{lem:gradient-mass}).
\end{proof}

\begin{lemma}[Long steps]
\label{lem:hormander-long}
\lean{Rd.inv_one_add_ofReal_two_pow_mul_le, Rd.setLIntegral_enorm_dilate_sub_le_inv}
\leanok
\uses{lem:schwartz-tails}
For every Schwartz function $G$, every $j \in \N$ and every $y \in V$,
\[
  \int_{\{\|x\| > 2\|y\|\}} |G_{2^j}(x - y) - G_{2^j}(x)| \, dx \leq
  2 \, T_{G,0,1} \, (2^j \|y\|)^{-1} .
\]
\end{lemma}

\begin{proof}
\leanok
For $\|x\| > 2 \|y\|$ both $\|x\|$ and $\|x - y\| \geq \|x\| - \|y\|$ are at least $\|y\|$, so
after the substitution $z = x - y$ each of the two terms is at most
$\int_{\{\|z\| \geq \|y\|\}} |G_{2^j}(z)| \, dz \leq T_{G,0,1} (1 + 2^j \|y\|)^{-1}$
(\Cref{lem:schwartz-tails}), and $(1 + t)^{-1} \leq t^{-1}$.
\end{proof}

\begin{definition}[H\"ormander constants]
\label{def:hormander-const}
\lean{Rd.hormanderConst, Rd.hormanderConst_ne_top, Rd.familyHormanderConst,
  Rd.familyHormanderConst_ne_top}
\leanok
\uses{def:gradient-mass, def:uniform-schwartz}
For every Schwartz function $G$ let $H_G = M_G + 2 \, T_{G,0,1}$, and for every family of bounds
$A$ let $H_A = M_A + 2 \, T_{A,0,1}$; both are finite.
\end{definition}

\begin{proposition}[The H\"ormander bound for one dilate]
\label{prop:hormander-one-dilate}
\lean{Rd.setLIntegral_enorm_dilate_sub_le_min, Rd.hormanderConst_le_familyHormanderConst}
\leanok
\uses{def:hormander-const, lem:hormander-short, lem:hormander-long, lem:gradient-mass,
  lem:uniform-schwartz-constants}
For every Schwartz function $G$, every $j \in \N$ and every $y \in V$,
\[
  \int_{\{\|x\| > 2\|y\|\}} |G_{2^j}(x - y) - G_{2^j}(x)| \, dx \leq
  H_G \min\bigl(2^j \|y\|, (2^j \|y\|)^{-1}\bigr) .
\]
For every uniform family $(G_j)$ with bounds $A$ and every $j \in \N$, $H_{G_j} \leq H_A$.
\end{proposition}

\begin{proof}
\leanok
For $2^j \|y\| \leq 1$ the minimum is $2^j \|y\|$, and \Cref{lem:hormander-short} applies on all
of $V$; otherwise it is $(2^j \|y\|)^{-1}$, and \Cref{lem:hormander-long} applies. The
comparison of the constants follows from \Cref{lem:gradient-mass} and
\Cref{lem:uniform-schwartz-constants}.
\end{proof}

\begin{theorem}[The H\"ormander condition for dilation families]
\label{thm:hormander-dilation-family}
\lean{Rd.lintegral_lqNorm_sub_le}
\leanok
\uses{prop:hormander-one-dilate, lem:dyadic-min-sum, def:lq-norm, lem:lq-basic}
For every uniform family $(G_j)$ with bounds $A$, $\theta_j = (G_j)_{2^j}$, and every
$y \in V$,
\[
  \int_{\{\|x\| > 2\|y\|\}} \|(\theta_j(x - y) - \theta_j(x))_j\|_{\ell_2} \, dx \leq 4 H_A .
\]
\end{theorem}

\begin{proof}
\leanok
The $\ell_2$ norm is at most the $\ell_1$ norm (\Cref{lem:lq-basic}). Exchanging sum and integral,
the $j$-th integral is at most $H_A \min(2^j \|y\|, (2^j \|y\|)^{-1})$
(\Cref{prop:hormander-one-dilate}), and the sum of the minima is at most $4$
(\Cref{lem:dyadic-min-sum}).
\end{proof}

\begin{definition}[Calder\'on--Zygmund families]
\label{def:cz-family}
\lean{Rd.IsCZFamily}
\leanok
\uses{def:lq-norm}
For all $B, H \in [0, \infty]$, a family $(\theta_j)_{j \in \N}$ of functions $V \to \C$ is a
Calder\'on--Zygmund family with Fourier overlap $B$ and H\"ormander constant $H$ if every
$\theta_j$ is integrable, continuous and bounded with integrable Fourier transform, if
$\sum_j |\widehat{\theta_j}(\xi)|^2 \leq B$ for every $\xi \in V$, and if
\[
  \int_{\{\|x\| > 2\|y\|\}} \|(\theta_j(x - y) - \theta_j(x))_j\|_{\ell_2} \, dx \leq H
\]
for every $y \in V$.
\end{definition}

\begin{lemma}[Larger constants]
\label{lem:cz-family-mono}
\lean{Rd.IsCZFamily.mono}
\leanok
\uses{def:cz-family}
For all $B \leq B'$ and $H \leq H'$ in $[0, \infty]$, every Calder\'on--Zygmund family with
overlap $B$ and H\"ormander constant $H$ is a Calder\'on--Zygmund family with overlap $B'$ and
H\"ormander constant $H'$.
\end{lemma}

\begin{proof}
\leanok
Both conditions are upper bounds.
\end{proof}

\begin{lemma}[Adjoint families]
\label{lem:cz-family-adjoint}
\lean{Rd.fourier_conj_comp_neg, Rd.IsCZFamily.adjoint}
\leanok
\uses{def:cz-family, lem:fourier-conjugation}
For every function $\theta : V \to \C$ and every $\xi \in V$, the Fourier transform of
$x \mapsto \overline{\theta(-x)}$ at $\xi$ is $\overline{\widehat\theta(\xi)}$. For every
Calder\'on--Zygmund family $(\theta_j)$ with overlap $B$ and H\"ormander constant $H$, the
family $\theta_j^*(x) = \overline{\theta_j(-x)}$ is a Calder\'on--Zygmund family with overlap
$B$ and H\"ormander constant $H$.
\end{lemma}

\begin{proof}
\leanok
$\overline{\widehat\theta} = \mathcal F^{-1} \overline\theta$ (\Cref{lem:fourier-conjugation}),
and $\mathcal F^{-1} f = \mathcal F (f(-\cdot))$ (Mathlib's
\texttt{Real.\allowbreak fourierInv\_\allowbreak eq\_\allowbreak fourier\_\allowbreak
comp\_\allowbreak neg}); so the overlap is unchanged. Integrability, continuity and boundedness
pass through $x \mapsto -x$ and conjugation. In the H\"ormander integral,
$|\theta_j^*(x - y) - \theta_j^*(x)| = |\theta_j(-x + y) - \theta_j(-x)|$; the substitution
$x \mapsto -x$ keeps the region $\{\|x\| > 2\|y\|\}$ and turns the integral into the H\"ormander
integral of $(\theta_j)$ at $-y$.
\end{proof}

\begin{theorem}[Dilation families are Calder\'on--Zygmund families]
\label{thm:cz-family-dilation}
\lean{Rd.IsCZFamily.of_dilation}
\leanok
\uses{def:cz-family, thm:hormander-dilation-family, lem:dilate, lem:fourier-dilate,
  lem:lp-kernels-basic}
For every uniform family $(G_j)$ with bounds $A$ and every $B \in [0, \infty]$ such that
$\theta_j = (G_j)_{2^j}$ satisfies $\sum_j |\widehat{\theta_j}(\xi)|^2 \leq B$ for every
$\xi \in V$, $(\theta_j)$ is a Calder\'on--Zygmund family with overlap $B$ and H\"ormander
constant $4 H_A$.
\end{theorem}

\begin{proof}
\leanok
Schwartz functions are bounded (\Cref{lem:lp-kernels-basic}); their dilations are integrable,
continuous and bounded (\Cref{lem:dilate}), and the Fourier transforms of the dilations are
dilations of Schwartz functions (\Cref{lem:fourier-dilate}). The
H\"ormander condition is \Cref{thm:hormander-dilation-family}.
\end{proof}

\begin{lemma}[Calder\'on--Zygmund families with exceptional members]
\label{lem:cz-family-exceptions}
\lean{Rd.enorm_fourier_dilate_le, Rd.IsCZFamily.of_dilation_of_forall_ge_eq}
\leanok
\uses{lem:uniform-schwartz-exceptions, thm:cz-family-dilation, lem:fourier-dilate,
  def:cz-family, def:hormander-const, lem:tsum-changed-terms}
For every Schwartz function
$G$, every real $c > 0$ and every $\xi \in V$,
$|\widehat{G_c}(\xi)| \leq \|\widehat G\|_{0,0}$. Let $(G_j)$ be a uniform family of Schwartz
functions with bounds $A$ such that $\sum_j |\widehat{(G_j)_{2^j}}(\xi)|^2 \leq B$ for every
$\xi \in V$, let $N \in \N$, and let $(G'_j)$ be Schwartz functions with $G'_j = G_j$ for every
$j \geq N$. Then $((G'_j)_{2^j})_j$ is a Calder\'on--Zygmund family with overlap
$B + \sum_{i < N} \|\widehat{G'_i}\|_{0,0}^2$ and H\"ormander constant $4 H_{A'}$, where
$A'_{k,l} = \max\bigl(A_{k,l}, \sum_{i < N} \|G'_i\|_{k,l}\bigr)$.
\end{lemma}

\begin{proof}
\leanok
First, $\widehat{G_c}(\xi) = \widehat G(c^{-1} \xi)$ (\Cref{lem:fourier-dilate}). The family
$(G'_j)$ is uniformly Schwartz with bounds $A'$ (\Cref{lem:uniform-schwartz-exceptions}); the
first claim and \Cref{lem:tsum-changed-terms} bound its overlap, and
\Cref{thm:cz-family-dilation} gives the family.
\end{proof}

\paragraph{The Littlewood--Paley families.}

\begin{lemma}[Values of the dilated cut-offs]
\label{lem:fourier-phi-values}
\lean{Rd.fourier_dilatedLpKernel_eq_one, Rd.fourier_dilatedLpKernel_eq_zero,
  Rd.exists_fourier_dilatedLpKernel_eq_ofReal, Rd.norm_fourier_dilatedLpKernel_le_one,
  Rd.norm_fourier_dilatedLpKernel_sub_le_one}
\leanok
\uses{def:lp-kernels, lem:lp-kernels-frequency, lem:lp-kernels-basic}
For all $k, l \in \N$ and every $\xi \in V$: $\widehat{\Phi_k}(\xi) = 1$ if
$\|\xi\| \leq 2^k$, $\widehat{\Phi_k}(\xi) = 0$ if $\|\xi\| \geq 2^{k+1}$,
$\widehat{\Phi_k}(\xi)$ is a real number in $[0, 1]$, $|\widehat{\Phi_k}(\xi)| \leq 1$ and
$|\widehat{\Phi_k}(\xi) - \widehat{\Phi_l}(\xi)| \leq 1$.
\end{lemma}

\begin{proof}
\leanok
$\widehat{\Phi_k}(\xi) = \chi_{1,2}(2^{-k} \xi)$ (\Cref{lem:lp-kernels-frequency}), and
the smooth bump $\chi_{1,2}$ has values in $[0, 1]$, equals $1$ on the closed unit ball and
vanishes outside the ball of radius $2$ (\Cref{lem:lp-kernels-basic}).
\end{proof}

\begin{definition}[The blocks as dilates]
\label{def:psi-base}
\lean{Rd.psiBase, Rd.psiBaseBound, Rd.psiBaseBound_nonneg}
\leanok
\uses{def:lp-kernels, def:schwartz-dilate}
Let $G^{\psi}_0 = \Phi$ and $G^{\psi}_j = \Phi - \delta_{1/2} \Phi$ for $j \geq 1$, and for
all $k, l \in \N$ let $A^{\psi}_{k,l} \geq 0$ be the larger of the seminorms
$\sup_x \|x\|^k \|D^l \Phi(x)\|$ and $\sup_x \|x\|^k \|D^l (\Phi - \delta_{1/2} \Phi)(x)\|$.
\end{definition}

\begin{lemma}[The blocks form a uniform family]
\label{lem:psi-base}
\lean{Rd.dilate_psiBase, Rd.psiS, Rd.coe_psiS, Rd.isUniformSchwartz_psiBase}
\leanok
\uses{def:psi-base, lem:dilate-sub-l1, def:moment-kernel, lem:uniform-schwartz-constants}
For every $j \in \N$, $\psi_j = (G^{\psi}_j)_{2^j}$, a Schwartz function, and $(G^{\psi}_j)$ is
uniformly Schwartz with bounds $A^{\psi}$.
\end{lemma}

\begin{proof}
\leanok
$\delta_c \delta_d = \delta_{cd}$ (\Cref{def:moment-kernel}), so
$\delta_{2^{j+1}} \delta_{1/2} \Phi = \Phi_j$ and
$\delta_{2^{j+1}} (\Phi - \delta_{1/2} \Phi) = \Phi_{j+1} - \Phi_j = \psi_{j+1}$
(\Cref{lem:dilate-sub-l1}). The family takes the two values $\Phi$ and
$\Phi - \delta_{1/2} \Phi$ (\Cref{lem:uniform-schwartz-constants}).
\end{proof}

\begin{theorem}[The blocks form a Calder\'on--Zygmund family]
\label{thm:cz-family-psi}
\lean{Rd.tsum_fourier_lpBlockKernel_sq_le, Rd.isCZFamily_lpBlockKernel}
\leanok
\uses{lem:psi-base, prop:bessel-lp-weights, thm:cz-family-dilation}
For every $\xi \in V$, $\sum_j |\widehat{\psi_j}(\xi)|^2 \leq 2$, and $(\psi_j)$ is a
Calder\'on--Zygmund family with overlap $2$ and H\"ormander constant $4 H_{A^{\psi}}$.
\end{theorem}

\begin{proof}
\leanok
The overlap bound is the case $s = 0$ of \Cref{prop:bessel-lp-weights}: at most two blocks are
nonzero at $\xi$, each of modulus at most $1$. \Cref{thm:cz-family-dilation} and
\Cref{lem:psi-base} give the rest.
\end{proof}

\begin{definition}[Synthesis kernels]
\label{def:psi-dual}
\lean{Rd.psiDual, Rd.psiDualBase, Rd.psiDualBaseBound, Rd.psiDualBaseBound_nonneg}
\leanok
\uses{def:lp-kernels, def:schwartz-dilate}
Let $\tilde\psi_0 = \Phi_1$, $\tilde\psi_1 = \Phi_2$ and
$\tilde\psi_j = \Phi_{j+1} - \Phi_{j-2}$ for $j \geq 2$. Let
$\tilde G_0 = \tilde G_1 = \delta_2 \Phi$ and $\tilde G_j = \delta_2 \Phi - \delta_{1/4} \Phi$ for
$j \geq 2$, and for all $k, l \in \N$ let $\tilde A_{k,l} \geq 0$ be the larger of the seminorms
$\sup_x \|x\|^k \|D^l \delta_2 \Phi(x)\|$ and
$\sup_x \|x\|^k \|D^l (\delta_2 \Phi - \delta_{1/4} \Phi)(x)\|$.
\end{definition}

\begin{lemma}[Properties of the synthesis kernels]
\label{lem:psi-dual-basic}
\lean{Rd.dilate_psiDualBase, Rd.isUniformSchwartz_psiDualBase, Rd.integrable_psiDual,
  Rd.continuous_psiDual, Rd.bddAbove_psiDual, Rd.integrable_fourier_psiDual,
  Rd.eLpNorm_psiDual_one_le, Rd.fourier_psiDual_add_two}
\leanok
\uses{def:psi-dual, lem:dilate-sub-l1, def:moment-kernel, lem:uniform-schwartz-constants,
  lem:dilate, lem:fourier-dilate, lem:lp-kernels-basic, lem:fourier-conv}
For every $j \in \N$, $\tilde\psi_j = (\tilde G_j)_{2^j}$; $(\tilde G_j)$ is uniformly Schwartz with
bounds $\tilde A$; $\tilde\psi_j$ is integrable, continuous and bounded with integrable Fourier
transform; $\|\tilde\psi_j\|_1 \leq 2 \|\Phi\|_1$; and
$\widehat{\tilde\psi_{j+2}} = \widehat{\Phi_{j+3}} - \widehat{\Phi_j}$.
\end{lemma}

\begin{proof}
\leanok
$\delta_{2^{j+2}} \delta_2 \Phi = \Phi_{j+3}$ and $\delta_{2^{j+2}} \delta_{1/4} \Phi = \Phi_j$
(\Cref{def:moment-kernel}: $\delta_c \delta_d = \delta_{cd}$), and $\delta_1 \delta_2 \Phi =
\Phi_1$, $\delta_2 \delta_2 \Phi = \Phi_2$. The family takes
two values (\Cref{lem:uniform-schwartz-constants}); the remaining properties hold for dilations
of Schwartz functions (\Cref{lem:dilate}, \Cref{lem:fourier-dilate}), and
$\|\Phi_k\|_1 = \|\Phi\|_1$ (\Cref{lem:lp-kernels-basic}).
\end{proof}

\begin{lemma}[The synthesis kernels reproduce the blocks]
\label{lem:psi-dual-reproduces}
\lean{Rd.fourier_psiDual_mul_lpBlockKernel, Rd.conv_psiDual_lpBlockKernel}
\leanok
\uses{lem:psi-dual-basic, lem:fourier-phi-values, lem:lp-kernels-frequency,
  lem:conv-of-fourier}
For every $j \in \N$, $\widehat{\tilde\psi_j} \, \widehat{\psi_j} = \widehat{\psi_j}$ and
$\tilde\psi_j * \psi_j = \psi_j$.
\end{lemma}

\begin{proof}
\leanok
Where $\widehat{\psi_j}(\xi) \neq 0$, $\|\xi\| < 2^{j+1}$, and for $j \geq 2$ also
$\|\xi\| > 2^{j-1}$ (\Cref{lem:lp-kernels-frequency}). Then $\widehat{\Phi_{j+1}}(\xi) = 1$,
and for $j \geq 2$ also $\widehat{\Phi_{j-2}}(\xi) = 0$ (\Cref{lem:fourier-phi-values}), so
$\widehat{\tilde\psi_j}(\xi) = 1$. The convolution identity follows by \Cref{lem:conv-of-fourier}.
\end{proof}

\begin{lemma}[Frequencies of the synthesis kernels]
\label{lem:psi-dual-frequency}
\lean{Rd.enorm_fourier_psiDual_le_one, Rd.norm_lt_of_fourier_psiDual_ne_zero,
  Rd.lt_norm_of_fourier_psiDual_ne_zero, Rd.mem_Icc_of_fourier_psiDual_ne_zero}
\leanok
\uses{lem:psi-dual-basic, lem:fourier-phi-values, lem:lp-kernels-frequency}
For every $j \in \N$ and every $\xi \in V$, $|\widehat{\tilde\psi_j}(\xi)| \leq 1$, and if
$\widehat{\tilde\psi_j}(\xi) \neq 0$, then $\|\xi\| < 2^{j+2}$, and $\|\xi\| > 2^{j-2}$ for
$j \geq 2$. For all $a, j \in \N$ and every $\xi \in V$ with $\|\xi\| < 2^{a+1}$ and
$\|\xi\| \geq 2^{i+1}$ for every $i < a$, if $\widehat{\tilde\psi_j}(\xi) \neq 0$, then
$a - 1 \leq j \leq a + 2$.
\end{lemma}

\begin{proof}
\leanok
The bound and the localisation follow from \Cref{lem:fourier-phi-values}: if both cut-offs
vanish at $\xi$, or both equal $1$, their difference vanishes. If
$\widehat{\tilde\psi_j}(\xi) \neq 0$, then $\|\xi\| < 2^{j+2}$, so $a \leq j + 1$ by the
minimality of $a$, and for $j \geq 2$, $2^{j-2} < \|\xi\| < 2^{a+1}$, so $j \leq a + 2$.
\end{proof}

\begin{theorem}[The synthesis kernels form a Calder\'on--Zygmund family]
\label{thm:cz-family-psi-dual}
\lean{Rd.tsum_fourier_psiDual_sq_le, Rd.isCZFamily_psiDual}
\leanok
\uses{lem:psi-dual-basic, lem:psi-dual-frequency, thm:cz-family-dilation}
For every $\xi \in V$, $\sum_j |\widehat{\tilde\psi_j}(\xi)|^2 \leq 4$, and $(\tilde\psi_j)$ is a
Calder\'on--Zygmund family with overlap $4$ and H\"ormander constant $4 H_{\tilde A}$.
\end{theorem}

\begin{proof}
\leanok
With $a$ the least index such that $\|\xi\| < 2^{a+1}$, at most the four kernels with
$a - 1 \leq j \leq a + 2$ are nonzero at $\xi$, each of modulus at most $1$
(\Cref{lem:psi-dual-frequency}). \Cref{thm:cz-family-dilation} and \Cref{lem:psi-dual-basic}
give the rest.
\end{proof}

\section{The adjoint of a convolution}
\label{sec:convolution-adjoint}

The strong type for $2 < p < \infty$ is proved by duality, through convolution with
reflected kernels, and the weak derivatives of
$\Sob{m}{p}(V)$ meet the Littlewood--Paley blocks through convolutions with Schwartz kernels.
This section moves a convolution kernel to the other side of an integral, moves derivatives
from a Schwartz kernel to the other Schwartz factor, and shows that weak derivatives commute
with convolution by Schwartz functions. Throughout, $K * f$ is the convolution of
\Cref{def:conv}, and for $v \in V^k$, $\partial^v$ is the iterated directional derivative of
\Cref{sec:schwartz-cutoff}.

\begin{definition}[Reflection]
\label{def:reflect}
\lean{Rd.reflect, Rd.reflect_apply}
\leanok
For every function $K : V \to \C$, the reflection of $K$ is $K^\vee(x) = K(-x)$.
\end{definition}

\begin{lemma}[Reflections and convolution integrals]
\label{lem:reflect-basic}
\lean{Rd.reflect_reflect, Rd.integrable_reflect, Rd.continuous_reflect, Rd.fourier_reflect,
  Rd.conv_apply_eq_integral_sub, Rd.conv_reflect_apply}
\leanok
\uses{def:reflect, def:conv}
For every function $K : V \to \C$, $(K^\vee)^\vee = K$, $K^\vee$ is integrable if $K$ is,
$K^\vee$ is continuous if $K$ is, and $\widehat{K^\vee}(\xi) = \widehat K(-\xi)$ for every
$\xi \in V$. For all functions $K, u : V \to \C$ and all $x, y \in V$,
\[
  (K * u)(x) = \int_V K(x - z) \, u(z) \, dz , \qquad
  (K^\vee * u)(y) = \int_V K(z - y) \, u(z) \, dz .
\]
\end{lemma}

\begin{proof}
\leanok
Lebesgue measure is invariant under $x \mapsto -x$. The Fourier transform of $K^\vee$ is the
substitution $x \mapsto -x$ in the defining integral (Mathlib's
\texttt{Real.\allowbreak fourierInv\_eq\_fourier\_comp\_neg}). The first convolution formula is
the substitution $z \mapsto x - z$ (Mathlib's
\texttt{MeasureTheory.\allowbreak convolution\_mul\_swap}), and the second is the first for
$K^\vee$.
\end{proof}

\begin{theorem}[The adjoint identity]
\label{thm:convolution-adjoint}
\lean{Rd.integrable_kernel_mul_mul, Rd.integral_conv_mul}
\leanok
\uses{lem:reflect-basic, thm:young}
For all $1 \leq p < \infty$ and $1 < q \leq \infty$ with $1/p + 1/q = 1$, every $K \in L_1(V)$,
every $f \in L_p(V)$ and every $g \in L_q(V)$, the function
$(x, y) \mapsto K(x - y) f(y) g(x)$ is integrable on $V \times V$, and
\[
  \int_V (K * f)(x) \, g(x) \, dx = \int_V f(y) \, (K^\vee * g)(y) \, dy .
\]
\end{theorem}

\begin{proof}
\leanok
For almost every $x$ the integral of $y \mapsto K(x - y) f(y)$ converges absolutely
(\Cref{thm:young}). By \Cref{lem:reflect-basic},
$\int_V |K(x - y)| \, |f(y)| \, dy = (|K| * |f|)(x)$ for every $x$, and $|K| * |f| \in L_p$ by
Young's inequality, so $x \mapsto (|K| * |f|)(x) \, |g(x)|$ is integrable by H\"older's
inequality. Together the two give the integrability on $V \times V$. Integrating first in $y$
gives the left side by \Cref{lem:reflect-basic}, integrating first in $x$ gives the right side,
and Fubini's theorem exchanges the integrations.
\end{proof}

With complex conjugation, the reflection becomes the adjoint kernel
$K^*(z) = \overline{K(-z)}$, the form in which the duality arguments for $p > 2$ use the
identity.

\begin{corollary}[The adjoint identity with conjugation]
\label{cor:convolution-adjoint-conj}
\lean{Rd.conv_reflect_conj, Rd.integral_conv_mul_conj}
\leanok
\uses{thm:convolution-adjoint, lem:reflect-basic}
For all functions $K, u : V \to \C$ and every $x \in V$,
$(K^\vee * \overline u)(x) = \overline{(K^* * u)(x)}$. For all $1 \leq p < \infty$ and
$1 < q \leq \infty$ with $1/p + 1/q = 1$, every $K \in L_1(V)$, every $f \in L_p(V)$ and every
$g \in L_q(V)$,
\[
  \int_V (K * f)(x) \, \overline{g(x)} \, dx
  = \int_V f(y) \, \overline{(K^* * g)(y)} \, dy .
\]
\end{corollary}

\begin{proof}
\leanok
By \Cref{lem:reflect-basic} both sides of the first identity are
$\int_V K(y - x) \, \overline{u(y)} \, dy$. \Cref{thm:convolution-adjoint} for $\overline g \in
L_q(V)$ and the first identity give the second.
\end{proof}

\begin{lemma}[Translates and reflections of Schwartz functions]
\label{lem:schwartz-shift-reflect}
\lean{SchwartzMap.lineDerivOp_compSubConstCLM, SchwartzMap.iteratedLineDerivOp_compSubConstCLM,
  Rd.lineDerivOp_schwartzReflect, Rd.iteratedLineDerivOp_schwartzReflect, Rd.shiftReflect,
  Rd.shiftReflect_apply, Rd.conv_eq_integral_shiftReflect, Rd.iteratedLineDerivOp_shiftReflect}
\leanok
\uses{def:schwartz-dilate, lem:reflect-basic}
For every Schwartz function $\eta$, all $x, y \in V$ and every $w \in V^k$, the functions
$z \mapsto \eta(z - y)$ and $z \mapsto \eta(x - z)$ are Schwartz functions,
\begin{gather*}
  \partial^w \bigl(\eta(\cdot - y)\bigr) = (\partial^w \eta)(\cdot - y), \qquad
  \partial^w \bigl(\eta(-\,\cdot)\bigr) = (-1)^k (\partial^w \eta)(-\,\cdot), \\
  \partial^w \bigl(\eta(x - \cdot)\bigr) = (-1)^k (\partial^w \eta)(x - \cdot),
\end{gather*}
and for every function $f : V \to \C$, $(\eta * f)(x)$ is the integral of $f$ against the
Schwartz function $\eta(x - \cdot)$:
$(\eta * f)(x) = \int_V \eta(x - z) \, f(z) \, dz$.
\end{lemma}

\begin{proof}
\leanok
The derivative of $z \mapsto \eta(z - y)$ at $z$ is $D\eta(z - y)$, and the directional
derivative of $\eta(-\,\cdot)$ in direction $v$ is $-(\partial_v \eta)(-\,\cdot)$ (Mathlib's
\texttt{SchwartzMap.\allowbreak lineDerivOp\_compCLMOfContinuousLinearEquiv}); induction on $k$
gives the iterated forms, and the third identity combines the first two. The convolution formula
is \Cref{lem:reflect-basic}.
\end{proof}

\begin{lemma}[Derivatives move from the kernel to the Schwartz factor]
\label{lem:conv-derivative-transfer}
\lean{Rd.conv_reflect_lineDerivOp, Rd.integral_conv_lineDerivOp_mul,
  Rd.integral_conv_iteratedLineDerivOp_mul}
\leanok
\uses{thm:convolution-adjoint, lem:schwartz-shift-reflect, lem:directional-derivatives-commute}
For all Schwartz functions $\eta$, $u$ and all $v, y \in V$,
$((\partial_v \eta)^\vee * u)(y) = -(\eta^\vee * \partial_v u)(y)$. For every
$1 \leq p < \infty$, every $f \in L_p(V)$, all Schwartz functions $\eta$, $u$ and every
$w \in V^k$,
\[
  \int_V \bigl((\partial^w \eta) * f\bigr)(x) \, u(x) \, dx
  = (-1)^k \int_V (\eta * f)(x) \, \partial^w u(x) \, dx .
\]
\end{lemma}

\begin{proof}
\leanok
By \Cref{lem:reflect-basic},
$((\partial_v \eta)^\vee * u)(y) = \int_V \partial_v \eta(z - y) \, u(z) \, dz$, and
integration by parts for Schwartz functions (Mathlib's
\texttt{SchwartzMap.\allowbreak integral\_mul\_lineDerivOp\_right\_eq\_neg\_left}), applied to
$\eta(\cdot - y)$ and $u$, gives the first identity by \Cref{lem:schwartz-shift-reflect}. For
$k = 1$, apply \Cref{thm:convolution-adjoint} with the conjugate exponent $q$ of $p$ to both
sides; Schwartz functions lie in every $L_q$, and the two reflected convolutions differ by the
sign. Induction on $k$ gives the general case, with
$\partial_a \partial^w u = \partial^w \partial_a u$
(\Cref{lem:directional-derivatives-commute}).
\end{proof}

\begin{theorem}[Weak derivatives commute with convolution]
\label{thm:weak-derivative-convolution}
\lean{Rd.conv_eq_conv_iteratedLineDerivOp, Rd.conv_eq_conv_iteratedLineDerivOp_of_memLp,
  Rd.hasWeakLineDeriv_conv, Rd.hasWeakLineDeriv_toLp_conv}
\leanok
\uses{thm:weak-derivative-tempered, lem:distribution-transpose, lem:schwartz-shift-reflect,
  lem:conv-derivative-transfer, lem:weak-derivative-integral, lem:weak-derivative-basic,
  def:test-schwartz, thm:young}
For all $1 \leq p, r \leq \infty$, every $f \in L_p(V)$, every $g \in L_r(V)$ and every
$v \in V^k$ such that $g$ is the weak derivative $\partial^v f$ on $V$, every Schwartz
function $\eta$ and every $x \in V$,
\[
  (\eta * g)(x) = \bigl((\partial^v \eta) * f\bigr)(x) .
\]
For every $1 \leq p < \infty$, every $f \in L_p(V)$, every Schwartz function $\eta$ and every
$v \in V^k$, $(\partial^v \eta) * f$ is the weak derivative $\partial^v (\eta * f)$ on $V$.
\end{theorem}

\begin{proof}
\leanok
Let $S_f$ and $S_g$ be the tempered distributions of $f$ and $g$, and for every Schwartz
function $\zeta$ let $\rho_x \zeta = \zeta(x - \cdot)$ be its reflected translate.
By \Cref{lem:schwartz-shift-reflect}, $(\eta * g)(x) = S_g(\rho_x \eta)$, which is
$\partial^v S_f(\rho_x \eta) = S_f\bigl((-1)^k \partial^v \rho_x \eta\bigr)$ by
\Cref{thm:weak-derivative-tempered} and \Cref{lem:distribution-transpose}. Again by
\Cref{lem:schwartz-shift-reflect}, $(-1)^k \partial^v \rho_x \eta = \rho_x(\partial^v \eta)$,
and $S_f\bigl(\rho_x(\partial^v \eta)\bigr) = ((\partial^v \eta) * f)(x)$.

For the second statement, $\eta * f$ and $(\partial^v \eta) * f$ lie in $L_p$ by
\Cref{thm:young}, so they are locally integrable. For every test function $\varphi$ on $V$,
\Cref{lem:conv-derivative-transfer} with the Schwartz function $u = \varphi$
(\Cref{def:test-schwartz}) is the identity of \Cref{lem:weak-derivative-integral}.
\end{proof}

\begin{lemma}[Weak derivatives of convolutions in $L_p$]
\label{lem:weak-derivative-convolution-lp}
\lean{Rd.hasWeakLineDeriv_convLp}
\leanok
\uses{thm:weak-derivative-convolution, lem:conv-op}
For all $1 \leq p, r < \infty$, every $k$, all directions $v = (v_1, \dots, v_k) \in V^k$, every
Schwartz function $\eta$, every $f \in L_p(V)$ and every $g \in L_r(V)$ with $g = \partial^v f$
weakly on $V$, $\eta * g = \partial^v (\eta * f)$ weakly on $V$.
\end{lemma}

\begin{proof}
\leanok
The weak derivative of $\eta * f$ is $(\partial^v \eta) * f$, and $(\partial^v \eta) * f =
\eta * g$ at every point (\Cref{thm:weak-derivative-convolution}).
\end{proof}

\begin{lemma}[Weak derivatives and convolution on open sets]
\label{lem:weak-derivative-convolution-local}
\lean{Rd.coe_shiftReflect_eq_comp, Rd.hasCompactSupport_shiftReflect,
  Rd.tsupport_shiftReflect_subset, Rd.conv_eq_conv_iteratedLineDerivOp_of_hasWeakLineDeriv}
\leanok
\uses{lem:schwartz-shift-reflect, lem:weak-derivative-complex-test}
Let $\Omega \subseteq V$ be open, $v \in V^k$, $g = \partial^v f$ weakly on $\Omega$, $\eta$ a
Schwartz function with compact support and $x \in V$ with $x - y \in \Omega$ for every $y$ in the
support of $\eta$. Then $(\eta * g)(x) = ((\partial^v \eta) * f)(x)$.
\end{lemma}

\begin{proof}
\leanok
The reflected translate $\rho_x \eta = \eta(x - \cdot)$ has compact support in $\Omega$, so
\Cref{lem:weak-derivative-complex-test} and \Cref{lem:schwartz-shift-reflect} give
$(\eta * g)(x) = \int \rho_x \eta \, g = (-1)^k \int (\partial^v \rho_x \eta) \, f
= \int \rho_x(\partial^v \eta) \, f = ((\partial^v \eta) * f)(x)$.
\end{proof}

\section{Plancherel bounds}
\label{sec:plancherel-bounds}

A family $(\theta_j)_{j \in \N}$ of integrable kernels on $V$ defines the analysis operator
$f \mapsto (\theta_j * f)_j$ and, for every finite set $s \subseteq \N$, the synthesis operator
$(g_j)_j \mapsto \sum_{j \in s} \theta_j * g_j$. On $L_2(V)$ both are bounded by the Fourier
overlap alone: Plancherel's theorem and the convolution theorem on $L_2$ turn them into
multiplications by the functions $\widehat{\theta_j}$. Throughout, $T_K f = K * f$ is the
convolution operator of \Cref{lem:conv-op} on $L_2(V)$, and $\widehat f$ the Fourier transform
of $f \in L_2(V)$ of \Cref{sec:sobolev-two-littlewood-paley}.

\begin{lemma}[Cauchy--Schwarz and squares of $L_2$ norms]
\label{lem:l2-squares}
\lean{enorm_sum_mul_sq_le, Rd.lintegral_enorm_conv_sq_eq}
\leanok
\uses{lem:plancherel-l2, lem:conv-op}
For every seminormed ring $R$, such as $\C$, every finite set $s$ and all $a_j, b_j \in R$,
$j \in s$,
\[
  \Bigl\|\sum_{j \in s} a_j b_j\Bigr\|^2 \leq \Bigl(\sum_{j \in s} \|a_j\|^2\Bigr)
  \Bigl(\sum_{j \in s} \|b_j\|^2\Bigr).
\]
For every integrable $K$ on $V$ and every $f \in L_2(V)$,
$\int_V |(K * f)(x)|^2 \, dx = \|T_K f\|_2^2$.
\end{lemma}

\begin{proof}
\leanok
The inequality is the Cauchy--Schwarz inequality for the norms (Mathlib's
\texttt{Finset.\allowbreak sum\_mul\_sq\_le\_sq\_mul\_sq}) after the triangle inequality and
$\|a b\| \leq \|a\| \, \|b\|$. The identity follows from \Cref{lem:plancherel-l2}, since
$T_K f = K * f$ almost everywhere (\Cref{lem:conv-op}).
\end{proof}

\begin{theorem}[The analysis bound on $L_2$]
\label{thm:l2-analysis-bound}
\lean{Rd.tsum_enorm_convLp_sq_le, Rd.lintegral_lqNorm_conv_sq_le}
\leanok
\uses{thm:convolution-theorem-l2, lem:plancherel-l2, lem:l2-squares, lem:fourier-integrable,
  def:lq-norm}
For every $B \in [0, \infty]$, every family $(\theta_j)_{j \in \N}$ of integrable kernels on $V$
with $\sum_j |\widehat{\theta_j}(\xi)|^2 \leq B$ for almost every $\xi \in V$, and every
$f \in L_2(V)$,
\[
  \sum_j \|\theta_j * f\|_2^2 \leq B \, \|f\|_2^2 , \qquad
  \int_V \|((\theta_j * f)(x))_j\|_{\ell_2}^2 \, dx \leq B \int_V |f(x)|^2 \, dx .
\]
\end{theorem}

\begin{proof}
\leanok
By \Cref{thm:convolution-theorem-l2} and monotone convergence,
\[
  \sum_j \|\theta_j * f\|_2^2
  = \int_V \Bigl(\sum_j |\widehat{\theta_j}(\xi)|^2\Bigr) |\widehat f(\xi)|^2 \, d\xi
  \leq B \int_V |\widehat f(\xi)|^2 \, d\xi = B \, \|f\|_2^2
\]
by \Cref{lem:plancherel-l2}; the integrand is measurable since every $\widehat{\theta_j}$ lies
in $L_\infty(V)$ (\Cref{lem:fourier-integrable}). The second inequality is the first for the
class of $f$, written with $\|a\|_{\ell_2}^2 = \sum_j a_j^2$ (\Cref{def:lq-norm}),
\Cref{lem:plancherel-l2}, \Cref{lem:l2-squares} and monotone convergence.
\end{proof}

\begin{lemma}[The Fourier transform of a sum of convolutions]
\label{lem:fourier-sum-conv}
\lean{Rd.fourier_sum_convLp}
\leanok
\uses{thm:convolution-theorem-l2}
For every family $(\theta_j)_{j \in \N}$ of integrable kernels on $V$, all $g_j \in L_2(V)$ and
every finite set $s \subseteq \N$,
$\mathcal F\bigl(\sum_{j \in s} \theta_j * g_j\bigr)
= \sum_{j \in s} \widehat{\theta_j} \, \widehat{g_j}$ almost everywhere.
\end{lemma}

\begin{proof}
\leanok
The Fourier transform on $L_2(V)$ is linear (Mathlib's
\texttt{Fourier\allowbreak Transform.\allowbreak fourier\_\allowbreak sum}); apply
\Cref{thm:convolution-theorem-l2} to each term.
\end{proof}

\begin{theorem}[The synthesis bound on $L_2$]
\label{thm:l2-synthesis-bound}
\lean{Rd.enorm_sum_convLp_sq_le, Rd.lintegral_enorm_sum_conv_sq_le}
\leanok
\uses{lem:fourier-sum-conv, lem:l2-squares, lem:plancherel-l2, def:lq-norm}
For every $B \in [0, \infty]$, every family $(\theta_j)_{j \in \N}$ of integrable kernels on $V$
with $\sum_j |\widehat{\theta_j}(\xi)|^2 \leq B$ for almost every $\xi \in V$, all
$g_j \in L_2(V)$ and every finite set $s \subseteq \N$,
\[
  \Bigl\|\sum_{j \in s} \theta_j * g_j\Bigr\|_2^2 \leq B \sum_{j \in s} \|g_j\|_2^2
\]
and
\[
  \int_V \Bigl|\sum_{j \in s} (\theta_j * g_j)(x)\Bigr|^2 \, dx
  \leq B \int_V \|(g_j(x))_j\|_{\ell_2}^2 \, dx .
\]
\end{theorem}

\begin{proof}
\leanok
For almost every $\xi$, \Cref{lem:fourier-sum-conv} and the Cauchy--Schwarz inequality of
\Cref{lem:l2-squares} give
\[
  \Bigl|\mathcal F\Bigl(\sum_{j \in s} \theta_j * g_j\Bigr)(\xi)\Bigr|^2
  \leq \Bigl(\sum_{j \in s} |\widehat{\theta_j}(\xi)|^2\Bigr)
  \sum_{j \in s} |\widehat{g_j}(\xi)|^2
  \leq B \sum_{j \in s} |\widehat{g_j}(\xi)|^2 .
\]
Integration and \Cref{lem:plancherel-l2} on both sides give the first inequality. The second is
the first for the classes of the $g_j$ (\Cref{lem:plancherel-l2}), with
$\sum_{j \in s} |g_j(x)|^2 \leq \|(g_j(x))_j\|_{\ell_2}^2$ (\Cref{def:lq-norm}).
\end{proof}

\begin{corollary}[Plancherel bounds for Calder\'on--Zygmund families]
\label{cor:cz-family-l2-bounds}
\lean{Rd.IsCZFamily.tsum_enorm_convLp_sq_le, Rd.IsCZFamily.lintegral_lqNorm_conv_sq_le,
  Rd.IsCZFamily.enorm_sum_convLp_sq_le, Rd.IsCZFamily.lintegral_enorm_sum_conv_sq_le}
\leanok
\uses{def:cz-family, thm:l2-analysis-bound, thm:l2-synthesis-bound}
For all $B, H \in [0, \infty]$ and every Calder\'on--Zygmund family $(\theta_j)_{j \in \N}$ on
$V$ with overlap $B$ and H\"ormander constant $H$,
\[
  \sum_j \|\theta_j * f\|_2^2 \leq B \, \|f\|_2^2 , \qquad
  \int_V \|((\theta_j * f)(x))_j\|_{\ell_2}^2 \, dx \leq B \int_V |f(x)|^2 \, dx
\]
for every $f \in L_2(V)$, and
\[
  \Bigl\|\sum_{j \in s} \theta_j * g_j\Bigr\|_2^2 \leq B \sum_{j \in s} \|g_j\|_2^2
\]
and
\[
  \int_V \Bigl|\sum_{j \in s} (\theta_j * g_j)(x)\Bigr|^2 \, dx
  \leq B \int_V \|(g_j(x))_j\|_{\ell_2}^2 \, dx
\]
for all $g_j \in L_2(V)$ and every finite set $s \subseteq \N$.
\end{corollary}

\begin{proof}
\leanok
The kernels of a Calder\'on--Zygmund family are integrable, and the overlap bound holds at
every $\xi$ (\Cref{def:cz-family}); apply \Cref{thm:l2-analysis-bound} and
\Cref{thm:l2-synthesis-bound}.
\end{proof}

\section{The weak type \texorpdfstring{$(1,1)$}{(1,1)} of vector-valued singular integrals}
\label{sec:vector-valued-weak-type}

Let $(\theta_j)_{j \in \N}$ be a Calder\'on--Zygmund family on $V$ with overlap $B$ and
H\"ormander constant $H$ (\Cref{def:cz-family}). This section proves the weak type $(1,1)$ of
the analysis operator $f \mapsto (\theta_j * f)_j$ and of the synthesis operators
$(g_j)_j \mapsto \sum_{j \in s} \theta_j * g_j$ for finite sets $s \subseteq \N$. At a height
$t$, the signed Calder\'on--Zygmund decomposition of \Cref{sec:calderon-zygmund-signed} splits
$f$ into the good part $\mathcal G f$, which the Plancherel bounds of
\Cref{cor:cz-family-l2-bounds} control in $L_2$, and the bad part $\mathcal B f$, a sum of
pieces $b_Q f$ with mean zero. A maximal big cube $Q = Q_{k,m}$ has centre $c_Q$ and lies in the
closed ball $\bar B(c_Q, \rho_Q)$ with $\rho_Q = n 2^{-k}$ (\Cref{lem:dyadic-cubes}), and the
enlarged exceptional set $\Omega^*$ is the union of the balls $\bar B(c_Q, 2 \rho_Q)$
(\Cref{def:calderon-zygmund-cubes}). Off $\Omega^*$ the mean zero of the pieces lets the
H\"ormander condition act. This is the classical argument for singular integrals with values in a
Hilbert space~\cite{BenedekCalderonPanzone}; see also Stein~\cite{SteinSingular} and
Grafakos~\cite{Grafakos}.

\begin{lemma}[The H\"ormander condition around a cube]
\label{lem:hormander-around-cube}
\lean{Rd.IsCZFamily.setLIntegral_compl_closedBall_le, Rd.IsCZFamily.measurable_lqNorm_sub,
  Rd.IsCZFamily.setLIntegral_compl_closedBall_setLIntegral_le}
\leanok
\uses{def:cz-family, def:lq-norm}
For all $B, H \in [0, \infty]$, every Calder\'on--Zygmund family $(\theta_j)_j$ on $V$ with
overlap $B$ and H\"ormander constant $H$, every $c \in V$, every $\rho \in \R$ and every
$y \in V$ with $\|y - c\| \leq \rho$,
\[
  \int_{\{\|x - c\| > 2\rho\}} \|(\theta_j(x - y) - \theta_j(x - c))_j\|_{\ell_2} \, dx \leq H .
\]
The function $(x, y) \mapsto \|(\theta_j(x - y) - \theta_j(x - c))_j\|_{\ell_2}$ is measurable
on $V \times V$, and for every measurable $Q \subseteq \bar B(c, \rho)$ and every almost
everywhere measurable $m : V \to [0, \infty]$,
\[
  \int_{\{\|x - c\| > 2\rho\}} \int_Q \|(\theta_j(x - y) - \theta_j(x - c))_j\|_{\ell_2} \,
  m(y) \, dy \, dx \leq H \int_Q m .
\]
\end{lemma}

\begin{proof}
\leanok
Every $x$ with $\|x - c\| > 2\rho$ satisfies $\|x - c\| > 2 \|y - c\|$. The substitution
$x' = x - c$ turns the integral into one over a subset of $\{\|x'\| > 2 \|y - c\|\}$ with
integrand $\|(\theta_j(x' - (y - c)) - \theta_j(x'))_j\|_{\ell_2}$, which is at most $H$ by the
H\"ormander condition at $y - c$ (\Cref{def:cz-family}). The kernels are continuous, so the
function on $V \times V$ is measurable. Tonelli's theorem exchanges the two integrals in the
last bound, and the inner integral is at most $H$ for every $y \in Q$.
\end{proof}

\begin{lemma}[Convolutions of pieces with mean zero]
\label{lem:mean-zero-pieces-conv}
\lean{Rd.enorm_conv_le_setLIntegral_sub, Rd.IsCZFamily.lqNorm_conv_le_setLIntegral,
  Rd.IsCZFamily.enorm_sum_conv_le_setLIntegral}
\leanok
\uses{def:cz-family, lem:reflect-basic, lem:minkowski-integral-lq,
  lem:lq-cauchy-schwarz-sequences}
For every continuous bounded $K : V \to \C$, every set $Q \subseteq V$, every integrable
$b : V \to \C$ with $b = 0$ outside $Q$ and $\int b = 0$, and all $c, x \in V$,
\[
  |(K * b)(x)| \leq \int_Q |K(x - y) - K(x - c)| \, |b(y)| \, dy .
\]
For all $B, H \in [0, \infty]$, every Calder\'on--Zygmund family $(\theta_j)_j$ with overlap
$B$ and H\"ormander constant $H$, every such $b$ and all $c, x \in V$,
\[
  \|((\theta_j * b)(x))_j\|_{\ell_2}
  \leq \int_Q \|(\theta_j(x - y) - \theta_j(x - c))_j\|_{\ell_2} \, |b(y)| \, dy ,
\]
and for every finite set $s \subseteq \N$ and all integrable $b_j : V \to \C$ with $b_j = 0$
outside $Q$ and $\int b_j = 0$,
\[
  \Bigl|\sum_{j \in s} (\theta_j * b_j)(x)\Bigr|
  \leq \int_Q \|(\theta_j(x - y) - \theta_j(x - c))_j\|_{\ell_2} \,
  \|(b_j(y))_j\|_{\ell_2} \, dy .
\]
\end{lemma}

\begin{proof}
\leanok
By \Cref{lem:reflect-basic}, $(K * b)(x) = \int K(x - y) b(y) \, dy$, and since $\int b = 0$
this is $\int (K(x - y) - K(x - c)) b(y) \, dy$; the integrand vanishes outside $Q$. The second
bound follows by Minkowski's integral inequality in $\ell_2$
(\Cref{lem:minkowski-integral-lq}). For the third, the triangle inequality and the first bound
for each $j \in s$ give
$\int_Q \sum_{j \in s} |\theta_j(x - y) - \theta_j(x - c)| \, |b_j(y)| \, dy$, and the
Cauchy--Schwarz inequality of \Cref{lem:lq-cauchy-schwarz-sequences} applies at each $y$.
\end{proof}

\begin{lemma}[Convolution of the bad part]
\label{lem:conv-bad-part}
\lean{Rd.hasSum_conv_czBadPart}
\leanok
\uses{def:calderon-zygmund-signed, lem:calderon-zygmund-bad-pieces,
  thm:calderon-zygmund-bad-pieces-mean-mass, lem:reflect-basic}
For every $G : V \to [0, \infty]$ and every $t \in [0, \infty]$, with $Q$ running over the
maximal big cubes of $G$ at height $t$: for every continuous bounded $K : V \to \C$, every
integrable $f : V \to \C$ and every $x \in V$,
\[
  (K * \mathcal B f)(x) = \sum_Q (K * b_Q f)(x) ,
\]
and the series converges absolutely.
\end{lemma}

\begin{proof}
\leanok
With $|K| \leq M$,
$\sum_Q \int |K(x - y)| \, |b_Q f(y)| \, dy \leq M \sum_Q \int |b_Q f| \leq 2 M \int |f|$
by \Cref{thm:calderon-zygmund-bad-pieces-mean-mass}, so the series of the integrals
$\int K(x - y) \, b_Q f(y) \, dy$ converges absolutely to the integral of the series, which is
$\int K(x - y) \, \mathcal B f(y) \, dy$ by \Cref{lem:calderon-zygmund-bad-pieces}. Both sides
are convolutions by \Cref{lem:reflect-basic}.
\end{proof}

\begin{lemma}[The bad part off the enlarged exceptional set]
\label{lem:bad-part-off-enlarged}
\lean{Rd.IsCZFamily.setLIntegral_compl_czEnlarged_tsum_le,
  Rd.IsCZFamily.setLIntegral_lqNorm_conv_czBadC_le,
  Rd.IsCZFamily.setLIntegral_enorm_sum_conv_czBadC_le}
\leanok
\uses{lem:hormander-around-cube, lem:mean-zero-pieces-conv, lem:conv-bad-part,
  lem:lp-series-minkowski, thm:calderon-zygmund-bad-pieces-mean-mass,
  thm:calderon-zygmund-signed-family, lem:dyadic-cubes, def:calderon-zygmund-cubes}
For all $B, H \in [0, \infty]$, every Calder\'on--Zygmund family $(\theta_j)_j$ with overlap
$B$ and H\"ormander constant $H$, every $G : V \to [0, \infty]$ and every $t \in [0, \infty]$,
with the enlarged exceptional set $\Omega^*$ of the maximal big cubes $Q$ of $G$ at height $t$:
for every family of almost everywhere measurable $m_Q : V \to [0, \infty]$,
\[
  \int_{V \setminus \Omega^*} \sum_Q \int_Q
  \|(\theta_j(x - y) - \theta_j(x - c_Q))_j\|_{\ell_2} \, m_Q(y) \, dy \, dx
  \leq H \sum_Q \int_Q m_Q ;
\]
for every integrable $f : V \to \C$,
\[
  \int_{V \setminus \Omega^*} \|((\theta_j * \mathcal B f)(x))_j\|_{\ell_2} \, dx
  \leq 2 H \int |f| ;
\]
and for every finite set $s \subseteq \N$ and all integrable $g_j : V \to \C$,
\[
  \int_{V \setminus \Omega^*} \Bigl|\sum_{j \in s} (\theta_j * \mathcal B g_j)(x)\Bigr| \, dx
  \leq 2 H \int \|(g_j(x))_j\|_{\ell_2} \, dx .
\]
\end{lemma}

\begin{proof}
\leanok
A point $x \notin \Omega^*$ satisfies $\|x - c_Q\| > 2\rho_Q$ for every maximal big cube $Q$,
and $Q \subseteq \bar B(c_Q, \rho_Q)$ (\Cref{lem:dyadic-cubes}). Monotone convergence and
\Cref{lem:hormander-around-cube} for each $Q$ give the first bound. By \Cref{lem:conv-bad-part},
Minkowski's inequality for series (\Cref{lem:lp-series-minkowski}) and
\Cref{lem:mean-zero-pieces-conv},
\[
  \|((\theta_j * \mathcal B f)(x))_j\|_{\ell_2}
  \leq \sum_Q \int_Q \|(\theta_j(x - y) - \theta_j(x - c_Q))_j\|_{\ell_2} \,
  |b_Q f(y)| \, dy
\]
for every $x$, so the first bound with $m_Q = |b_Q f|$ and
\Cref{thm:calderon-zygmund-bad-pieces-mean-mass} give the second. The third follows in the same
way with $m_Q(y) = \|(b_Q g_j(y))_j\|_{\ell_2}$ and \Cref{thm:calderon-zygmund-signed-family}.
\end{proof}

\begin{lemma}[The good part and the splitting]
\label{lem:good-part-split}
\lean{Rd.IsCZFamily.ae_lqNorm_conv_le_add, Rd.IsCZFamily.ae_enorm_sum_conv_le_add,
  Rd.IsCZFamily.lintegral_lqNorm_conv_czGoodC_sq_le,
  Rd.IsCZFamily.lintegral_enorm_sum_conv_czGoodC_sq_le}
\leanok
\uses{cor:cz-family-l2-bounds, thm:calderon-zygmund-signed-good-bounds,
  thm:calderon-zygmund-signed-family, lem:calderon-zygmund-signed-good,
  lem:calderon-zygmund-bad-pieces, lem:conv-assoc, lem:weak-type-elementary, lem:lq-basic}
For all $B, H \in [0, \infty]$, every Calder\'on--Zygmund family $(\theta_j)_j$ with overlap
$B$ and H\"ormander constant $H$, every $G : V \to [0, \infty]$ and every $t \in [0, \infty]$:
for every integrable $f : V \to \C$ and almost every $x \in V$,
\[
  \|((\theta_j * f)(x))_j\|_{\ell_2} \leq \|((\theta_j * \mathcal G f)(x))_j\|_{\ell_2}
  + \|((\theta_j * \mathcal B f)(x))_j\|_{\ell_2} ,
\]
and for every finite set $s \subseteq \N$, all integrable $g_j : V \to \C$ and almost every
$x \in V$,
\[
  \Bigl|\sum_{j \in s} (\theta_j * g_j)(x)\Bigr|
  \leq \Bigl|\sum_{j \in s} (\theta_j * \mathcal G g_j)(x)\Bigr|
  + \Bigl|\sum_{j \in s} (\theta_j * \mathcal B g_j)(x)\Bigr| .
\]
If moreover $n \geq 1$ and $0 < t < \infty$, then for every integrable $f : V \to \C$ and the
cubes of $|f|$ at height $t$,
\[
  \int \|((\theta_j * \mathcal G f)(x))_j\|_{\ell_2}^2 \, dx \leq B \, 2^n t \int |f| ,
\]
and for every finite set $s \subseteq \N$, all almost everywhere strongly measurable
$g_j : V \to \C$ with $\int \|(g_j(x))_j\|_{\ell_2} \, dx < \infty$ and the cubes of
$x \mapsto \|(g_j(x))_j\|_{\ell_2}$ at height $t$,
\[
  \int \Bigl|\sum_{j \in s} (\theta_j * \mathcal G g_j)(x)\Bigr|^2 \, dx
  \leq B \, 2^n t \int \|(g_j(x))_j\|_{\ell_2} \, dx .
\]
\end{lemma}

\begin{proof}
\leanok
By \Cref{lem:calderon-zygmund-bad-pieces}, $f = \mathcal G f + \mathcal B f$, and both parts
are integrable
(\Cref{lem:calderon-zygmund-signed-good}, \Cref{lem:calderon-zygmund-bad-pieces}); convolution
is additive almost everywhere (\Cref{lem:conv-assoc}), and Minkowski's inequality in $\ell_2$
(\Cref{lem:lq-basic}) gives the first bound, the triangle inequality the second. For the good
parts, $\int |\mathcal G f|^2 \leq 2^n t \int |f| < \infty$
(\Cref{thm:calderon-zygmund-signed-good-bounds}), and
$|\mathcal G g_j| \leq \|(\mathcal G g_i)_i\|_{\ell_2}$, whose square has integral at most
$2^n t \int \|(g_i)_i\|_{\ell_2}$ (\Cref{thm:calderon-zygmund-signed-family}). So the good parts
lie in $L_2$ (\Cref{lem:weak-type-elementary}), and \Cref{cor:cz-family-l2-bounds} gives the
last two bounds.
\end{proof}

\begin{definition}[The weak type constant]
\label{def:vector-weak-const}
\lean{Rd.vectorWeakConst, Rd.vectorWeakConst_ne_top}
\leanok
For all $B, H \in [0, \infty]$,
\[
  C_V(B, H) = 4 \cdot 2^n B + (2n)^n \, |B(0, 1)| + 4 H ,
\]
where $|B(0, 1)|$ is the volume of the unit ball of $V$. It is finite if $B$ and $H$ are.
\end{definition}

\begin{theorem}[The weak type $(1,1)$ of vector-valued singular integrals]
\label{thm:vector-valued-weak-type}
\lean{Rd.IsCZFamily.volume_lqNorm_conv_gt_le, Rd.IsCZFamily.volume_enorm_sum_conv_gt_le}
\leanok
\uses{lem:level-sets-split-function, lem:good-part-split, lem:bad-part-off-enlarged,
  lem:calderon-zygmund-averages, def:vector-weak-const}
Let $n \geq 1$. For all $B, H \in [0, \infty]$, every Calder\'on--Zygmund family
$(\theta_j)_j$ on $V$ with overlap $B$ and H\"ormander constant $H$ and every
$t \in (0, \infty)$:
for every integrable $f : V \to \C$,
\[
  \bigl|\{x \in V : \|((\theta_j * f)(x))_j\|_{\ell_2} > t\}\bigr|
  \leq C_V(B, H) \, t^{-1} \int |f| ,
\]
and for every finite set $s \subseteq \N$ and all almost everywhere strongly measurable
$g_j : V \to \C$ with $\int \|(g_j(x))_j\|_{\ell_2} \, dx < \infty$,
\[
  \Bigl|\Bigl\{x \in V : \Bigl|\sum_{j \in s} (\theta_j * g_j)(x)\Bigr| > t\Bigr\}\Bigr|
  \leq C_V(B, H) \, t^{-1} \int \|(g_j(x))_j\|_{\ell_2} \, dx .
\]
\end{theorem}

\begin{proof}
\leanok
Decompose $f$ along the maximal big cubes of $|f|$ at height $t$, and every $g_j$ along those
of $\|(g_j)_j\|_{\ell_2}$; the $g_j$ are integrable, since $|g_j| \leq \|(g_i)_i\|_{\ell_2}$.
Apply \Cref{lem:level-sets-split-function} with the sizes of the good and of the bad parts as $a$
and $b$ and with $S = \Omega^*$: \Cref{lem:good-part-split} gives the splitting and $A = B 2^n$,
\Cref{lem:calderon-zygmund-averages} gives $D = (2n)^n |B(0, 1)|$, and
\Cref{lem:bad-part-off-enlarged} gives $K = 2H$. Then $4A + D + 2K = C_V(B, H)$.
\end{proof}

\section{The strong type of vector-valued singular integrals}
\label{sec:vector-valued-strong-type}

Let $(\theta_j)_{j \in \N}$ be a Calder\'on--Zygmund family on $V$ with overlap $B$ and
H\"ormander constant $H$ (\Cref{def:cz-family}). This section proves, for every
$1 < p < \infty$, the strong type $(p, p)$ of the analysis operator
$f \mapsto \|((\theta_j * f)(x))_j\|_{\ell_2}$ and of the synthesis operators
$(g_j)_j \mapsto \sum_{j \in s} \theta_j * g_j$ for finite sets $s \subseteq \N$, with explicit
constants. For $p = 2$ these are the Plancherel bounds. For $1 < p < 2$ Marcinkiewicz
interpolation between the weak type $(1,1)$ of \Cref{thm:vector-valued-weak-type} and the
$L_2$ bounds applies: to the analysis operator in the form of the formalisation of Carleson's
theorem (\Cref{thm:marcinkiewicz-interpolation}), and to the synthesis operators in the form
for an abstract input type (\Cref{thm:marcinkiewicz-abstract}), whose inputs are families of
measurable functions. For $2 < p < \infty$, the adjoint identity exchanges the two operators of
the adjoint family $\theta_j^*(x) = \overline{\theta_j(-x)}$, which has the same overlap and
H\"ormander constant, at the conjugate exponent $p' < 2$, and \Cref{thm:lp-norm-duality}
returns to $L_p$. These are the classical $L_p$ bounds of vector-valued singular
integrals~\cite{BenedekCalderonPanzone}; see also Stein~\cite{SteinSingular}.

\begin{theorem}[The strong type $(2,2)$]
\label{thm:cz-strong-type-two}
\lean{Rd.IsCZFamily.hasStrongType_lqNorm_conv_two, Rd.IsCZFamily.eLpNorm_lqNorm_conv_two_le,
  Rd.IsCZFamily.eLpNorm_sum_conv_two_le}
\leanok
\uses{cor:cz-family-l2-bounds, lem:lp-norm-from-power-integrals, def:weak-type-operators,
  lem:lp-norm-nonnegative}
For all $B, H \in [0, \infty]$, every Calder\'on--Zygmund family $(\theta_j)_j$ on $V$ with
overlap $B$ and H\"ormander constant $H$, every $f \in L_2(V)$, every finite set
$s \subseteq \N$ and all $g_j \in L_2(V)$,
\[
  \bigl\| \|((\theta_j * f)(x))_j\|_{\ell_2} \bigr\|_2 \leq B^{1/2} \|f\|_2 , \qquad
  \Bigl\| \sum_{j \in s} \theta_j * g_j \Bigr\|_2
  \leq B^{1/2} \bigl\| \|(g_j(x))_j\|_{\ell_2} \bigr\|_2 .
\]
In particular the analysis operator has strong type $(2, 2)$ with constant $B^{1/2}$.
\end{theorem}

\begin{proof}
\leanok
The function $x \mapsto \|((\theta_j * f)(x))_j\|_{\ell_2}$ is almost everywhere measurable
(\Cref{lem:lp-norm-nonnegative}). \Cref{lem:lp-norm-from-power-integrals} with $p = 2$ and
$C = B$ turns the integral bounds of \Cref{cor:cz-family-l2-bounds} into the norm bounds.
\end{proof}

\begin{lemma}[Subadditivity and weak type $(1,1)$ of the analysis operator]
\label{lem:analysis-subadditive-weak-type}
\lean{Rd.IsCZFamily.aeSubadditiveOn_lqNorm_conv, Rd.IsCZFamily.hasWeakType_lqNorm_conv}
\leanok
\uses{def:cz-family, lem:conv-assoc, lem:lq-basic, thm:vector-valued-weak-type,
  def:weak-type-operators, def:vector-weak-const, lem:distribution-weak-norm-basic}
For all $B, H \in [0, \infty]$ and every Calder\'on--Zygmund family $(\theta_j)_j$ on $V$ with
overlap $B$ and H\"ormander constant $H$, the analysis operator
$f \mapsto \|((\theta_j * f)(x))_j\|_{\ell_2}$ is almost everywhere subadditive with constant
$1$ on the functions in $L_1(V)$ or $L_2(V)$. If $n \geq 1$, it has weak type $(1, 1)$ with
constant $C_V(B, H)$.
\end{lemma}

\begin{proof}
\leanok
For $f \in L_p(V)$ and $g \in L_r(V)$ with $p, r \in \{1, 2\}$,
$\theta_j * (f + g) = \theta_j * f + \theta_j * g$ almost everywhere for every $j$
(\Cref{lem:conv-assoc}), and Minkowski's inequality in $\ell_2$ (\Cref{lem:lq-basic}) gives the
subadditivity. For $t > 0$, \Cref{thm:vector-valued-weak-type} gives
$t \, |\{\|((\theta_j * f)(x))_j\|_{\ell_2} > t\}| \leq C_V(B, H) \int |f|$, and for $t = 0$
the left side vanishes; the supremum over $t$ is the weak $L_1$ norm
(\Cref{lem:distribution-weak-norm-basic}).
\end{proof}

\begin{definition}[The strong type constant below $2$]
\label{def:cz-const-below}
\lean{Rd.czConstBelow, Rd.czConstBelow_ne_top}
\leanok
\uses{def:marcinkiewicz-const, def:vector-weak-const}
For $1 < p < 2$ and $B, H \in [0, \infty]$,
\[
  S^-_{V,p}(B, H) = C^{\mathrm M}_p\bigl(C_V(B, H), B\bigr)^{1/p} .
\]
It is finite if $B$ and $H$ are.
\end{definition}

\begin{theorem}[The analysis operator for $1 < p < 2$]
\label{thm:cz-analysis-below-two}
\lean{Rd.IsCZFamily.hasStrongType_lqNorm_conv_of_lt, Rd.IsCZFamily.eLpNorm_lqNorm_conv_le_of_lt}
\leanok
\uses{thm:marcinkiewicz-interpolation, lem:analysis-subadditive-weak-type,
  thm:cz-strong-type-two, def:cz-const-below}
Let $n \geq 1$ and $1 < p < 2$. For all $B, H \in [0, \infty]$ and every Calder\'on--Zygmund
family $(\theta_j)_j$ on $V$ with overlap $B$ and H\"ormander constant $H$, the analysis
operator has strong type $(p, p)$ with constant $S^-_{V,p}(B, H)$: for every $f \in L_p(V)$,
\[
  \bigl\| \|((\theta_j * f)(x))_j\|_{\ell_2} \bigr\|_p \leq S^-_{V,p}(B, H) \, \|f\|_p .
\]
\end{theorem}

\begin{proof}
\leanok
\Cref{thm:marcinkiewicz-interpolation} applies with $A = 1$, $C_1 = C_V(B, H)$ and
$C_2 = B^{1/2}$ (\Cref{lem:analysis-subadditive-weak-type}, \Cref{thm:cz-strong-type-two}), and
$N_p(1, C_V(B, H), B^{1/2}) = S^-_{V,p}(B, H)$.
\end{proof}

\begin{theorem}[The synthesis operators for $1 < p < 2$]
\label{thm:cz-synthesis-below-two}
\lean{Rd.IsCZFamily.eLpNorm_sum_conv_le_of_lt_of_measurable,
  Rd.IsCZFamily.eLpNorm_sum_conv_le_of_lt}
\leanok
\uses{thm:marcinkiewicz-abstract, thm:vector-valued-weak-type, cor:cz-family-l2-bounds,
  lem:finite-families-lp-l2, lem:fefferman-stein-split, lem:conv-assoc, def:conv,
  def:cz-const-below}
Let $n \geq 1$ and $1 < p < 2$. For all $B, H \in [0, \infty]$, every Calder\'on--Zygmund family
$(\theta_j)_j$ on $V$ with overlap $B$ and H\"ormander constant $H$, every finite set
$s \subseteq \N$ and all $g_j \in L_p(V)$, $j \in \N$,
\[
  \Bigl\| \sum_{j \in s} \theta_j * g_j \Bigr\|_p
  \leq S^-_{V,p}(B, H) \, \bigl\| \|(g_j(x))_j\|_{\ell_2} \bigr\|_p .
\]
\end{theorem}

\begin{proof}
\leanok
First let every $g_j$ be measurable with
$\bigl\| \|(g_j)_j\|_{\ell_2} \bigr\|_p < \infty$. Apply \Cref{thm:marcinkiewicz-abstract} to
the families of measurable functions, with the size $G_g = \|(g_j)_j\|_{\ell_2}$, the
restriction $g|_S = (\mathbf 1_S g_j)_j$, whose size is $\mathbf 1_S G_g$
(\Cref{lem:fefferman-stein-split}), and
$T g = \bigl|\sum_{j \in s} \theta_j * g_j\bigr|$. If $\int G_g^p < \infty$, every $g_j$ lies in
$L_p(V)$ (\Cref{lem:finite-families-lp-l2}), so convolution is additive almost everywhere
(\Cref{lem:conv-assoc}) and $T g \leq T(g|_S) + T(g|_{V \setminus S})$.
\Cref{thm:vector-valued-weak-type} gives the weak bound with $C_1 = C_V(B, H)$ and
\Cref{cor:cz-family-l2-bounds} the $L_2$ bound with $C_2 = B$. In general, replace $g_j$ by a
measurable function equal to it almost everywhere for $j \in s$ and by $0$ for $j \notin s$: the
sum is unchanged (\Cref{def:conv}), the right side can only decrease, and it becomes finite
(\Cref{lem:finite-families-lp-l2}).
\end{proof}

\begin{lemma}[Pairings through the adjoint family]
\label{lem:cz-pairings-adjoint}
\lean{Rd.IsCZFamily.enorm_sum_integral_conv_mul_conj_le,
  Rd.IsCZFamily.enorm_integral_sum_conv_mul_conj_le}
\leanok
\uses{cor:convolution-adjoint-conj, lem:cz-family-adjoint, thm:cz-analysis-below-two,
  thm:cz-synthesis-below-two, lem:holder-lp-norms, lem:lq-cauchy-schwarz-sequences,
  thm:young}
Let $n \geq 1$, $1 < q < 2$ and $p = q/(q - 1)$. For all $B, H \in [0, \infty]$, every
Calder\'on--Zygmund family $(\theta_j)_j$ on $V$ with overlap $B$ and H\"ormander constant $H$
and every finite set $s \subseteq \N$: for every $f \in L_p(V)$ and all $g_j \in L_q(V)$,
\[
  \Bigl| \sum_{j \in s} \int_V (\theta_j * f) \, \overline{g_j} \Bigr|
  \leq S^-_{V,q}(B, H) \, \|f\|_p \, \bigl\| \|(g_j(x))_j\|_{\ell_2} \bigr\|_q ,
\]
and for all $g_j \in L_p(V)$ and every $u \in L_q(V)$,
\[
  \Bigl| \int_V \Bigl(\sum_{j \in s} \theta_j * g_j\Bigr) \overline u \Bigr|
  \leq S^-_{V,q}(B, H) \, \bigl\| \|(g_j(x))_j\|_{\ell_2} \bigr\|_p \, \|u\|_q .
\]
\end{lemma}

\begin{proof}
\leanok
By \Cref{cor:convolution-adjoint-conj}, every pairing moves $\theta_j$ to the other factor as
$\theta_j^*$, and $(\theta_j^*)_j$ is a Calder\'on--Zygmund family with overlap $B$ and
H\"ormander constant $H$ (\Cref{lem:cz-family-adjoint}). The first pairing becomes
$\int_V f \, \overline{\sum_{j \in s} \theta_j^* * g_j}$, and H\"older's inequality
(\Cref{lem:holder-lp-norms}) and \Cref{thm:cz-synthesis-below-two} for the adjoint family
at $q$ bound it. The second becomes
$\sum_{j \in s} \int_V g_j \, \overline{\theta_j^* * u}$; the Cauchy--Schwarz inequality at
every point (\Cref{lem:lq-cauchy-schwarz-sequences}), H\"older's inequality and
\Cref{thm:cz-analysis-below-two} for the adjoint family at $q$ bound it.
\end{proof}

\begin{theorem}[The strong type for $2 < p < \infty$]
\label{thm:cz-above-two}
\lean{Rd.IsCZFamily.eLpNorm_lqNorm_conv_le_of_two_lt, Rd.IsCZFamily.eLpNorm_sum_conv_le_of_two_lt}
\leanok
\uses{lem:cz-pairings-adjoint, thm:lp-norm-duality, lem:countable-minkowski, thm:young,
  lem:cancel-conjugate-power, lem:lp-series-minkowski, def:cz-const-below}
Let $n \geq 1$, $2 < p < \infty$ and $p' = p/(p - 1)$. For all $B, H \in [0, \infty]$ and every
Calder\'on--Zygmund family $(\theta_j)_j$ on $V$ with overlap $B$ and H\"ormander constant $H$:
for every $f \in L_p(V)$,
\[
  \bigl\| \|((\theta_j * f)(x))_j\|_{\ell_2} \bigr\|_p \leq S^-_{V,p'}(B, H) \, \|f\|_p ,
\]
and for every finite set $s \subseteq \N$ and all $g_j \in L_p(V)$,
\[
  \Bigl\| \sum_{j \in s} \theta_j * g_j \Bigr\|_p
  \leq S^-_{V,p'}(B, H) \, \bigl\| \|(g_j(x))_j\|_{\ell_2} \bigr\|_p .
\]
\end{theorem}

\begin{proof}
\leanok
By \Cref{lem:cancel-conjugate-power}, $1 < p' < 2$. For every $J \in \N$, the family
$u_j = \theta_j * f$ for $j < J$ and $u_j = 0$ for $j \geq J$ lies in $L_p(V)$
(\Cref{thm:young}), and the first bound of \Cref{lem:cz-pairings-adjoint} is the hypothesis of
\Cref{thm:lp-norm-duality} with $A = S^-_{V,p'}(B, H) \|f\|_p$. The truncated $\ell_2$ norms
increase to $\|((\theta_j * f)(x))_j\|_{\ell_2}$ (\Cref{lem:countable-minkowski}), and monotone
convergence in $L_p$ (\Cref{lem:lp-series-minkowski}) gives the first inequality. The second
is \Cref{thm:lp-norm-duality} for
$v = \sum_{j \in s} \theta_j * g_j \in L_p(V)$ (\Cref{thm:young}) with the second bound of
\Cref{lem:cz-pairings-adjoint}.
\end{proof}

\begin{definition}[The strong type constant]
\label{def:cz-const}
\lean{Rd.czConst, Rd.czConst_ne_top}
\leanok
\uses{def:cz-const-below}
For $1 < p < \infty$ and $B, H \in [0, \infty]$,
\[
  S_{V,p}(B, H) =
  \begin{cases}
    S^-_{V,p}(B, H) & \text{for } p < 2, \\
    B^{1/2} & \text{for } p = 2, \\
    S^-_{V,p/(p-1)}(B, H) & \text{for } p > 2 .
  \end{cases}
\]
It is finite if $B$ and $H$ are.
\end{definition}

\begin{theorem}[The strong type of vector-valued singular integrals]
\label{thm:vector-valued-strong-type}
\lean{Rd.IsCZFamily.eLpNorm_lqNorm_conv_le, Rd.IsCZFamily.eLpNorm_sum_conv_le}
\leanok
\uses{thm:cz-strong-type-two, thm:cz-analysis-below-two, thm:cz-synthesis-below-two,
  thm:cz-above-two, def:cz-const}
Let $n \geq 1$ and $1 < p < \infty$. For all $B, H \in [0, \infty]$, every
Calder\'on--Zygmund family $(\theta_j)_j$ on $V$ with overlap $B$ and H\"ormander constant $H$,
every $f \in L_p(V)$, every finite set $s \subseteq \N$ and all $g_j \in L_p(V)$,
\[
  \bigl\| \|((\theta_j * f)(x))_j\|_{\ell_2} \bigr\|_p \leq S_{V,p}(B, H) \, \|f\|_p
\]
and
\[
  \Bigl\| \sum_{j \in s} \theta_j * g_j \Bigr\|_p
  \leq S_{V,p}(B, H) \, \bigl\| \|(g_j(x))_j\|_{\ell_2} \bigr\|_p .
\]
\end{theorem}

\begin{proof}
\leanok
\Cref{thm:cz-analysis-below-two} and \Cref{thm:cz-synthesis-below-two} for $p < 2$,
\Cref{thm:cz-strong-type-two} for $p = 2$, and \Cref{thm:cz-above-two} for $p > 2$.
\end{proof}

\section{The square function theorem}
\label{sec:square-function}

For $f \in L_p(V)$ the square function $x \mapsto \|((\psi_j * f)(x))_j\|_{\ell_2}$ collects the
Littlewood--Paley blocks at every point, and its $L_p$ norm is $\|f\|_{\TL{0}{p}{2}}$ by the
definition of the Triebel--Lizorkin norm. The square function theorem is the two-sided bound of
this norm by $\|f\|_p$ for $1 < p < \infty$, so $\TL{0}{p}{2} = L_p(V)$. The upper bound is the
strong type of the analysis operator of the Calder\'on--Zygmund family $(\psi_j)_j$; the lower
bound combines the reproducing formula
$\sum_{j < J} \tilde\psi_j * (\psi_j * f) \to f$ with the strong type of the synthesis
operators of the family $(\tilde\psi_j)_j$ (\Cref{thm:vector-valued-strong-type}). Stein's
book~\cite{SteinSingular} proves the Littlewood--Paley theorem through vector-valued singular
integrals, with the lower bound by duality, and Triebel~\cite{Triebel} states it as
$\TL{0}{p}{2} = L_p(V)$. Grafakos~\cite{Grafakos} gives a textbook account of the
Littlewood--Paley theorem and of vector-valued singular integrals.

\begin{definition}[The square function constants]
\label{def:square-function-consts}
\lean{Rd.squareFunctionConst, Rd.squareFunctionConst_ne_top, Rd.squareSynthesisConst,
  Rd.squareSynthesisConst_ne_top}
\leanok
\uses{def:cz-const, def:hormander-const, def:psi-base, def:psi-dual}
For $1 < p < \infty$,
\[
  Q_{V,p} = S_{V,p}\bigl(2, 4 H_{A^{\psi}}\bigr) , \qquad
  \tilde Q_{V,p} = S_{V,p}\bigl(4, 4 H_{\tilde A}\bigr) .
\]
Both are finite.
\end{definition}

The constants are formulas in the Schwartz seminorms $A^\psi$ and $\tilde A$ of the fixed blocks
and synthesis kernels.

\begin{theorem}[The upper square function bound]
\label{thm:square-function-upper}
\lean{Rd.triebelLizorkinNorm_zero_two_le, Rd.memTriebelLizorkin_zero_two}
\leanok
\uses{def:triebel-lizorkin, prop:tl-elementary, thm:cz-family-psi,
  thm:vector-valued-strong-type, def:square-function-consts}
For every $n \geq 1$, every $1 < p < \infty$ and every $f \in L_p(V)$,
\[
  \|f\|_{\TL{0}{p}{2}} \leq Q_{V,p} \, \|f\|_p ,
\]
and in particular $f \in \TL{0}{p}{2}$.
\end{theorem}

\begin{proof}
\leanok
By \Cref{prop:tl-elementary} the norm $\|f\|_{\TL{0}{p}{2}}$ is the $L_p$ norm of the square
function. The blocks form a Calder\'on--Zygmund family with overlap $2$ and H\"ormander
constant $4 H_{A^{\psi}}$ (\Cref{thm:cz-family-psi}), and \Cref{thm:vector-valued-strong-type}
bounds its analysis operator.
\end{proof}

\begin{lemma}[The reproducing formula]
\label{lem:square-function-reproducing}
\lean{Rd.conv_psiDual_blockFun_ae, Rd.sum_conv_psiDual_blockFun_ae,
  Rd.tendsto_eLpNorm_sum_conv_psiDual_blockFun_sub}
\leanok
\uses{def:triebel-lizorkin, lem:conv-assoc, lem:psi-dual-reproduces, lem:lp-blocks,
  lem:lp-kernels-basic, lem:psi-dual-basic}
Let $1 \leq p < \infty$ and $f \in L_p(V)$. For every $j \in \N$,
$\tilde\psi_j * (\psi_j * f) = \psi_j * f$ almost everywhere. For every finite set
$s \subseteq \N$,
\[
  \sum_{j \in s} \tilde\psi_j * (\psi_j * f) = \sum_{j \in s} \Delta_j f
\]
almost everywhere, and $\sum_{j < J} \tilde\psi_j * (\psi_j * f) \to f$ in $L_p(V)$ as
$J \to \infty$.
\end{lemma}

\begin{proof}
\leanok
The kernels $\psi_j$ and $\tilde\psi_j$ are integrable (\Cref{lem:lp-kernels-basic},
\Cref{lem:psi-dual-basic}). By associativity (\Cref{lem:conv-assoc}) and
$\tilde\psi_j * \psi_j = \psi_j$ (\Cref{lem:psi-dual-reproduces}),
$\tilde\psi_j * (\psi_j * f) = (\tilde\psi_j * \psi_j) * f = \psi_j * f$ almost everywhere, and
$\psi_j * f$ represents $\Delta_j f$ (\Cref{def:triebel-lizorkin}). The partial sums
$\sum_{j < J} \Delta_j f$ tend to $f$ in $L_p(V)$ (\Cref{lem:lp-blocks}).
\end{proof}

\begin{theorem}[The lower square function bound]
\label{thm:square-function-lower}
\lean{Rd.enorm_le_triebelLizorkinNorm_zero_two}
\leanok
\uses{lem:square-function-reproducing, thm:cz-family-psi-dual, thm:vector-valued-strong-type,
  thm:square-function-upper, def:square-function-consts, lem:lp-blocks, lem:tl-pointwise-blocks}
Let $n \geq 1$ and $1 < p < \infty$. For every $f \in L_p(V)$,
\[
  \|f\|_p \leq \tilde Q_{V,p} \, \|f\|_{\TL{0}{p}{2}} .
\]
\end{theorem}

\begin{proof}
\leanok
The blocks $\psi_j * f$ lie in $L_p(V)$ (\Cref{lem:tl-pointwise-blocks}). For every $J \in \N$,
\Cref{lem:square-function-reproducing} and
\Cref{thm:vector-valued-strong-type} for the Calder\'on--Zygmund family $(\tilde\psi_j)_j$ with
overlap $4$ and H\"ormander constant $4 H_{\tilde A}$ (\Cref{thm:cz-family-psi-dual}) give
\begin{multline*}
  \Bigl\| \sum_{j < J} \Delta_j f \Bigr\|_p
  = \Bigl\| \sum_{j < J} \tilde\psi_j * (\psi_j * f) \Bigr\|_p \\
  \leq \tilde Q_{V,p} \, \bigl\| \|((\psi_j * f)(x))_j\|_{\ell_2} \bigr\|_p
  = \tilde Q_{V,p} \, \|f\|_{\TL{0}{p}{2}} ,
\end{multline*}
and the left side tends to $\|f\|_p$ as $J \to \infty$ (\Cref{lem:lp-blocks}).
\end{proof}

\begin{corollary}[The square function theorem]
\label{cor:square-function-theorem}
\lean{Rd.triebelLizorkinSpace_zero_two_eq_top, Rd.memTriebelLizorkin_two_of_memApprox}
\leanok
\uses{thm:square-function-upper, thm:square-function-lower, prop:tl-elementary,
  thm:tl-besov-sandwich}
Let $n \geq 1$ and $1 < p < \infty$. Then $\TL{0}{p}{2} = L_p(V)$, and for every
$f \in L_p(V)$,
\[
  \|f\|_p \leq \tilde Q_{V,p} \, \|f\|_{\TL{0}{p}{2}} , \qquad
  \|f\|_{\TL{0}{p}{2}} \leq Q_{V,p} \, \|f\|_p .
\]
For every $s \geq 0$, $\Besov{s}{p}{\min(p,2)} \subseteq \TL{s}{p}{2}$.
\end{corollary}

\begin{proof}
\leanok
\Cref{thm:square-function-upper} and \Cref{thm:square-function-lower}. For $s > 0$ the inclusion
is the Besov sandwich for $q = 2$ (\Cref{thm:tl-besov-sandwich}); for $s = 0$ every
$f \in L_p(V)$ lies in $\TL{0}{p}{2}$.
\end{proof}

\section{Vector-valued singular integrals on the torus}
\label{sec:torus-singular-integrals}

This section transfers the strong type of vector-valued singular integrals
(\Cref{thm:vector-valued-strong-type}) to the torus $\T^d$, on which kernels $\theta_j$ on
$\R^d$ act by $(\theta_j * f)(t) = \int_{\R^d} \theta_j(y) f(t - \pi y) \, dy$
(\Cref{def:torus-econv}). The transfer localises the periodic lift, as for the Fefferman--Stein
inequality (\Cref{thm:fefferman-stein-periodic}). At a point $x$ of the period box $P_0$,
$(\theta_j * f)(\pi x) = (\theta_j * (f \circ \pi))(x)$ (\Cref{lem:torus-lift-conv}). The part
$\mathbf 1_T (f \circ \pi)$ of the lift on the box $T = (-1, 2]^d$ is a function on $\R^d$ with
$L_p$ norm $3^{\#d/p} \|f\|_p$, to which the bounds on $\R^d$ apply; the rest is seen by
$\theta_j$ only at points $\|y\| \geq 1$, and the mass of the kernels there bounds it by Young's
inequality. Throughout, $\#d$ is the dimension of the torus, and $\mathbf 1_T$ is the indicator
function of $T$.

\begin{definition}[The far mass and the far convolution]
\label{def:torus-far-mass}
\lean{Torus.farMass, Torus.farConv, Torus.aemeasurable_tsum_enorm, Torus.aemeasurable_farConv}
\leanok
\uses{def:torus-econv, lem:torus-proj}
For every family $\theta = (\theta_j)_{j \in \N}$ of functions $\R^d \to \C$, the far mass is
\[
  M_\theta = \int_{\{\|y\| \geq 1\}} \sum_j |\theta_j(y)| \, dy \in [0, \infty],
\]
and for every $G : \T^d \to [0, \infty]$ the far convolution is
\[
  F_\theta G(t) = \int_{\{\|y\| \geq 1\}} \Bigl( \sum_j |\theta_j(y)| \Bigr) G(t - \pi y) \, dy .
\]
If every $\theta_j$ is almost everywhere measurable, then $y \mapsto \sum_j |\theta_j(y)|$ is
almost everywhere measurable, and $F_\theta G$ is almost everywhere measurable for every
measurable $G$.
\end{definition}

\begin{lemma}[The far mass of dilation families]
\label{lem:torus-far-mass-dilate}
\lean{Rd.inv_one_add_ofReal_two_pow_le, Torus.farMass_dilate_le, Torus.farMass_lpBlockKernel_le,
  Torus.farMass_psiDual_le}
\leanok
\uses{def:torus-far-mass, def:uniform-schwartz, lem:schwartz-tails,
  lem:uniform-schwartz-constants, def:psi-base, lem:psi-base, def:psi-dual, lem:psi-dual-basic}
For every uniform Schwartz family $(G_j)_j$ on $\R^d$ with bounds $A$, the dilates
$\theta_j = (G_j)_{2^j}$ satisfy
\[
  M_\theta \leq 2 \, T_{A,0,1} .
\]
In particular $M_\psi \leq 2 T_{A^\psi,0,1}$ for the blocks and
$M_{\tilde\psi} \leq 2 T_{\tilde A,0,1}$ for the synthesis kernels.
\end{lemma}

\begin{proof}
\leanok
By monotone convergence $M_\theta$ is the sum over $j$ of the tails
$\int_{\{\|y\| \geq 1\}} |(G_j)_{2^j}(y)| \, dy$. \Cref{lem:schwartz-tails} with $a = 1$ and
$N = 1$ and \Cref{lem:uniform-schwartz-constants} bound the $j$-th tail by
$T_{A,0,1} (1 + 2^j)^{-1} \leq T_{A,0,1} 2^{-j}$, and $\sum_j 2^{-j} = 2$. The blocks and the
synthesis kernels are such dilates by \Cref{lem:psi-base} and \Cref{lem:psi-dual-basic}.
\end{proof}

\begin{lemma}[Young's inequality for the far convolution]
\label{lem:torus-far-conv-young}
\lean{Torus.lintegral_rpow_lintegral_mul_comp_sub_proj_le, Torus.eLpNorm_farConv_le,
  Torus.tsum_setLIntegral_eq_farConv}
\leanok
\uses{def:torus-far-mass, lem:weighted-jensen, lem:lp-norm-nonnegative}
For every $q \geq 1$, every s-finite measure $\nu$ on $\R^d$, every almost everywhere measurable
$a : \R^d \to [0, \infty]$ and every measurable $G : \T^d \to [0, \infty]$,
\[
  \int_{\T^d} \Bigl( \int a(y) G(t - \pi y) \, d\nu(y) \Bigr)^q dt
  \leq \Bigl( \int a \, d\nu \Bigr)^q \int_{\T^d} G^q .
\]
For every $1 \leq p < \infty$, every family $\theta$ of almost everywhere measurable functions and
every measurable $G : \T^d \to [0, \infty]$, $\|F_\theta G\|_p \leq M_\theta \|G\|_p$ and, for
every $t \in \T^d$,
\[
  \sum_j \int_{\{\|y\| \geq 1\}} |\theta_j(y)| \, G(t - \pi y) \, dy = F_\theta G(t) .
\]
\end{lemma}

\begin{proof}
\leanok
This is the kernel form of Young's inequality (\Cref{lem:weighted-jensen}) for
$K(t, y) = G(t - \pi y)$, since $\int_{\T^d} G(t - \pi y)^q \, dt = \int_{\T^d} G^q$ for every
$y$ by the translation invariance of the Haar measure of $\T^d$. The
bound of $F_\theta G$ is the case $q = p$, $\nu$ the Lebesgue measure on $\{\|y\| \geq 1\}$ and
$a = \sum_j |\theta_j|$, with $\|G\|_p^p = \int G^p$ (\Cref{lem:lp-norm-nonnegative}); the sum
over $j$ is monotone convergence.
\end{proof}

\begin{lemma}[The box of side $3$ and the far points]
\label{lem:torus-triple-box}
\lean{Torus.one_le_norm_of_sub_notMem_tripleBox}
\leanok
\uses{def:period-box}
For every $x \in P_0$ and every $y \in \R^d$ with $x - y \notin T$, $\|y\| \geq 1$.
\end{lemma}

\begin{proof}
\leanok
Some coordinate has $(x - y)_i \leq -1$ or $(x - y)_i > 2$; with $0 < x_i \leq 1$ this gives
$y_i \geq 1$ or $y_i < -1$, and $|y_i| \leq \|y\|$.
\end{proof}

\begin{lemma}[Localisation of convolutions]
\label{lem:torus-conv-localisation}
\lean{Torus.enorm_econv_proj_sub_conv_le, Torus.ae_integrable_econv_proj,
  Torus.lqNorm_econv_proj_le, Torus.enorm_sum_econv_proj_le}
\leanok
\uses{lem:torus-triple-box, lem:torus-lift-conv, def:conv, lem:torus-econv-young,
  lem:period-box-measure, lem:lq-basic, lem:torus-far-conv-young, def:torus-far-mass}
For every $\theta : \R^d \to \C$, every $f : \T^d \to \C$ and every $x \in P_0$ such that
$y \mapsto \theta(y) f(\pi x - \pi y)$ is integrable,
\[
  \bigl| (\theta * f)(\pi x) - \bigl(\theta * (\mathbf 1_T (f \circ \pi))\bigr)(x) \bigr|
  \leq \int_{\{\|y\| \geq 1\}} |\theta(y)| \, |f(\pi x - \pi y)| \, dy ,
\]
with the convolution on $\R^d$ in the second term. For every family of integrable
$\theta_j : \R^d \to \C$ and every family of measurable $g_j \in L_1(\T^d)$, the functions
$y \mapsto \theta_j(y) g_j(\pi x - \pi y)$ are integrable for almost every $x \in \R^d$ and all
$j$ at once. For every family of almost everywhere measurable $\theta_j$, every measurable
$f : \T^d \to \C$ and every $x \in P_0$ at which all these functions are integrable,
\[
  \|((\theta_j * f)(\pi x))_j\|_{\ell_2}
  \leq \bigl\| \bigl(\bigl(\theta_j * (\mathbf 1_T (f \circ \pi))\bigr)(x)\bigr)_j \bigr\|_{\ell_2}
  + F_\theta |f| (\pi x) ,
\]
and for every family of measurable $g_j : \T^d \to \C$, every finite set $s \subseteq \N$ and
every $x \in P_0$ at which all the integrals for $g_j$ converge absolutely, with
$G(t) = \|(g_j(t))_j\|_{\ell_2}$,
\[
  \Bigl| \sum_{j \in s} (\theta_j * g_j)(\pi x) \Bigr|
  \leq \Bigl| \sum_{j \in s} \bigl(\theta_j * (\mathbf 1_T (g_j \circ \pi))\bigr)(x) \Bigr|
  + F_\theta G (\pi x) .
\]
\end{lemma}

\begin{proof}
\leanok
By \Cref{lem:torus-lift-conv}, $(\theta * f)(\pi x) = \int \theta(y) f(\pi x - \pi y) \, dy$, and
the convolution with $\mathbf 1_T (f \circ \pi)$ is the integral of the same function over the
$y$ with $x - y \in T$. The difference is the integral over the other $y$, which have
$\|y\| \geq 1$ (\Cref{lem:torus-triple-box}). The integrability holds for almost every $t$ at
$t = \pi x$ (\Cref{lem:torus-econv-young}), and lifts to almost every $x$
(\Cref{lem:period-box-measure}). For the $\ell_2$ bound, the triangle inequality in $\ell_2$
separates the two parts, $\|\cdot\|_{\ell_2} \leq \|\cdot\|_{\ell_1}$ (\Cref{lem:lq-basic}) bounds
the far parts by their sum, and that sum is $F_\theta |f| (\pi x)$
(\Cref{lem:torus-far-conv-young}). For the sums, $|g_j| \leq G$ at every point, and the far parts
over $j \in s$ are at most their sum over all $j$.
\end{proof}

\begin{lemma}[From the period box to the torus]
\label{lem:torus-period-box-split}
\lean{Torus.eLpNorm_le_of_comp_proj_le}
\leanok
\uses{lem:period-box-measure}
For every $1 \leq p \leq \infty$, all almost everywhere measurable $F, G : \T^d \to [0, \infty]$
and every $A : \R^d \to [0, \infty]$ with $F(\pi x) \leq A(x) + G(\pi x)$ for almost every
$x \in P_0$,
\[
  \|F\|_{L_p(\T^d)} \leq \|A\|_{L_p(\R^d)} + \|G\|_{L_p(\T^d)} .
\]
\end{lemma}

\begin{proof}
\leanok
$\|F\|_{L_p(\T^d)} = \|F \circ \pi\|_{L_p(P_0)}$ and
$\|G \circ \pi\|_{L_p(P_0)} = \|G\|_{L_p(\T^d)}$ (\Cref{lem:period-box-measure}); monotonicity
and the triangle inequality in $L_p(P_0)$, and $\|A\|_{L_p(P_0)} \leq \|A\|_{L_p(\R^d)}$ finish.
\end{proof}

\begin{definition}[The strong type constant on the torus]
\label{def:periodic-cz-const}
\lean{Torus.periodicCZConst, Torus.periodicCZConst_ne_top}
\leanok
\uses{def:cz-const}
For $1 < p < \infty$ and $B, H, D \in [0, \infty]$,
\[
  S^{\T}_{p}(B, H, D) = 3^{\#d/p} \, S_{\R^d,p}(B, H) + D ,
\]
with the constant $S_{\R^d,p}(B, H)$ of \Cref{def:cz-const}. It is finite if $B$, $H$ and $D$
are.
\end{definition}

\begin{theorem}[The strong type of vector-valued singular integrals on the torus]
\label{thm:torus-vector-valued-strong-type}
\lean{Torus.eLpNorm_lqNorm_econv_le, Torus.eLpNorm_sum_econv_le}
\leanok
\uses{thm:vector-valued-strong-type, def:periodic-cz-const, lem:torus-conv-localisation,
  lem:torus-period-box-split, lem:triple-box-norm, lem:torus-far-conv-young, def:cz-family,
  lem:torus-econv-young}
Let $\#d \geq 1$ and $1 < p < \infty$. For all $B, H, D \in [0, \infty]$, every
Calder\'on--Zygmund family $\theta = (\theta_j)_j$ on $\R^d$ with overlap $B$ and H\"ormander
constant $H$ and with $M_\theta \leq D$, every $f \in L_p(\T^d)$,
\[
  \bigl\| t \mapsto \|((\theta_j * f)(t))_j\|_{\ell_2} \bigr\|_{L_p(\T^d)}
  \leq S^{\T}_{p}(B, H, D) \, \|f\|_p ,
\]
and for every finite set $s \subseteq \N$ and all measurable $g_j \in L_p(\T^d)$,
\[
  \Bigl\| \sum_{j \in s} \theta_j * g_j \Bigr\|_{L_p(\T^d)}
  \leq S^{\T}_{p}(B, H, D) \, \bigl\| t \mapsto \|(g_j(t))_j\|_{\ell_2} \bigr\|_{L_p(\T^d)} .
\]
\end{theorem}

\begin{proof}
\leanok
On the torus $L_p \subseteq L_1$, so \Cref{lem:torus-conv-localisation} applies at almost every
point of $P_0$, and \Cref{lem:torus-period-box-split} bounds the left side by the sum of two
terms. The first is the $L_p(\R^d)$ norm of the analysis operator, or of a synthesis operator,
of the family on $\R^d$ applied to $\mathbf 1_T (f \circ \pi)$, or to the
$\mathbf 1_T (g_j \circ \pi)$, which lie in $L_p(\R^d)$ (\Cref{lem:triple-box-norm}).
\Cref{thm:vector-valued-strong-type} bounds it by $S_{\R^d,p}(B, H)$ times
$\|\mathbf 1_T (f \circ \pi)\|_p = 3^{\#d/p} \|f\|_p$, respectively times
$\|\mathbf 1_T (G \circ \pi)\|_p = 3^{\#d/p} \|G\|_p$ with $G(t) = \|(g_j(t))_j\|_{\ell_2}$, since
$\|(\mathbf 1_T (g_j \circ \pi)(x))_j\|_{\ell_2} = \mathbf 1_T (G \circ \pi)(x)$. The second is
the norm of a far convolution, at most $M_\theta \|f\|_p$, respectively $M_\theta \|G\|_p$
(\Cref{lem:torus-far-conv-young}).
\end{proof}

\section{The square function theorem on the torus}
\label{sec:torus-square-function}

The square function theorem of \Cref{sec:square-function} holds on the torus with the same
proof: the blocks $\psi_j$ and the synthesis kernels $\tilde\psi_j$ are Calder\'on--Zygmund
families of dilates on $\R^d$ (\Cref{thm:cz-family-psi}, \Cref{thm:cz-family-psi-dual}), their
far masses are finite (\Cref{lem:torus-far-mass-dilate}), and the strong type on the torus
(\Cref{thm:torus-vector-valued-strong-type}) replaces the strong type on $\R^d$. For
$1 < p < \infty$ this gives $\TL{0}{p}{2}(\T^d) = L_p(\T^d)$ with equivalent norms.

\begin{definition}[The square function constants on the torus]
\label{def:torus-square-function-consts}
\lean{Torus.squareFunctionConst, Torus.squareFunctionConst_ne_top, Torus.squareSynthesisConst,
  Torus.squareSynthesisConst_ne_top}
\leanok
\uses{def:periodic-cz-const, def:hormander-const, def:uniform-schwartz, def:psi-base,
  def:psi-dual}
For $1 < p < \infty$,
\[
  Q^{\T}_{p} = S^{\T}_{p}\bigl(2, 4 H_{A^\psi}, 2 T_{A^\psi,0,1}\bigr) , \qquad
  \tilde Q^{\T}_{p} = S^{\T}_{p}\bigl(4, 4 H_{\tilde A}, 2 T_{\tilde A,0,1}\bigr) .
\]
Both are finite.
\end{definition}

\begin{theorem}[The upper square function bound on the torus]
\label{thm:torus-square-function-upper}
\lean{Torus.triebelLizorkinNorm_zero_two_le, Torus.memTriebelLizorkin_zero_two}
\leanok
\uses{def:torus-triebel-lizorkin, prop:torus-tl-elementary, thm:cz-family-psi,
  lem:torus-far-mass-dilate, thm:torus-vector-valued-strong-type,
  def:torus-square-function-consts}
Let $\#d \geq 1$ and $1 < p < \infty$. For every $f \in L_p(\T^d)$,
\[
  \|f\|_{\TL{0}{p}{2}(\T^d)} \leq Q^{\T}_{p} \, \|f\|_p ,
\]
and in particular $f \in \TL{0}{p}{2}(\T^d)$.
\end{theorem}

\begin{proof}
\leanok
By \Cref{prop:torus-tl-elementary} the norm $\|f\|_{\TL{0}{p}{2}(\T^d)}$ is the $L_p$ norm of the
square function $t \mapsto \|((\psi_j * f)(t))_j\|_{\ell_2}$. The blocks form a
Calder\'on--Zygmund family with overlap $2$ and H\"ormander constant $4 H_{A^\psi}$
(\Cref{thm:cz-family-psi}) and far mass at most
$2 T_{A^\psi,0,1}$ (\Cref{lem:torus-far-mass-dilate}), and
\Cref{thm:torus-vector-valued-strong-type} bounds its analysis operator.
\end{proof}

\begin{lemma}[The reproducing formula on the torus]
\label{lem:torus-square-function-reproducing}
\lean{Torus.econvLp_psiDual_lpBlock, Torus.sum_econv_psiDual_lpBlock_ae}
\leanok
\uses{lem:torus-econv-identities, lem:psi-dual-reproduces, def:torus-lp-blocks,
  lem:torus-econv-young, lem:psi-dual-basic}
Let $1 \leq p < \infty$ and $f \in L_p(\T^d)$. For every $j \in \N$,
$\tilde\psi_j * \Delta_j f = \Delta_j f$, and for every finite set $s \subseteq \N$,
\[
  \sum_{j \in s} \tilde\psi_j * \Delta_j f = \sum_{j \in s} \Delta_j f
\]
almost everywhere, with the convolution integrals of representatives of the $\Delta_j f$ on the
left.
\end{lemma}

\begin{proof}
\leanok
$\tilde\psi_j * (\psi_j * f) = (\tilde\psi_j * \psi_j) * f$ (\Cref{lem:torus-econv-identities})
and $\tilde\psi_j * \psi_j = \psi_j$ (\Cref{lem:psi-dual-reproduces}). The convolution integral
represents the operator of \Cref{lem:torus-econv-young}.
\end{proof}

\begin{theorem}[The lower square function bound on the torus]
\label{thm:torus-square-function-lower}
\lean{Torus.enorm_le_triebelLizorkinNorm_zero_two}
\leanok
\uses{lem:torus-square-function-reproducing, thm:cz-family-psi-dual, lem:torus-far-mass-dilate,
  thm:torus-vector-valued-strong-type, prop:torus-tl-elementary,
  def:torus-square-function-consts, lem:torus-lp-blocks, def:torus-triebel-lizorkin}
Let $\#d \geq 1$ and $1 < p < \infty$. For every $f \in L_p(\T^d)$,
\[
  \|f\|_p \leq \tilde Q^{\T}_{p} \, \|f\|_{\TL{0}{p}{2}(\T^d)} .
\]
\end{theorem}

\begin{proof}
\leanok
For every $J \in \N$, \Cref{lem:torus-square-function-reproducing} and
\Cref{thm:torus-vector-valued-strong-type} for the Calder\'on--Zygmund family
$(\tilde\psi_j)_j$ with overlap $4$, H\"ormander constant $4 H_{\tilde A}$
(\Cref{thm:cz-family-psi-dual}) and far mass at most $2 T_{\tilde A,0,1}$
(\Cref{lem:torus-far-mass-dilate}), applied to $g_j = \Delta_j f$, give
\[
  \Bigl\| \sum_{j < J} \Delta_j f \Bigr\|_p
  \leq \tilde Q^{\T}_{p} \, \bigl\| t \mapsto \|(\Delta_j f(t))_j\|_{\ell_2} \bigr\|_p
  = \tilde Q^{\T}_{p} \, \|f\|_{\TL{0}{p}{2}(\T^d)} ,
\]
since $\Delta_j f = \psi_j * f$ almost everywhere for all $j$ at once
(\Cref{def:torus-triebel-lizorkin}, \Cref{prop:torus-tl-elementary}). The left side tends
to $\|f\|_p$ as $J \to \infty$ (\Cref{lem:torus-lp-blocks}).
\end{proof}

\begin{corollary}[The square function theorem on the torus]
\label{cor:torus-square-function-theorem}
\lean{Torus.triebelLizorkinSpace_zero_two_eq_top, Torus.memTriebelLizorkin_two_of_memBesov}
\leanok
\uses{thm:torus-square-function-upper, thm:torus-square-function-lower,
  prop:torus-tl-elementary, thm:torus-tl-besov-sandwich}
Let $\#d \geq 1$ and $1 < p < \infty$. Then $\TL{0}{p}{2}(\T^d) = L_p(\T^d)$, and for every
$f \in L_p(\T^d)$,
\[
  \|f\|_p \leq \tilde Q^{\T}_{p} \, \|f\|_{\TL{0}{p}{2}(\T^d)} , \qquad
  \|f\|_{\TL{0}{p}{2}(\T^d)} \leq Q^{\T}_{p} \, \|f\|_p .
\]
For every $s \geq 0$, $\Besov{s}{p}{\min(p,2)}(\T^d) \subseteq \TL{s}{p}{2}(\T^d)$.
\end{corollary}

\begin{proof}
\leanok
\Cref{thm:torus-square-function-upper} and \Cref{thm:torus-square-function-lower}. For $s > 0$
the inclusion is the Besov sandwich (\Cref{thm:torus-tl-besov-sandwich}); for $s = 0$ every
$f \in L_p(\T^d)$ lies in $\TL{0}{p}{2}(\T^d)$.
\end{proof}
